\documentclass[11pt]{article}
\usepackage[margin=1in]{geometry}
\usepackage{amsmath,amssymb,amsthm,mathtools,bm,mathrsfs}
\usepackage[hidelinks]{hyperref}
\newtheorem{theorem}{Theorem}[section]
\newtheorem{lemma}[theorem]{Lemma}
\newtheorem{proposition}[theorem]{Proposition}
\newtheorem{corollary}[theorem]{Corollary}
\newcommand{\SU}{\operatorname{SU}}
\newcommand{\tr}{\operatorname{tr}}
\newcommand{\Var}{\operatorname{Var}}
\newcommand{\Rea}{\operatorname{Re}}
\newcommand{\dd}{\,\mathrm d}
\newcommand{\HH}{\mathscr H}
\newcommand{\KK}{\mathscr K}
\newcommand{\LL}{\mathscr L}
\newcommand{\VV}{\mathscr V}
\newcommand{\EE}{\mathsf E}
\newcommand{\PP}{\mathsf P}
\newcommand{\NN}{\mathcal N}
\title{Extensive true-vacuum quantum blocking in the full spatial Wilson box\\
Exact path-product maps, associative memory, and physical spectral channels}
\author{}
\date{8 September 2026}
\begin{document}
\maketitle

\section{Original objects and the finite interacting operator}

Let $L\geq2$ be an integer. The vertex set is
$\mathsf V_L=\{-L,-L+1,\ldots,L\}^3$. For each $i\in\{1,2,3\}$,
the physical edge $e=(n,i)$ goes from $n$ to $n+\mathbf e_i$ when both
vertices are present. Its variable is $U_i(n)\in\SU(2)$. Inverse traversal
uses $U_i(n)^{-1}$. Every such edge belongs to $\EE_L$. Every elementary
face $p=(n;i,j)$, $i<j$, with four vertices in the box belongs to $\PP_L$;
its ordered holonomy and trace are
\begin{equation}\label{eq:plaquette}
 U_p=U_i(n)U_j(n+\mathbf e_i)U_i(n+\mathbf e_j)^{-1}U_j(n)^{-1},
 \qquad W_p=\tr U_p.
\end{equation}
The counts are $N=3(2L)(2L+1)^2$ and $M=3(2L)^2(2L+1)$: fix a
direction or a face plane and count its initial and transverse coordinates.
The configuration space is $\mathcal Q=\SU(2)^{\EE_L}$ with product
Haar probability measure $\lambda$. Put
\begin{gather}
 T_a=-i\sigma_a/2,\qquad
 X_{e,a}f=\left.\frac{d}{dt}f(\ldots,e^{tT_a}U_e,\ldots)\right|_{t=0},
 \qquad E_e=-\sum_{a=1}^3X_{e,a}^2,\label{eq:fields}\\
 H=\kappa\sum_{e\in\EE_L}E_e+b\sum_{p\in\PP_L}(2-W_p),\qquad
 \kappa=\frac{2g^2}{a},\quad b=\frac{1}{2g^2a},\quad
 \xi=\frac b\kappa=\frac1{4g^4},\quad a,g>0.\label{eq:H}
\end{gather}
Here $a$ in $a,g>0$ denotes spatial spacing; the color index in
(2) is summed only when indicated. Haar integration by parts
makes $X_{e,a}$ skew-adjoint. The Pauli identities give
$-2\tr(T_aT_b)=\delta_{ab}$ and $-\sum_aT_a^2=3I/4$.
The original metric $c(A,B)=-\tr(AB)/2$ has orthonormal basis $2T_a$;
hence $X^c=2X$, $\Delta^c=4\Delta^T$, and the same kinetic term is
$-(\kappa/4)\sum_e\Delta_e^c$. All faces, including boundary faces
and faces in every plane, remain in (3).
In particular the exact scalar term $2bM=M/(g^2a)$ is part of $H$.

Gauge transformations are $(h\cdot U)_e=h_{s(e)}^{-1}U_eh_{t(e)}$.
The form of $H$ on $H^1(\mathcal Q)$ is
\[
 q_H(u,v)=\kappa\sum_{e,a}\int\overline{X_{e,a}u}X_{e,a}v\dd\lambda
        +b\sum_p\int(2-W_p)\overline uv\dd\lambda.
\]
It is closed, since the smooth potential is bounded between $0$ and
$4bM$. Its associated self-adjoint operator has domain $H^2$ by elliptic
regularity; smooth functions are operator and form cores. Compact Sobolev
embedding gives compact resolvent. Taking the modulus of a minimizing
vector does not increase its form: apply the chain rule first to
$(|u|^2+\epsilon^2)^{1/2}$ and then let $\epsilon$ decrease to zero.
The resulting nonnegative lowest eigenfunction is smooth by elliptic
regularity and strictly positive by the elliptic strong maximum principle.
Write it as $\psi>0$, with $H\psi=E_0\psi$ and $\int\psi^2\dd\lambda=1$.
This fixes the state norm; none of the physical coefficients is changed.

Set $\rho=\psi^2$, $\mu=\rho\lambda$, and $\HH=L^2(\mu)$.
Multiplication $\mathcal Uu=\psi u$ is a unitary $\HH\to L^2(\lambda)$,
with inverse division by $\psi$. The product rule and $H\psi=E_0\psi$ give
\begin{align}
 \LL&=\mathcal U^{-1}(H-E_0)\mathcal U
 =-\kappa\sum_{e,a}\bigl[X_{e,a}^2+(X_{e,a}\log\rho)X_{e,a}\bigr],
 \label{eq:gs}\\
 q(u,v)&=\kappa\sum_{e,a}\int\overline{X_{e,a}u}X_{e,a}v\dd\mu,
 \qquad D(q)=\VV=H^1(\mathcal Q).\label{eq:q}
\end{align}
Integration by parts proves the form identity. Multiplication and division
by the smooth positive $\psi$ preserve every integer Sobolev space, so
the form and operator domains are respectively $H^1,H^2$.
If $q(u,u)=0$, all its derivatives vanish and connectedness makes $u$
constant. This proves uniqueness of $\psi$, and its gauge invariance
follows because gauge transformations preserve $H$, positivity and norm.
It also proves that the finite first nonzero eigenvalue $\delta$ of
$\LL$ is strictly positive. No uniform lower bound for $\delta$ is used
in the uniform estimates below. Inner products are conjugate-linear in
their first arguments.

\section{A uniform single-link estimate for the actual vacuum}

This section supplies estimates independent of $L$, with every $b>0$
allowed. Define the constants
\begin{equation}\label{eq:uniformconstants}
 D=16b,\quad d=D/\kappa=16\xi=4/g^4,\quad
 t_0=\frac{2\log12}{\kappa},\quad
 R=3e^{Dt_0}=3\,12^{2d},\quad \alpha=R^{-2},
\end{equation}
and
\begin{equation}\label{eq:G}
 G=R\exp\!\left(\frac d{1+d}\right)
       \left(1+\frac{2d}{1+d}\right)\sqrt{\frac{1+d}{2}},
 \qquad C_*=\frac32+2G.
\end{equation}
Here $D$ has units of energy and $t_0$ of inverse energy.
The constants $d,R,\alpha,G,C_*$ are dimensionless and depend on the
original coupling $g$, but not on the number of links, the blocking
length, or $a$ separately.

\paragraph{Lemma 2.1.}
For the single-link heat operator $P_t=e^{-t\kappa E}$, its kernel at
$t=t_0$ lies between $1/2$ and $3/2$. For every real smooth $f$,
\begin{equation}\label{eq:reverse}
 2\kappa t\sum_a|X_aP_tf|^2
 \leq P_t(f^2)-(P_tf)^2,\qquad
 \|\nabla P_tf\|_\infty\leq\frac{\|f\|_\infty}{\sqrt{2\kappa t}}.
\end{equation}

\paragraph{Proof.}
The spin-$j$ representation has dimension $2j+1$ and Casimir $j(j+1)$
in (2). One can verify the latter by the eigenvalues
$m=-j,\ldots,j$ of $iT_3$ and the raising/lowering coefficients
$\sqrt{(j\mp m)(j\pm m+1)}$; their quadratic sum is $j(j+1)$.
Peter--Weyl completeness therefore gives the heat kernel
$p_t(U)=\sum_{j\geq0}(2j+1)e^{-\kappa tj(j+1)}\chi_j(U)$.
This series is absolutely uniformly convergent for $t>0$, since
$|\chi_j|\leq2j+1$. Writing $k=2j$, using
$j(j+1)=k(k+2)/4\geq k/2$ and $(k+1)^2\leq4^k$ for $k\geq1$, gives
\[
 |p_{t_0}-1|\leq\sum_{k\geq1}4^k e^{-\kappa t_0k/2}
 =\sum_{k\geq1}3^{-k}=\frac12.
\]
For the gradient assertion, the bi-invariant Casimir commutes with
each $X_a$: the commutator of its sum of squares with $X_b$ contracts
antisymmetric Lie structure constants with $X_aX_c+X_cX_a$ and is zero.
Thus $X_aP_s=P_sX_a$ on smooth functions, and Jensen's inequality gives
$\Gamma(P_sf)\leq P_s\Gamma(f)$ for $\Gamma(f)=\sum_a|X_af|^2$.
Differentiate $P_s[(P_{t-s}f)^2]$ and use the product rule:
\[
 P_t(f^2)-(P_tf)^2=2\kappa\int_0^tP_s\Gamma(P_{t-s}f)\dd s
 \geq2\kappa t\Gamma(P_tf).
\]
The last inequality applies the preceding gradient contraction with
input $P_{t-s}f$. The supremum estimate follows because
$P_t(f^2)\leq\|f\|_\infty^2$. This is the exact single-group version
of the reverse Poincare inequality \href{https://doi.org/10.4171/RMI/470}{[2, Eq.~(1.7)]}, with the time
and generator coefficient displayed explicitly.



\paragraph{Theorem 2.2.}
Fix one physical link $e$ and write all other variables as $y$. For all
$u,u'\in\SU(2)$ and $y$,
\begin{equation}\label{eq:ratio}
 R^{-1}\leq\frac{\psi(u,y)}{\psi(u',y)}\leq R,\qquad
 \alpha\leq\frac{\rho(u,y)}{\int\rho(v,y)\dd v}\leq\alpha^{-1}.
\end{equation}
Moreover,
\begin{equation}\label{eq:scorebound}
 \left(\sum_a|X_{e,a}\log\psi|^2\right)^{1/2}\leq G,
 \qquad |\nabla_e\log\rho|\leq2G.
\end{equation}
All bounds are uniform in $e$ and $L$ at the fixed positive $a,g$.

\paragraph{Proof.}
Let $V_e=b\sum_{p\ni e}(2-W_p)$ and let $H_{\neg e}$ contain all
other kinetic terms and all faces not incident to $e$. There are at
most four incident faces in the full three-dimensional box, including
the boundary cases with fewer faces. Therefore $0\leq V_e\leq D$,
and the exact operator is
\[
 H=A_e+V_e,\qquad A_e=\kappa E_e+H_{\neg e},\qquad
 T_t=e^{-tA_e}=P_t^{(e)}\otimes e^{-tH_{\neg e}}.
\]
The two factors of $T_t$ commute because $H_{\neg e}$ is independent
of $u$. They have positive kernels. The parabolic comparison principle
with the bounded multiplication $0\leq V_e\leq D$ gives, for $f\geq0$,
\begin{equation}\label{eq:comparison}
 e^{-Dt}T_tf\leq e^{-tH}f\leq T_tf.
\end{equation}
For example $T_tf$ is a supersolution for
$\partial_t+A_e+V_e$, while $e^{-Dt}T_tf$ is a subsolution; their
initial data agree. The maximum principle gives (11).
This proves the comparison for this compact group product directly;
no Euclidean configuration-space identification is needed.

Apply it to $\psi$. At $t_0$, positivity of the other-link kernel and
Lemma~2.1 imply
\[
 \frac{(T_{t_0}\psi)(u,y)}{(T_{t_0}\psi)(u',y)}\leq3.
\]
Since $e^{-tH}\psi=e^{-tE_0}\psi$, comparison at both $u,u'$ gives
$\psi(u,y)/\psi(u',y)\leq3e^{Dt_0}=R$.
Interchanging $u,u'$ proves both sides. Squaring and averaging the
denominator over $u'$ proves the conditional density bounds.
The extensive energy $E_0$ cancels from this ratio exactly.

Here is a derivative proof that does not differentiate an unknown
vacuum estimate. If a positive function $F(u,y)$ satisfies the first
ratio bound, so does
$f_t(u)=(e^{-tH_{\neg e}}F(u,\cdot))(y)$ at each fixed $y$, by
positivity. Hence $\|f_t\|_\infty\leq R\min f_t\leq RP_t f_t(u)$.
Lemma~2.1 yields
\[
 |\nabla_e T_t\psi|\leq\frac{R}{\sqrt{2\kappa t}}T_t\psi.
\]
The function $(e^{-tH_{\neg e}}(V_e\psi))(u,y)$ is nonnegative and
bounded above by $D(e^{-tH_{\neg e}}\psi)(u,y)$. Its supremum in $u$
is consequently at most $DR\min_u(e^{-tH_{\neg e}}\psi)(u,y)$.
The same heat estimate gives
\[
 |\nabla_e T_t(V_e\psi)|\leq\frac{DR}{\sqrt{2\kappa t}}T_t\psi.
\]
The exact Duhamel identity on the smooth eigenvector is
\begin{equation}\label{eq:duhamel}
 \psi=e^{E_0t}T_t\psi-\int_0^t e^{E_0\tau}T_\tau(V_e\psi)\dd\tau.
\end{equation}
It follows either by differentiating $e^{E_0t}T_t\psi$ and using
$A_e\psi=(E_0-V_e)\psi$, or by the bounded-potential Duhamel formula.
Comparison gives $e^{E_0\tau}T_\tau\psi\leq e^{D\tau}\psi$.
The two gradient estimates, whose $\tau^{-1/2}$ singularity is
integrable, justify differentiation of (12) and give
\[
 \frac{|\nabla_e\psi|}{\psi}
 \leq R\left[\frac{e^{Dt}}{\sqrt{2\kappa t}}
       +D\int_0^t\frac{e^{D\tau}}{\sqrt{2\kappa\tau}}\dd\tau\right]
 \leq\frac{Re^{Dt}(1+2Dt)}{\sqrt{2\kappa t}}.
\]
Choose the actual time $t=(\kappa+D)^{-1}$. This last expression is
exactly $G$ in (7). Finally $X\log\rho=2X\log\psi$.
No gap bound, perturbation disk, or correlation factorization was used.



\section{A specified extensive spatial blocking sequence}

Let $m\geq1$ divide $2L$, and put $q_m=2L/m$. Retain as skeleton links
those direction-$i$ fine edges whose two transverse coordinates lie in
$-L+m\{0,1,\ldots,q_m\}$. Every longitudinal fine edge on each such
line is included. Denote this subset by $R_m$ and its complement by $S_m$.
Its exact count is
\begin{equation}\label{eq:counts}
 |R_m|=3(2L)(q_m+1)^2,\qquad
 |S_m|=3(2L)\bigl[(2L+1)^2-(q_m+1)^2\bigr].
\end{equation}
Thus a fixed $m>1$ eliminates an extensive number of links as the box
increases. If $m\mid n\mid2L$, then $R_n\subset R_m$. In particular
$m=1,2,4,\ldots,2^k$ gives a repeated blocking sequence whenever $2^k\mid2L$.
The Hamiltonian still has all $N$ link and $M$ face terms.

The coarse vertices are the points
$x\in[-L,L]^3$ whose coordinates all lie in $-L+m\mathbb Z$.
A coarse direction-$i$ edge $\gamma=(x,i)$ is the ordered path of $m$
fine edges from $x$ to $x+m\mathbf e_i$. The number of coarse edges is
\[
 N_m=3q_m(q_m+1)^2,\qquad |R_m|=mN_m.
\]
The chains have disjoint physical-edge supports and meet only at coarse
vertices. Write their variables as $u_{\gamma,1},\ldots,u_{\gamma,m}$
in positive order, and define
\begin{equation}\label{eq:path}
 z_\gamma=u_{\gamma,1}\cdots u_{\gamma,m},\qquad
 \pi_m:\mathcal Q\longrightarrow\mathcal Q_m:=\SU(2)^{N_m}.
\end{equation}
No original coordinate was replaced by a unit coarse spacing: the physical
coarse spacing is $ma$, while the original coordinates remain $x$.

For every chain put $v_{\gamma,k}=u_{\gamma,k}$ for $k<m$.
Let $y$ comprise all $v$ and all untouched variables on $S_m$.
The full typed coordinate map is
\[
 \Phi_m:\mathcal Q\longrightarrow\mathcal Y_m\times\mathcal Q_m,
 \quad U\longmapsto(y,z),\qquad
 \mathcal Y_m=\SU(2)^{S_m}\times
                \prod_{\gamma=1}^{N_m}\SU(2)^{m-1}.
\]
It is a global smooth bijection. Its inverse
$\Phi_m^{-1}:\mathcal Y_m\times\mathcal Q_m\to\mathcal Q$ uses
\[
 u_{\gamma,m}=(v_{\gamma,1}\cdots v_{\gamma,m-1})^{-1}z_\gamma,
\]
and keeps every other variable. Substituting the original ordered
product cancels its prefix on the left and recovers $u_{\gamma,m}$.
Conversely multiplying the retained prefix by the displayed last link
recovers $z_\gamma$. Every $v$ and every $S_m$ coordinate is fixed,
so both compositions are the identity on their respective full spaces.
There are $|S_m|+(m-1)N_m=N-N_m$ factors in $\mathcal Y_m$;
none of those coordinates is omitted from $\Phi_m$.
Integrating $u_{\gamma,m}$ last,
left invariance of Haar proves the exact measure identity
\begin{equation}\label{eq:haarcoordinates}
 \dd\lambda(U)=\dd\lambda_y(y)\dd\lambda_m(z).
\end{equation}
This proves the map and its measure factor, rather than postulating
independent coarse links in the interacting vacuum.

The complete kinetic operator in these coordinates is as follows.
Define the prefix and the real
orthogonal matrix $O_{\gamma,k}$ by
\[
 p_{\gamma,k}=v_{\gamma,1}\cdots v_{\gamma,k-1},\qquad
 p_{\gamma,k}T_a p_{\gamma,k}^{-1}
   =\sum_b O_{\gamma,k;ba}T_b.
\]
Orthogonality follows by applying $-2\tr(AB)$ to both conjugated factors;
its determinant is $1$ by connectedness. Let $D_{\gamma,b}$ be left
multiplication differentiation on $z_\gamma$, with the same $T_b$.
For $k<m$ the transformed fine vector field is
\[
 X_{\gamma,k,a}=X_{v_{\gamma,k},a}
            +\sum_bO_{\gamma,k;ba}D_{\gamma,b};
\]
for $k=m$ it is $\sum_bO_{\gamma,m;ba}D_{\gamma,b}$. The prefix
$p_{\gamma,k}$ does not involve $v_{\gamma,k}$, and no $O$ depends on $z$.
The full sum of squares is consequently
\begin{align}\label{eq:fulltensor}
 \sum_{e,a}X_{e,a}^2
 ={}&\sum_{e\in S_m,a}X_{e,a}^2
 +\sum_{\gamma}\left[\sum_{k<m,a}X_{v_{\gamma,k},a}^2
                      +m\sum_bD_{\gamma,b}^2\right]\nonumber\\
 &+2\sum_{\gamma,k<m,a,b}O_{\gamma,k;ba}
                    X_{v_{\gamma,k},a}D_{\gamma,b}.
\end{align}
Each square before expansion is the square of the displayed original
field, so this formula includes all mixed derivatives with their signs
and coefficients. Every original $W_p$ in (3) is evaluated at
the explicit inverse coordinate map. It is not replaced by a coarse face
trace or truncated to selected words.

Gauge transformations on a chain cancel at its internal vertices and give
$z_\gamma\mapsto h_x^{-1}z_\gamma h_{x+m\mathbf e_i}$.
The marginal density on $R_m$ is invariant under the internal gauge
transformations, by gauge invariance of $\rho$ and Haar substitution on
$S_m$. A non-coarse skeleton vertex is bivalent in the skeleton. For a
fixed chain with endpoints fixed, set successively
$h_{x+k\mathbf e_i}=(u_1\cdots u_k)^{-1}$, $k<m$.
This sends its first $m-1$ links to the identity and its last to $z_\gamma$.
Different chain interiors are disjoint, so these choices are simultaneous.
It follows that this skeleton marginal is a function of $z$ alone.
This statement concerns the marginal; the full conditional distribution
of $y$ given $z$ need not factorize.

\section{The actual coarse density, domains, and conditional coupling}

Let $\widetilde\rho_m(y,z)=\rho(U(y,z))$ and define
\begin{gather}
 r_m(z)=\int\widetilde\rho_m(y,z)\dd\lambda_y,
 \quad \mu_m=r_m\lambda_m,\quad
 p_m(y\mid z)=\widetilde\rho_m(y,z)/r_m(z),\label{eq:coarsedensity}\\
 \KK_m=L^2(\mu_m),\quad J_m f=f\circ\pi_m,\quad
 P_m=J_mJ_m^*,\quad Q_m=I-P_m,\nonumber\\
 (J_m^*u)(z)=\int p_m(y\mid z)u(U(y,z))\dd\lambda_y.\label{eq:J}
\end{gather}
In particular the types are $J_m:\KK_m\to\HH$ and
$J_m^*:\HH\to\KK_m$. Their full coordinate expressions are
$J_mf(U)=f(z(U))$ and the fibre integral (18);
the operator $P_m:\HH\to J_m\KK_m\subset\HH$ is their composition.
Smoothness and strict positivity follow from compact integration.
Fubini proves the adjoint formula and $J_m^*J_m=I$; Jensen proves that
$J_m^*$ is a contraction. Thus $P_m$ is the true-vacuum orthogonal
projection. All these maps intertwine the stated fine and coarse gauge
actions. The restriction of the conditional projection onto skeleton
functions, on the fine invariant Hilbert space, equals this path projection:
its output is internally gauge invariant, hence depends only on $z$, as
proved in the preceding section. The exact invariant Hilbert-space
unitary, its almost-everywhere inverse and the equality of the full
complementary operator domains are proved in
Proposition~10.1.

For each chain and color set, in the original coordinates,
\begin{equation}\label{eq:Y}
 Y_{\gamma,b}=\frac1m\sum_{k=1}^m\sum_a
                   O_{\gamma,k;ba}X_{\gamma,k,a}.
\end{equation}
Then $Y_{\gamma,b}J_mf=J_mD_{\gamma,b}f$.
Indeed each summand differentiates $z_\gamma$ in the direction
$\operatorname{Ad}_{p_{\gamma,k}}T_a$, and orthogonality of $O$ gives
one copy of $D_{\gamma,b}$ per $k$.
Also $Y_{\gamma,b}$ has zero Haar divergence: every coefficient
$O_{\gamma,k;ba}$ is independent of its own differentiated fine edge.
Consequently integration by parts against a coarse smooth test function
gives the exact derivative identities
\begin{align}
 D_{\gamma,b}\log r_m&=J_m^*(Y_{\gamma,b}\log\rho),\label{eq:marginalscore}\\
 D_{\gamma,b}J_m^*u
   &=J_m^*(Y_{\gamma,b}u)+J_m^*(\beta_{\gamma,b}u),\label{eq:derivativeJ}\\
 \beta_{\gamma,b}&=Y_{\gamma,b}\log\rho
                       -J_mD_{\gamma,b}\log r_m,\qquad J_m^*\beta_{\gamma,b}=0.
 \label{eq:beta}
\end{align}
For example integrate $Y_{\gamma,b}(J_m\varphi\,\rho)$ over the full
configuration space; its integral is zero. The resulting identity is
$\int(D_{\gamma,b}\varphi)r_m=-\int\varphi J_m^*(Y_{\gamma,b}\log\rho)r_m$,
which proves (20). Apply the same calculation to
$u\rho$ to prove (21). This also avoids differentiating
fibre coordinates while holding an incompatible fine coordinate fixed.

Theorem~2.2, the average in (19), and conditional
Jensen give
\begin{equation}\label{eq:coarsescorebound}
 |Y_\gamma\log\rho|\leq2G,\qquad
 |D_\gamma\log r_m|\leq2G,\qquad |\beta_\gamma|\leq4G.
\end{equation}
Here each absolute value is the Euclidean norm of its three color
components. The first inequality uses the average of $m$ rotated vectors,
each with norm at most $2G$; it does not assume independence.

\paragraph{Proposition 4.1.}
The exact lifted form and its generator are
\begin{align}
 a_m(f,g)&=q(J_mf,J_mg)
    =\kappa m\sum_{\gamma,b}\int\overline{D_{\gamma,b}f}D_{\gamma,b}g\dd\mu_m,
       &D(a_m)&=\VV_m=H^1(\mathcal Q_m),\label{eq:am}\\
 A_m&=-\kappa m\,r_m^{-1}\sum_{\gamma,b}D_{\gamma,b}(r_mD_{\gamma,b}),
       &D(A_m)&=H^2(\mathcal Q_m).\label{eq:Am}
\end{align}
The maps $J_m,J_m^*$ preserve all integer Sobolev spaces, and $P_m,Q_m$
preserve $\VV$. On $\VV_m$ the full eliminated coupling is
\begin{equation}\label{eq:Bm}
 B_mf=-\kappa m\sum_{\gamma,b}\beta_{\gamma,b}J_mD_{\gamma,b}f,
 \quad J_m^*B_mf=0,\quad \LL J_mf=J_mA_mf+B_mf\quad(f\in H^2).
\end{equation}
In particular the retained kinetic coefficient is exactly $\kappa m$.

\paragraph{Proof.}
The fine derivative of $J_mf$ on edge $(\gamma,k)$ is the orthogonal
rotation $O_{\gamma,k}^TD_\gamma f$; on $S_m$ it is zero.
Summing the squared derivatives over the $m$ edges proves (24).
The smooth positive weight gives closedness and the $H^2$ elliptic domain
of (25). The compact smooth coordinate bijection
(15) preserves each integer Sobolev class with
finite norm bounds, since its derivatives and inverse derivatives are
bounded. The fibre integral with smooth positive kernel $p_m$ has bounded
derivatives of every finite order. This proves the Sobolev mapping claims.

There is also a useful explicit first-order estimate. Define
\[
 \eta_m=\kappa m\left\|\sum_{\gamma,b}\beta_{\gamma,b}^2\right\|_\infty
             \leq16\kappa mN_mG^2.
\]
Orthogonality of the rotations and Cauchy--Schwarz in the $k$ sum give
the pointwise inequality
\[
 m\sum_{\gamma,b}|Y_{\gamma,b}u|^2
 \leq\sum_{e\in R_m,a}|X_{e,a}u|^2.
\]
Conditional Jensen and (21) consequently give
\begin{equation}\label{eq:projectionbound}
 a_m(J_m^*u,J_m^*u)\leq2q(u,u)+2\eta_m\|u\|_\mu^2.
\end{equation}
This extends by smooth approximation. The same derivative identities
show that $J_mf\in H^1$ implies $f\in H^1$, so the retained form space
has exactly the claimed domain.
On a lifted smooth $f$, (16) gives
$\sum_{e,a}X_{e,a}^2J_mf=mJ_m\sum_{\gamma,b}D_{\gamma,b}^2f$.
The drift in (4) is
$-\kappa m\sum_{\gamma,b}(Y_{\gamma,b}\log\rho)J_mD_{\gamma,b}f$.
Subtract (25), and use (22); this proves
(26) for smooth $f$. Density gives the stated $H^2$ identity
and the bounded map $B_m:\VV_m\to Q_m\HH$, with
\begin{equation}\label{eq:Bbound}
 \|B_mf\|_\mu^2\leq\eta_m a_m(f,f).
\end{equation}



The full correlation tensor of the induced first-order terms is
\begin{equation}\label{eq:Fisher}
 I_{\gamma b,\eta c}(z)=J_m^*(\beta_{\gamma,b}\beta_{\eta,c})(z),\quad
 \|B_mf\|^2=(\kappa m)^2\int
 \sum_{\gamma,b,\eta,c}I_{\gamma b,\eta c}
              \overline{D_{\gamma,b}f}D_{\eta,c}f\dd\mu_m.
\end{equation}
Every off-diagonal index is retained. This is the conditional covariance
of the vectors $Y_{\gamma,b}\log\rho$. It is positive semidefinite
because its quadratic form is a conditional expectation of a squared
absolute value. Positivity does not assert that its different entries
vanish or that it is uniformly diagonal-dominant.

There is an exact Haar-space description of the instantaneous marginal
operator. Multiplication $f\mapsto\sqrt{r_m}f$ is a unitary
$\KK_m\to L^2(\lambda_m)$, with inverse
$v\mapsto v/\sqrt{r_m}$. Both maps preserve $H^1$ and $H^2$,
since $r_m$ is smooth and strictly positive on the compact group
product. Direct product differentiation gives
\begin{equation}\label{eq:marginalpotential}
 \sqrt{r_m}A_m r_m^{-1/2}
 =-\kappa m\sum_{\gamma,b}D_{\gamma,b}^2
       +\kappa m\frac{\sum_{\gamma,b}D_{\gamma,b}^2\sqrt{r_m}}{\sqrt{r_m}}.
\end{equation}
The last scalar function equals
$\kappa m\sum_{\gamma,b}[D_{\gamma,b}^2\log r_m/2
 +(D_{\gamma,b}\log r_m)^2/4]$.
It depends on the entire induced density. No equality with a sum of
one-plaquette Wilson terms is asserted. Moreover (30)
is only the instantaneous compression; the next section gives the rest
of the physical operator relation.

\section{Exact memory, resolvent reconstruction, and physical norms}

Set $\HH_{Q_m}=Q_m\HH$ and $\VV_{Q_m}=\VV\cap\HH_{Q_m}$.
The restriction $c_m(h,k)=q(h,k)$ on $\VV_{Q_m}$ is closed and densely
defined. Closedness uses the $L^2$-closed subspace, and density follows
by applying $Q_m$ to smooth approximants, using (27).
Define its self-adjoint operator $C_m$ by
\begin{align}\label{eq:Cdomain}
 D(C_m)=\{h\in\VV_{Q_m}:&\ \exists v\in\HH_{Q_m}\ \text{such that}
 \ c_m(k,h)=\langle k,v\rangle\ \forall k\in\VV_{Q_m}\},
 \qquad C_mh=v.
\end{align}
For smooth $h\in\HH_{Q_m}$ this is $Q_m\LL h$.
Every such $h$ is orthogonal to $1=J_m1$, so $C_m\geq\delta I>0$.
For form-domain arguments no formal compression domain is substituted
for (31).

For $f,g\in\VV_m$ and $h,k\in\VV_{Q_m}$, integration by parts on a
smooth core, followed by continuity, gives
\begin{equation}\label{eq:block}
 q(J_mf+h,J_mg+k)=a_m(f,g)+\langle B_mf,k\rangle
                         +\langle h,B_mg\rangle+c_m(h,k).
\end{equation}
For smooth $h\in\HH_{Q_m}$ the adjoint coupling is
\begin{equation}\label{eq:Bstar}
 B_m^*h=\kappa m\,r_m^{-1}\sum_{\gamma,b}
           D_{\gamma,b}\bigl(r_mJ_m^*(\beta_{\gamma,b}h)\bigr).
\end{equation}
To prove the sign, integrate the negative derivative in (26)
against $h$ using Haar integration by parts in $z$. Differentiating the
conditional integral by (21) expands each derivative
in (33) into terms with $Yh$, $Y\beta$, $\beta^2h$ and
$(D\log r_m)\beta h$. All coefficients are smooth and bounded in a fixed
box, so $B_m^*:\VV_{Q_m}\to\KK_m$ is bounded. Without that regularity
the adjoint is the bounded dual map $\HH_{Q_m}\to\VV_m^*$.

For $z>0$ define on $\VV_m$ the sesquilinear form
\begin{equation}\label{eq:schur}
 s_{m,z}(f,g)=a_m(f,g)+z\langle f,g\rangle
             -\langle B_mf,(C_m+z)^{-1}B_mg\rangle.
\end{equation}
The use of $z$ as a positive spectral shift here is distinct from the
tuple $z_\gamma$ of coarse link variables. The argument of a resolvent
always denotes the scalar shift.

\paragraph{Theorem 5.1.}
The form $s_{m,z}$ is closed with domain $\VV_m$, and its associated
operator $S_{m,z}$ satisfies $S_{m,z}\geq zI$. Its exact full resolvent
and reconstruction maps are
\begin{align}
 J_m^*(\LL+z)^{-1}J_m&=S_{m,z}^{-1},\label{eq:resolvent}\\
 (\LL+z)^{-1}J_mg
 &=\bigl[J_m-(C_m+z)^{-1}B_m\bigr]S_{m,z}^{-1}g.\label{eq:reconstruction}
\end{align}
Every product in these formulas is defined by the displayed form domains.

\paragraph{Proof.}
Completing the square in (32) gives
\begin{align}\label{eq:square}
 q(J_mf+h,J_mf+h)+z(\|f\|^2+\|h\|^2)
 =s_{m,z}(f,f)+\|(C_m+z)^{1/2}[h+(C_m+z)^{-1}B_mf]\|^2.
\end{align}
Its unique minimizing $h$ is $-(C_m+z)^{-1}B_mf$, which lies in
$D(C_m)\subset\VV_{Q_m}$. Nonnegativity of $q$ proves
$s_{m,z}(f,f)\geq z\|f\|^2$. Apply (27) to
$u=J_mf+h$. It gives
\[
 a_m(f,f)+\kappa\|f\|^2\leq2q(u,u)+(2\eta_m+\kappa)\|u\|^2
 \leq K_{m,z}\,[q(u,u)+z\|u\|^2],
\]
where $K_{m,z}=\max\{2,(2\eta_m+\kappa)/z\}$ in the form norm with
reference energy $\kappa$. This constant is dimensionless.
Taking the minimizing $h$ proves a lower form-norm bound. The upper
bound $s_{m,z}(f,f)\leq a_m(f,f)+z\|f\|^2$ follows by choosing $h=0$.
Their equivalence proves closedness. For $v=(\LL+z)^{-1}J_mg$, its
form-domain decomposition $v=J_mf+h$ is legitimate. Testing the weak
resolvent equation against $k\in\VV_{Q_m}$ gives
$h=-(C_m+z)^{-1}B_mf$; testing against $J_mf'$ gives
$s_{m,z}(f',f)=\langle f',g\rangle$. This proves both identities.



The notation in (34) is a closed-form implementation of the
Feshbach--Schur map. The general method and reconstruction are established
in \href{https://doi.org/10.4171/ECR/18-1/5}{[1, Theorem~1.2, Eqs.~(1.10)--(1.16)]}. That theorem's bounded
effective operator assumption is not silently applied to our infinite-rank
projection; the coercivity and domain proof above supplies the needed
extension for this particular operator.

For smooth initial $f_0$, define
$K_m(t)=J_m^*e^{-t\LL}J_m$, $f(t)=K_m(t)f_0$.
The full evolution and (32) give
\begin{align}
 f'(t)&=-A_m f(t)+\int_0^t B_m^*e^{-(t-s)C_m}B_m f(s)\dd s,
       &f(0)&=f_0,\label{eq:memory}\\
 K_m(t)&=J_m^*e^{-t\LL}J_m,\qquad
 \int_0^\infty e^{-zt}K_m(t)\dd t=S_{m,z}^{-1}.\label{eq:laplace}
\end{align}
Here is the domain justification and the sign calculation. Smooth initial
data lie in all powers of the compact elliptic $\LL$. Spectral calculus
and elliptic regularity keep $v(t)=e^{-t\LL}J_mf_0$, $f(t)=J_m^*v(t)$
and $h(t)=Q_mv(t)$ smooth, continuously in the required Sobolev norms.
They satisfy $f'=-A_mf-B_m^*h$, $h'=-B_mf-C_mh$, $h(0)=0$.
Variation of constants gives
$h(t)=-\int_0^te^{-(t-s)C_m}B_mf(s)\dd s$.
Since $\sup_{x\geq0}\sqrt{x}e^{-\tau x}=(2e\tau)^{-1/2}$,
$e^{-\tau C_m}$ maps $\HH_{Q_m}$ to its form space with norm at most
$\sqrt\kappa+(2e\tau)^{-1/2}$ when the squared form norm is
$c_m(h,h)+\kappa\|h\|^2$. The singularity is integrable; the bounded form-space
map (33) makes (38) a Bochner integral in
$\KK_m$. The positive sign follows from the two negative cross terms.
The Laplace identity follows by the spectral theorem and
Theorem~5.1. This explicitly realizes projection memory of
the kind introduced in \href{https://doi.org/10.1143/PTP.33.423}{[4]}, with this Euclidean generator.

For every smooth $f$, orthogonality of the two components of
$\LL J_mf$ gives the full second moment
\begin{equation}\label{eq:second}
 \|\LL J_mf\|^2=\|A_mf\|^2+\|B_mf\|^2.
\end{equation}
Thus $K_m(t)$ is replaced by $e^{-tA_m}$ for all $t$ only when $B_m$
vanishes: compare their scalar second derivatives at $t=0$ on a smooth
core. Conversely $B_m=0$ makes (32) a direct-sum form and
proves that replacement exactly. We do not make that replacement.

The dependence on energy also records the actual reconstructed norm.
For fixed $0\ne f\in\VV_m$ with $\mu_m(f)=0$ and scalar $z\geq0$, put
\begin{equation}\label{eq:dressed}
 w_z=(C_m+z)^{-1}B_mf,\quad
 u_z=\psi(J_mf-w_z),\quad d_f=\|f\|^2.
\end{equation}
The $z=0$ case exists because $C_m\geq\delta>0$ in every fixed box.
The vector is in the original form domain, is orthogonal to $\psi$,
and is gauge invariant when $f$ is. Its exact norm and energy are
\begin{align}
 \|u_z\|_\lambda^2&=d_f+\|w_z\|_\mu^2
                  =\partial_z s_{m,z}(f,f),\label{eq:normderivative}\\
 \mathfrak q_{H-E_0}[u_z]
 &=a_m(f,f)-\langle B_mf,w_z\rangle-z\|w_z\|^2
   =s_{m,z}(f,f)-z\partial_zs_{m,z}(f,f).\label{eq:dressedenergy}
\end{align}
To verify these, orthogonality gives the norm; differentiation of the
bounded resolvent gives $\partial_z(C_m+z)^{-1}=-(C_m+z)^{-2}$.
Testing $(C_m+z)w_z=B_mf$ against $w_z$ gives
$c_m(w_z,w_z)=\langle B_mf,w_z\rangle-z\|w_z\|^2$; insert it in
(32). Nonnegativity of $q$ proves the numerator is nonnegative.
For $B_mf\ne0$ the quotient in (43) divided by
(42) is strictly below $a_m(f,f)/d_f$.
These are energies of actual full-system vectors; no artificial effective
vacuum norm has replaced their eliminated components.

\section{Associativity of repeated blocking with all prior memory}

Let $m\mid n\mid2L$. A coarse $n$-edge is the ordered concatenation of
$n/m$ coarse $m$-edges along its line. Let $\pi_{m,n}$ take these products
and discard the other $m$-coarse edges. Literal ordered multiplication gives
\begin{equation}\label{eq:tower}
 \pi_n=\pi_{m,n}\circ\pi_m,\quad J_n=J_m I_{m,n},\quad
 \mu_n=(\pi_{m,n})_*\mu_m,\quad
 J_n^*=I_{m,n}^*J_m^*.
\end{equation}
Here $I_{m,n}f=f\circ\pi_{m,n}$ is an isometry $\KK_n\to\KK_m$;
its adjoint is conditional integration using $r_m$, not Haar integration
in place of that weight. Fubini proves each equality, including the adjoint
identity. In the full Hilbert space $P_nP_m=P_mP_n=P_n$.
All maps preserve the relevant Sobolev spaces by the same compact
path-product coordinate construction. Coefficient composition in the
instantaneous forms is exactly $\kappa m(n/m)=\kappa n$.

Let $\widehat Q=I-I_{m,n}I_{m,n}^*$ on $\KK_m$. On
$\widehat Q\KK_m\cap\VV_m$, let $t_{22,z}$ be the restriction of the
already completed form $s_{m,z}$. Define its other block forms by
\begin{align*}
 t_{11,z}(f,g)&=s_{m,z}(I_{m,n}f,I_{m,n}g),\\
 t_{21,z}(h,f)&=s_{m,z}(h,I_{m,n}f),\\
 t_{12,z}(f,h)&=s_{m,z}(I_{m,n}f,h).
\end{align*}
The operator $T_{22,z}$ associated to $t_{22,z}$ is positive with lower
bound $z$. Its inverse also acts from its form dual to its form space,
by the coercive weak equation. In that sense the next eliminated component
is $h_f=-T_{22,z}^{-1}t_{21,z}(\cdot,f)$, and the next form is
\begin{equation}\label{eq:recursive}
 s_{n,z}(f,g)=t_{11,z}(f,g)
       -t_{12,z}\bigl(f,T_{22,z}^{-1}t_{21,z}(\cdot,g)\bigr).
\end{equation}
This notation means substitution of the unique weak solution, so does
not require multiplying unspecified unbounded compressions.

\paragraph{Theorem 6.1.}
Equation (45) equals the direct fine-to-$n$ form
(34). In particular
\begin{equation}\label{eq:resolventtower}
 S_{n,z}^{-1}=I_{m,n}^*S_{m,z}^{-1}I_{m,n}
            =J_n^*(\LL+z)^{-1}J_n.
\end{equation}
No memory generated at an earlier level is lost in the next level.

\paragraph{Proof.}
For prescribed $f\in\VV_n$, every $u\in\VV$ with $J_n^*u=f$ has a
unique orthogonal decomposition
$u=J_n f+J_mh+k$, where $h\in\widehat Q\KK_m\cap\VV_m$ and
$k\in\HH_{Q_m}\cap\VV$; obtain it by the nested projections in
(44). Put $v_h=J_m(I_{m,n}f+h)$ to keep both arguments
of the sesquilinear form explicit. The exact variational identity
(37) gives
\begin{align*}
 s_{n,z}(f,f)
 &=\inf_{J_n^*u=f}[q(u,u)+z\|u\|^2]\\
 &=\inf_h\inf_k[q(v_h+k,v_h+k)+z\|v_h+k\|^2]\\
 &=\inf_h s_{m,z}(I_{m,n}f+h,I_{m,n}f+h).
\end{align*}
The infima have unique minimizers by coercivity and the closed form-space
constraints. The weak equation for the last minimizer is precisely the
$T_{22,z}$ equation above. Expanding the last form and polarizing proves
(45). Alternatively compress
(35) by $I_{m,n}$ and use (44); this proves
(46) directly and confirms the same unique closed
form. Both arguments include all intermediate eliminated components.



To display the generated terms explicitly, write
$\mathcal M_{m,z}(f,g)=\langle B_mf,(C_m+z)^{-1}B_mg\rangle$.
Then each of $t_{11,z},t_{12,z},t_{21,z},t_{22,z}$ is the corresponding
block of $a_m+z\langle\cdot,\cdot\rangle-\mathcal M_{m,z}$.
In particular the two off-diagonal blocks contain the corresponding
off-diagonal terms of $-\mathcal M_{m,z}$, and $T_{22,z}$ contains its
entire diagonal term. Setting these terms to zero would change
(45). The exact recursion is independent of the order of
the specified nested eliminations; it does not reinitialize the dynamics
at each marginal $A_m$.

\section{A large-block Wilson channel with uniform norm and spectral bounds}

Let $C$ be a closed coarse edge word using each coarse physical edge at
most once. Its fine path then uses each of its fine physical edges once;
let $r$ be its number of coarse edges and $\ell=mr$ its number of fine
edges. For example, every elementary coarse face has $r=4$, fine perimeter
$\ell=4m$, and physical side length $ma$. Keep its original ordered word,
with inverses on the negative traversals, and set
\begin{gather}
 W_C(z)=\tr\prod_{(\gamma,\epsilon)\in C}z_\gamma^\epsilon,\quad
 \overline W_C=\mu_m(W_C),\quad r_C=\mu_m(W_C^2),\nonumber\\
 f_C=W_C-\overline W_C,\quad u_C=\psi J_mf_C,\quad
 d_C=\|u_C\|_\lambda^2=r_C-(\overline W_C)^2.\label{eq:loop}
\end{gather}
SU(2) traces are real. The fine and coarse gauge actions both fix the
closed trace, so $u_C$ is a physical invariant vector, orthogonal to
$\psi$. The subtraction is the actual interacting mean.

\paragraph{Lemma 7.1.}
For every such loop, at arbitrary positive coupling,
\begin{equation}\label{eq:loopnorm}
 \alpha\leq d_C\leq4,\qquad
 \alpha\leq r_C\leq4-3\alpha.
\end{equation}
The exact unnormalized excitation energy and bounds are
\begin{equation}\label{eq:loopenergy}
 \mathfrak q_{H-E_0}[u_C]
   =\kappa\ell\left(1-\frac{r_C}{4}\right),\qquad
 \frac34\alpha\kappa\ell\leq\mathfrak q_{H-E_0}[u_C]\leq\kappa\ell.
\end{equation}
Consequently
\begin{equation}\label{eq:loopquotient}
 \frac{3\alpha}{16}\kappa\ell\leq
 \frac{\mathfrak q_{H-E_0}[u_C]}{d_C}
 =\kappa\ell\frac{1-r_C/4}{r_C-(\overline W_C)^2}
 \leq\frac{\kappa\ell}{\alpha}.
\end{equation}

\paragraph{Proof.}
Fix any one fine edge used by the loop and condition on all other fine
variables. The ordered holonomy is $A u^{\pm1}B$ for fixed group matrices
$A,B$, so its Haar distribution is Haar by invariance and inversion.
Writing $U=u_0I+i\sum_a u_a\sigma_a$ identifies Haar measure with the
uniform measure on the unit sphere in $\mathbb R^4$. Its coordinate
second moments are $1/4$ by sign and permutation symmetry and
$\sum_{a=0}^3u_a^2=1$. Thus
$\int W_C\dd u=0$, $\int W_C^2\dd u=1$, and
$\int(1-W_C^2/4)\dd u=3/4$.
The conditional density is bounded below by $\alpha$ from
Theorem~2.2. Integrating the two nonnegative functions
$W_C^2$ and $4-W_C^2$ proves the two bounds on $r_C$.
For every real constant $c$, conditional integration gives
$\int(W_C-c)^2p(u\mid y)\dd u\geq\alpha\int(W_C-c)^2\dd u
=\alpha(1+c^2)\geq\alpha$.
Taking $c$ to be that conditional mean proves conditional variance
at least $\alpha$. The law of total variance, obtained by expanding
around the conditional mean and integrating the cross term to zero,
proves $d_C\geq\alpha$. Since $|W_C|\leq2$, $d_C\leq4$.

For a positive occurrence, a based form of the derivative is
$\tr(T_aA_C)$ with $A_C\in\SU(2)$ having trace $W_C$; for a negative
occurrence it is $-\tr(T_aA_C')$. This follows from
$X_a u^{-1}=-u^{-1}T_a$ and cyclic trace invariance. The Pauli expression
above gives $\sum_a|\tr(T_aA_C)|^2=1-W_C^2/4$.
There are exactly $\ell$ used fine edges. Insert this identity in
(5); centering leaves derivatives unchanged. The exact energy
and its lower and upper bounds follow. Divide using
$\alpha\leq d_C\leq4$ to obtain (50).



These bounds are for the same density used by the block projection.
For an additional consistency check, that coarse density itself has a
single-coarse-link conditional lower bound $\alpha$, independent of $m$.
Indeed change only the last fine edge of a specified chain, holding all
$y$ and all other coarse links fixed in (15).
Theorem~2.2 bounds the ratio of the two full densities
by $R^2$. Integrating the same $y$ on both sides bounds the ratio
$r_m(z_\gamma,z_{\ne\gamma})/r_m(z'_\gamma,z_{\ne\gamma})$ by $R^2$.
Averaging over $z'_\gamma$ proves the assertion. This exact argument
does not equate single-link conditional bounds with independence or a
global Poincare inequality.

\paragraph{Theorem 7.2.}
Define the unnormalized actual full-Hamiltonian spectral measure
\begin{equation}\label{eq:nu}
 \mathsf M_C(B)=\langle u_C,\mathbf1_B(H-E_0)u_C\rangle,\qquad
 \mathsf M_C([0,\infty))=d_C.
\end{equation}
No division by $d_C$ is made in this measure. It has no atom at zero.
Its full first and second moments satisfy
\begin{gather}
 \int\omega\dd\mathsf M_C=\kappa\ell(1-r_C/4),\nonumber\\
 \int\omega^2\dd\mathsf M_C
 =\|A_mf_C\|^2+\|B_mf_C\|^2
 \leq C_*^2\kappa^2\ell^2.\label{eq:moments}\\
 \Omega_C=\frac{3\alpha}{32}\kappa\ell,\qquad
 c_H=\frac{9\alpha^3}{1024C_*^2}>0,\qquad
 \boxed{\mathsf M_C([\Omega_C,\infty))\geq c_Hd_C.}\label{eq:highweight}
\end{gather}
In particular $c_H$ is independent of the increasing blocking length $m$
and of total box volume. The full compressed resolvent, including its
memory, obeys at the scalar shift $z=\kappa\ell$
\begin{equation}\label{eq:resolventdeficit}
 z\langle f_C,S_{m,z}^{-1}f_C\rangle
 \leq d_C\left[1-c_H\frac{3\alpha/32}{1+3\alpha/32}\right].
\end{equation}

\paragraph{Proof.}
The state is smooth and orthogonal to the unique vacuum. Spectral
calculus and (49) prove the first moment and the
absence of a zero atom. Every used edge has fundamental Casimir $3/4$,
including inverse occurrences, so
$-\sum_{e,a}X_{e,a}^2W_C=3\ell W_C/4$.
Combining (4) with (10) and
$|\nabla_e W_C|\leq1$ gives the pointwise bound
\[
 |\LL J_mf_C|=|\LL J_mW_C|
 \leq\kappa\ell\left(\frac32+2G\right)=\kappa\ell C_*.
\]
The second-moment bound follows by integration and
(40). In particular the term $\|B_mf_C\|^2$ has not been
deleted from this spectral measure.

Put $A=3\alpha/16$. The first moment is at least $A\kappa\ell d_C$
by (50). With $\Omega_C=A\kappa\ell/2$, its part on
$[\Omega_C,\infty)$ is at least $A\kappa\ell d_C/2$, since the integral
over $[0,\Omega_C)$ is at most $\Omega_C d_C$.
Cauchy--Schwarz against the unnormalized $\mathsf M_C$ gives
\[
 (A\kappa\ell d_C/2)^2
 \leq\left(\int\omega^2\dd\mathsf M_C\right)
                       \mathsf M_C([\Omega_C,\infty)).
\]
Using (52) gives the stronger intermediate bound
$\mathsf M_C([\Omega_C,\infty))\geq A^2d_C^2/(4C_*^2)$.
Since $d_C\geq\alpha$, this is at least
$A^2\alpha d_C/(4C_*^2)=c_Hd_C$, with the state mass retained.
Finally (35) gives
$z\langle f_C,S_{m,z}^{-1}f_C\rangle
=\int z/(\omega+z)\dd\mathsf M_C$.
On $[\Omega_C,\infty)$ its integrand is at most
$z/(\Omega_C+z)$, and elsewhere it is at most $1$.
Its upper bound is therefore
$d_C-\Omega_C\mathsf M_C([\Omega_C,\infty))/(\Omega_C+z)$.
Insert (53) and $z=\kappa\ell$ to obtain
(54).



The same channel can be dressed by the complete eliminated resolvent
with a quantitative bound that continues to hold for growing $m$.
Define
\begin{equation}\label{eq:tau}
 \tau_* =\frac{8G}{\sqrt{\alpha c_H}},\qquad
 z_C=\tau\kappa\ell,\quad\tau\geq\tau_*.
\end{equation}
For $f=f_C$ use the full $u_{z_C}$ of (41), with all
fine eliminated degrees of freedom in $C_m$.
\paragraph{Corollary 7.3.}
This is a nonzero physical vector with exact norm and energy
(42)--(43). Quantitatively,
\begin{gather}
 d_C\leq\|u_{z_C}\|^2\leq d_C+\alpha c_H/4,\nonumber\\
 \|\mathbf1_{[\Omega_C,\infty)}(H-E_0)u_{z_C}\|^2
 \geq\frac{c_H}{4+c_H}\|u_{z_C}\|^2,\label{eq:dressedweight}\\
 \frac{\Omega_Cc_H}{4+c_H}
 \leq\frac{\mathfrak q_{H-E_0}[u_{z_C}]}{\|u_{z_C}\|^2}
 \leq\kappa\ell\frac{1-r_C/4}{d_C}.
\end{gather}

\paragraph{Proof.}
By (23), each used coarse edge contributes at
most $4G\kappa m$ to the pointwise absolute value of $B_mf_C$;
there are $r$ such edges and $|D_\gamma W_C|\leq1$.
Thus $\|B_mf_C\|\leq4G\kappa\ell$.
Since $C_m\geq0$, its shifted inverse has norm at most $1/z_C$.
Consequently
\[
 \|w_{z_C}\|\leq4G/\tau\leq\tfrac12\sqrt{\alpha c_H}
                              \leq\tfrac12\sqrt{d_Cc_H}.
\]
Orthogonality gives the norm bound. Multiplication by $\psi$ is unitary,
and an orthogonal spectral projection is contractive. Hence the triangle
inequality and (53) give
\[
 \|\mathbf1_{[\Omega_C,\infty)}(H-E_0)u_{z_C}\|
 \geq\sqrt{d_Cc_H}-\|w_{z_C}\|
 \geq\tfrac12\sqrt{d_Cc_H}.
\]
The retained norm satisfies
$\|u_{z_C}\|^2=d_C+\|w_{z_C}\|^2\leq d_C(1+c_H/4)$,
so the squared lower bound $d_Cc_H/4$ is at least
$c_H\|u_{z_C}\|^2/(4+c_H)$.
Integrating $\omega\geq\Omega_C$ over this weight proves the lower
energy bound; (43) proves the upper bound.
The large positive shift here is explicitly fixed in (55);
this argument does not assert the same lower bound for the separately
defined vector minimizing the energy numerator at $z=0$. Its full
Rayleigh variational problem is calculated in Section~8.




\section{Zero-shift reconstruction and affine physical states}
\label{sec:zero}

We now calculate the zero-shift channel without transferring the
positive-shift estimate by assertion. Gauge invariance is imposed
throughout this section. Let $\HH^{\rm G}$ be the closed subspace of
gauge-invariant functions in $\HH$, and let $\KK_m^{\rm G}$ be the
corresponding coarse subspace. The gauge group is the compact product
$\SU(2)^{\mathsf V_L}$; its Haar average is an orthogonal projection.
The intertwining of the path products and the gauge descent of the
conditional density show that $J_m,J_m^*,P_m,Q_m$ commute with the
corresponding gauge actions. Consequently $B_m,C_m$ and their
resolvents preserve these subspaces. All physical vectors below are
obtained by the same unitary multiplication by $\psi$.

Take $m\geq2$ dividing $2L$. In this section write $C$ for the restriction
of $C_m$ to $\HH_{Q_m}^{\rm G}$, retaining precisely its form domain
$\VV_{Q_m}\cap\HH^{\rm G}$. This is a reducing restriction of the
operator, not a change of interaction. Put
\[
 c=\inf\sigma(C),\qquad
 \delta_{\rm G}=\inf_{\substack{0\ne u\in\VV\cap\HH^{\rm G}\\\mu u=0}}
                   \frac{q(u,u)}{\|u\|^2}.
\]
This complementary physical space is nonzero. For example, the fine
face at $n=(-L,-L,-L)$ in the first two directions contains the
direction-2 edge with first coordinate $-L+1$, which is outside $R_m$.
Varying that edge with every other fine edge fixed leaves every coarse
product unchanged and changes the fine face trace. Positivity of the
conditional density implies $Q_mW_p\ne0$, and this function is gauge
invariant. Compact embedding of the restricted form domain into $L^2$
gives compact resolvent for $C$: a bounded sequence in its form norm
is bounded in the original $H^1$ norm. Thus $c$ is an eigenvalue.

\paragraph{Proposition 8.1.}
With the original number of fine edges $N=3(2L)(2L+1)^2$, define
\begin{equation}\label{eq:finitegap}
 \Delta_L=\frac{3\kappa}{4}\alpha^N
          =\frac{3g^2}{2a}\alpha^{\,3(2L)(2L+1)^2}.
\end{equation}
Then $\delta_{\rm G}\geq\delta\geq\Delta_L>0$ and $c\geq\delta_{\rm G}$.
For mean-zero $f\in\VV_m\cap\KK_m^{\rm G}$, define the linear map
\begin{equation}\label{eq:Rzero}
 \mathcal R_m f=J_mf-C^{-1}B_mf,\quad
 s_{m,0}(f,g)=a_m(f,g)-\langle B_mf,C^{-1}B_mg\rangle.
\end{equation}
The form $s_{m,0}$ is closed on the indicated retained form space,
with
\begin{equation}\label{eq:zerocoercivity}
 \Delta_L\|f\|^2\leq s_{m,0}(f,f)\leq a_m(f,f),\qquad
 a_m(f,f)\leq2(1+\eta_m/\Delta_L)s_{m,0}(f,f).
\end{equation}
The map $\mathcal R_m$ takes this space into the original physical
mean-zero form domain and satisfies $J_m^*\mathcal R_mf=f$.
For every $h\in\VV_{Q_m}\cap\HH^{\rm G}$,
\begin{equation}\label{eq:zeroorthogonal}
 q(\mathcal R_m f,h)=0,\qquad
 q(\mathcal R_m f,\mathcal R_m g)=s_{m,0}(f,g).
\end{equation}
For its associated mean-zero operator $S_{m,0}$, the exact domain map is
\begin{equation}\label{eq:zerodomain}
 f\in D(S_{m,0})\ \Longleftrightarrow
 \mathcal R_m f\in D(\LL),\qquad
 \LL\mathcal R_m f=J_mS_{m,0}f.
\end{equation}
The equivalence in (62) is for $f$ already in the
retained mean-zero form domain.

\paragraph{Proof.}
Change fine variables one at a time between any two configurations.
Theorem~2.2 gives
$\rho_{\max}/\rho_{\min}\leq R^{2N}=\alpha^{-N}$.
The product Haar Laplacian has first positive eigenvalue $3/4$:
tensor products of the complete matrix-coefficient bases in
Lemma~2.1 have eigenvalues
$\sum_e j_e(j_e+1)$, whose least positive value is $3/4$.
Expanding a smooth function in that basis, and then using $H^1$
density, gives the Haar Poincare inequality. Hence
\[
 \Var_\mu u
 \leq\rho_{\max}\Var_\lambda u
 \leq\frac{4\rho_{\max}}3\sum_{e,a}\int|X_{e,a}u|^2\dd\lambda
 \leq\frac{4\rho_{\max}}{3\kappa\rho_{\min}}q(u,u).
\]
The first inequality evaluates the infimum over constants defining
$\Var_\mu$ at the Haar mean. This proves (58) for all
mean-zero functions, and restricting to gauge-invariant functions
preserves the lower bound. Every complementary physical function has
zero mean, so $c\geq\delta_{\rm G}$.

The bounded resolvent $C^{-1}$ maps its Hilbert space to $D(C)$ and
therefore to the form domain. Completing the square at zero gives
\[
 q(J_mf+h,J_mf+h)
 =s_{m,0}(f,f)+c_m(h+C^{-1}B_mf,h+C^{-1}B_mf).
\]
The unique minimizer is $\mathcal R_m f$, proving
(61) and the upper bound in
(60). Its mean is zero and its squared norm is
$\|f\|^2+\|C^{-1}B_mf\|^2$. The already proved full Poincare inequality
gives the lower bound and
$\|\mathcal R_m f\|^2\leq s_{m,0}(f,f)/\Delta_L$.
Apply (27) to $\mathcal R_m f$ for the second
inequality in (60). These two-sided form-norm
bounds prove closedness with exactly the retained $H^1$ domain.

If $f\in D(S_{m,0})$, decompose any physical form test vector as
$J_mg+h$. Equations (61) give
\[
 q(J_mg+h,\mathcal R_m f)=s_{m,0}(g,f)
                       =\langle J_mg+h,J_mS_{m,0}f\rangle .
\]
Constants give zero on both sides. This is the defining weak operator
identity. It also holds against arbitrary, noninvariant test vectors:
average that test vector over the gauge group, since the operator and
the vector being tested are invariant. Thus $\mathcal R_m f$ belongs
to the full $D(\LL)=H^2$ and has the displayed image.
Conversely if $\mathcal R_m f\in D(\LL)$, tests in the complementary
space show $Q_m\LL\mathcal R_m f=0$; retained tests then prove
$f\in D(S_{m,0})$ and the claimed image.



The exponential dependence on $N$ in (58) is explicit.
It provides a proved finite-box inverse bound, not a volume-uniform gap.

Fix a nonzero, mean-zero, gauge-invariant $f\in\VV_m$. Keep its norm
unchanged and set
\[
 d_f=\|f\|^2,\quad a_f=a_m(f,f),\quad b_f=B_mf,\quad
 \sigma_f(I)=\|\mathbf1_I(C)b_f\|^2 .
\]
This $\sigma_f$ is a finite positive measure on $[c,\infty)$ of total
mass $\|b_f\|^2$, with no division by that mass. For every real $z>-c$,
put
\begin{equation}\label{eq:inversemoments}
 M_j(z)=\int_{[c,\infty)}(t+z)^{-j}\dd\sigma_f(t),\quad j=1,2,3,
 \qquad v_z=J_mf-(C+z)^{-1}b_f .
\end{equation}
All these moments are finite by $t+z\geq c+z>0$.

\paragraph{Theorem 8.2.}
The actual physical state $\psi v_z$ has exact squared norm, excitation
energy and Rayleigh quotient
\begin{equation}\label{eq:affinemoments}
 N_f(z)=d_f+M_2(z),\quad
 e_f(z)=a_f-M_1(z)-zM_2(z),\quad
 Q_f(z)=\frac{e_f(z)}{N_f(z)}.
\end{equation}
For all $z>-c$,
\begin{equation}\label{eq:rayleighderivative}
 Q_f'(z)=\frac{2M_3(z)}{N_f(z)}[z+Q_f(z)].
\end{equation}
In particular $Q_f$ is strictly increasing on $[0,\infty)$ whenever
$b_f\ne0$. Its zero-shift value is strictly smaller than the lifted
quotient $a_f/d_f$, and some negative shifts give a quotient strictly
smaller still. Zero shift minimizes the energy numerator at fixed
retained data, but it does not minimize the physical quotient when
$b_f\ne0$.

The full affine variational problem, with all physical eliminated
states allowed, is determined as follows:
\begin{gather}
 \chi_f=\inf_{h\in\VV_{Q_m}\cap\HH^{\rm G}}
           \frac{q(J_mf+h,J_mf+h)}{d_f+\|h\|^2},\nonumber\\
 F_f(E)=a_f-E d_f-M_1(-E),\qquad 0\leq E<c,\qquad
 F_f'(E)=-d_f-M_2(-E)<0,\label{eq:affineroot}\\
 F_f(0)>0,\quad
 \mathfrak l_f=\lim_{E\uparrow c}F_f(E)\in[-\infty,\infty).\nonumber
\end{gather}
If $\mathfrak l_f<0$, there is exactly one root $E_f\in(0,c)$ and
\begin{gather}
 \chi_f=E_f,\qquad
 h_f=-(C-E_f)^{-1}b_f,\qquad
 \|\psi(J_mf+h_f)\|^2=-F_f'(E_f),\nonumber\\
 \mathfrak q_{H-E_0}[\psi(J_mf+h_f)]=-E_fF_f'(E_f).
 \label{eq:affineattain}
\end{gather}
This is the unique affine minimizer. If $\mathfrak l_f\geq0$, then
$\chi_f=c$. In this case $b_f$ is orthogonal to $\ker(C-c)$.
The value $c$ is attained exactly when $\mathfrak l_f=0$, and all
minimizers are
\begin{equation}\label{eq:boundaryattain}
 h=-(C-c)^\dagger b_f+k,\qquad k\in\ker(C-c).
\end{equation}
Here the dagger is the inverse on $\ker(C-c)^\perp$ and zero on the
kernel. When $\mathfrak l_f>0$, no finite affine vector attains $c$;
the infimum is approached by $h=t k$, $t\to\infty$, for a nonzero
$k\in\ker(C-c)$. These alternatives are exhaustive for the actual
interacting $C$ and $b_f$, including $b_f=0$.

\paragraph{Proof.}
The spectral theorem identifies $M_1$ with
$\langle b_f,(C+z)^{-1}b_f\rangle$ and $M_2$ with
$\|(C+z)^{-1}b_f\|^2$. The orthogonal decomposition of $v_z$ and
the block form (32) prove (64)
for negative as well as positive $z$. Differentiating these bounded
resolvents gives $M_1'=-M_2$, $M_2'=-2M_3$,
$N_f'=-2M_3$, and $e_f'=2zM_3$. The quotient rule proves
(65). Each $v_z$ is nonzero and mean-zero,
so $Q_f(z)\geq\delta_{\rm G}>0$.
For $b_f\ne0$, $M_3>0$ and the derivative is positive at zero
and at every positive shift. Also
\[
 \frac{a_f}{d_f}-Q_f(0)
 =\frac{M_1(0)+(a_f/d_f)M_2(0)}{d_f+M_2(0)}>0 .
\]
Continuity of the derivative at zero proves strict improvement for
sufficiently small negative shifts. The statement about the numerator
follows independently from (61); increasing the
norm is part of the exact quotient calculation.

For $0\leq E<c$, complete the shifted square for arbitrary $h$:
\begin{equation}\label{eq:affinesquare}
 q(J_mf+h,J_mf+h)-E(d_f+\|h\|^2)
 =F_f(E)+\|(C-E)^{1/2}[h+(C-E)^{-1}b_f]\|^2.
\end{equation}
No discarded term or approximation enters this identity.
The sign and derivative statements in (66) follow
from Proposition~8.1 and the spectral derivative.
Monotonicity bounds the limit above by $F_f(0)$ and allows the value
$-\infty$ below. The intermediate value theorem proves the unique
root when the limit is negative. At this root
(69) bounds every quotient below by $E_f$,
and equality holds precisely for the stated $h_f$. The norm and
energy in (67) follow from
(64) and $F_f(E_f)=0$.

If the limit is nonnegative, $F_f(E)>0$ for every $E<c$.
Equation (69), followed by $E\uparrow c$,
bounds every affine quotient below by $c$. A complementary
eigenvector $k$ of eigenvalue $c$ yields
\[
 \frac{a_f+2t\Rea\langle b_f,k\rangle+t^2c\|k\|^2}
      {d_f+t^2\|k\|^2}\longrightarrow c,
\]
proving equality of the infimum. If $b_f$ had a nonzero projection
on that eigenspace, its contribution to $M_1(-E)$ would diverge as
$(c-E)^{-1}$, contradicting $\mathfrak l_f\geq0$.
Compact resolvent isolates $c$ from the rest of $\sigma(C)$, so
$(C-c)^\dagger$ is bounded on the complementary eigenspaces and
maps into $D(C)$. Completing the square at $E=c$ now gives
\[
 q(J_mf+h,J_mf+h)-c(d_f+\|h\|^2)
 =\mathfrak l_f+
 \|(C-c)^{1/2}[h+(C-c)^\dagger b_f]\|^2 .
\]
It proves both the attainment criterion and every vector in
(68).



The scalar root in (66) solves the affine problem for
this fixed retained direction. It is not, in general, an eigenvalue of
the full operator: a full eigenvector also requires the weak equation
$s_{m,-E}(g,f)=0$ for every retained test direction $g$.
Here $s_{m,-E}$ denotes the sesquilinear expression
(34) at $z=-E$, restricted to physical retained functions;
it is well defined on $\VV_m\cap\KK_m^{\rm G}$ because $E<c$.
Indeed complementary testing gives $h=-(C-E)^{-1}B_mf$, and retained
testing gives exactly that additional equation. Conversely these
equations imply the full weak eigenvector equation and its $H^2$
operator domain. This proves the exact map in both directions, and
explains precisely which test directions the affine minimization uses.
The energy-dependent eigenvalue and reconstruction principle agrees
with \href{https://doi.org/10.4171/ECR/18-1/5}{[1, Theorem~1.2 and Corollary~1.3]}; the affine classification
and the physical domains above were proved directly.

\paragraph{Corollary 8.3.}
For the loop $f=f_C$, the zero-shift state
$u_0=\psi\mathcal R_m f_C$ obeys the exact identities
\begin{align}
 \|u_0\|^2&=d_C+M_2(0),\nonumber\\
 \mathfrak q_{H-E_0}[u_0]&=\kappa\ell(1-r_C/4)-M_1(0).
 \label{eq:zeroloopactual}
\end{align}
\begin{gather}
 \alpha\leq d_C\leq\|u_0\|^2
 \leq d_C+\min\left\{
       \frac{\kappa\ell(1-r_C/4)}{\Delta_L},
       \frac{16G^2\kappa^2\ell^2}{\Delta_L^2}\right\},
 \label{eq:zeroloopnorm}\\
 \|u_0\|^2\leq4+\frac{4\ell}{3\alpha^N},\nonumber\\
 \Delta_L\leq
 \frac{\mathfrak q_{H-E_0}[u_0]}{\|u_0\|^2}
 \leq\frac{\kappa\ell(1-r_C/4)}{d_C}.\label{eq:zeroloopenergy}
\end{gather}
For an elementary coarse face, the last norm bound is
$4+16m/(3\alpha^N)$, and its energy numerator is exactly
$8g^2m(1-r_C/4)/a-M_1(0)$. The same $\Delta_L=3g^2\alpha^N/(2a)$
bounds the quotient below. These bounds are finite and explicit;
their volume dependence is retained.

\paragraph{Proof.}
The identities are (64) at zero.
Positivity gives $M_1(0)\leq a_f$, and the spectral inequality
$t^{-2}\leq c^{-1}t^{-1}$ gives
$M_2(0)\leq a_f/c\leq a_f/\Delta_L$.
The other estimate uses $M_2(0)\leq\|B_mf_C\|^2/\Delta_L^2$
and the already proved $\|B_mf_C\|\leq4G\kappa\ell$.
Insert $a_f=\kappa\ell(1-r_C/4)\leq\kappa\ell$ and
$d_C\leq4$. The lower quotient bound is the full Poincare inequality,
and the upper one follows from the exact numerator and norm.
Finally use $\ell=4m$ and $\kappa=2g^2/a$.



\section{Exact comparison of all physical low-energy states}
\label{sec:allsoft}

The previous section treated a fixed retained observable. We now give
a quantitative relation encompassing all physical states, with no
assumed complementary gap uniform in volume. Define the actual
zero-shift reconstructed infimum
\begin{equation}\label{eq:reconstructedgap}
 \lambda_m^{\rm R}
 =\inf_{\substack{0\ne f\in\VV_m\cap\KK_m^{\rm G}\\\mu_m f=0}}
 \frac{s_{m,0}(f,f)}
      {\|f\|^2+\|C^{-1}B_mf\|^2}.
\end{equation}
This denominator is the full physical norm. In every fixed box it is
positive, and $\lambda_m^{\rm R}\geq\delta_{\rm G}\geq\Delta_L$.

\paragraph{Theorem 9.1.}
For every allowed $L,m,a,g$ with $m\geq2$, the following comparison
holds with numerical constants independent of all four parameters:
\begin{equation}\label{eq:harmonicgap}
 \frac{\lambda_m^{\rm R}c}{\lambda_m^{\rm R}+c}
 \leq\delta_{\rm G}\leq\min\{\lambda_m^{\rm R},c\},\qquad
 \frac12\min\{\lambda_m^{\rm R},c\}\leq\delta_{\rm G}.
\end{equation}
It has the following explicit state map. For every mean-zero physical
$u\in\VV$, put
\begin{equation}\label{eq:softdecomposition}
 f=J_m^*u,\quad v=\mathcal R_m f,\quad
 k=u-v=Q_mu+C^{-1}B_mf .
\end{equation}
Then $k\in\VV_{Q_m}\cap\HH^{\rm G}$ and
\begin{equation}\label{eq:softidentities}
 q(u,u)=s_{m,0}(f,f)+c_m(k,k),\qquad
 \|u\|^2=\|v\|^2+\|k\|^2+2\Rea\langle v,k\rangle .
\end{equation}
At least one of the nonzero components $v,k$ has physical Rayleigh
quotient at most $2q(u,u)/\|u\|^2$. Conversely each such component is an
actual admissible physical state under multiplication by $\psi$.

For every sequence of allowed boxes, block lengths, positive spacings
and couplings, $\delta_{\rm G}$ tends to zero if and only if
$\min\{\lambda_m^{\rm R},c\}$ tends to zero. Uniform positive lower
bounds for $\delta_{\rm G}$ are equivalent to simultaneous uniform
positive lower bounds for these two actual sector infima. These are
equivalences proved by (74), not assumed estimates
for either sector.

\paragraph{Proof.}
Sobolev preservation of $J_m^*$ and Proposition~8.1
make all components in (75) well defined in
their stated form domains. The identity $J_m^*v=f$ puts $k$ in the
complement. Equation (61) gives $q(v,k)=0$,
which proves the energy identity; the norm identity retains its
cross term. Cauchy--Schwarz in two scalar components gives
\[
 \|u\|\leq\|v\|+\|k\|
 \leq\sqrt{\frac{q(v,v)}{\lambda_m^{\rm R}}}
      +\sqrt{\frac{c_m(k,k)}c}
 \leq\sqrt{\left(\frac1{\lambda_m^{\rm R}}+\frac1c\right)q(u,u)}.
\]
Zero components contribute zero. Taking the infimum over $u$ proves
the first lower bound. The upper bound follows because the full
variational set contains both the reconstructed vectors and all
complementary physical vectors. For positive $x,y$,
$xy/(x+y)\geq\min(x,y)/2$, proving the last bound.

For the explicit factor-two state assertion,
$\|u\|^2\leq2(\|v\|^2+\|k\|^2)$ and the energy identity show
\[
 \min_{\substack{w\in\{v,k\}\\w\ne0}}
       \frac{q(w,w)}{\|w\|^2}
 \leq\frac{q(v,v)+q(k,k)}{\|v\|^2+\|k\|^2}
 \leq2\frac{q(u,u)}{\|u\|^2}.
\]
This construction changes no physical coefficient or state norm.
All assertions about sequences and uniform bounds follow directly
from the two-sided inequalities for each member of the sequence.



This comparison calculates how a putative collective low-energy sector
must appear under the exact block map: it appears either among the
full zero-shift reconstructed states, with their eliminated norm
retained, or among physical states in the complementary sector.
The decomposition (75) exhibits the vectors
and controls their energy quotients by the numerical factor two.
It does not assert which alternative occurs in the interacting
large-volume limit, or that either one actually becomes soft.
The fixed-loop high-energy fraction in Theorem~7.2
continues to hold in its full scope. Equations
(70)--(72) and
Theorem~9.1 now supply the proved zero-shift and
all-state relations without extending that fraction to an unsupported
family of vectors. For $m=1$, $J_1$ is the identity, $Q_1=B_1=0$,
and reconstruction is the identity; the physical gap relation is then
exactly $\lambda_1^{\rm R}=\delta_{\rm G}$ with no complementary sector.

\section{Physical skeleton morphism and uniform complementary inverse}
\label{sec:uniforminverse}

This section uses the complete fixed-point and physical-resolvent proof
retained in \nolinkurl{references/LANE_A_ELIMINATION_INPUTS.md},
Sections 1, 10, 11, 12 and 15. The input concerns the identical open
box, every original link and face, and the actual positive unit vacuum.
Here is its exact dictionary. Its $g_{\rm YM}$ is our $g$, its
$\mathcal E$ is our $E_0$, and its coordinate action
$\alpha_h(U)_e=h_{s(e)}U_eh_{t(e)}^{-1}$ obeys
\[
 (h\cdot U)_e=h_{s(e)}^{-1}U_eh_{t(e)}
                  =(\alpha_{h^{-1}}U)_e.
\]
Thus its Hilbert-space representation
$\mathsf T_hu=u\circ\alpha_{h^{-1}}$ is exactly
$\mathsf T_hu(U)=u(h\cdot U)$ here. The identity map on functions
intertwines these representations, on $L^2$, $H^1$ and $H^2$;
no assertion that parameter inversion is a group homomorphism is needed.
The two Hamiltonians and their domains coincide. Uniqueness of the
positive unit vacuum then proves equality of the two $\psi$, of
$\rho=\psi^2$, and of the ground energies. The source scalar in its
(12.7) will be written $c_{\rm vac}=\langle1,\psi\rangle$ here;
it is not the density constant $\alpha$ in (58).

For precision, the proved input is
\begin{equation}\label{eq:stronggap}
 \delta_{\rm G}\geq\Delta_{\rm sc}
 :=3\kappa\left(1-\frac{512}{3}\xi\right)
 =\kappa(3-512\xi)\geq\frac{287\kappa}{96},
 \qquad 0<\xi\leq\frac1{49152}.
\end{equation}
Its entire proof is retained, not a proposed additional hypothesis.
The exact source construction maps a collection of nonconstant
Haar sectors $z_I$ to $\phi=e^{\mathfrak C}1$, where
$\mathfrak C=\sum_{I\ne\varnothing}|z_I\rangle\langle1_I|
\otimes I_{I^c}$ is bounded and nilpotent in each finite box.
The inverse is $e^{-\mathfrak C}$, with both exponentials finite
polynomials preserving $D(H_0)=H^2$. Its full fixed-point proof gives
\[
 \psi=c_{\rm vac}\phi,\quad
 c_{\rm vac}>0,\quad c_{\rm vac}^2\|\phi\|^2=1,\quad
 e^{-\mathfrak C}(H-E_0)e^{\mathfrak C}=H_0+\mathcal K,
 \qquad H_0=\kappa\sum_eE_e.
\]
These maps preserve the vacuum scalar and all of $2bM$ in the
eigenvalue equation. The source proves that every iterate, hence both
similarities, commutes with every vertex gauge transformation.
In the orthogonal sector decomposition a physical nonempty support
has no degree-one vertex: averaging at such a vertex would give both
the same invariant vector and its zero one-link Haar average.
Every nonempty surviving support therefore contains at least four
edges of this open bipartite cubic graph. With the original
$E_e$ Casimir, $H_0\geq3\kappa$ on the physical nonconstant sectors.
The complete interaction estimate proved there is
$\|\mathcal K f\|_{\oplus,1}\leq(512\xi/3)\|H_0f\|_{\oplus,1}$,
where $\|f\|_{\oplus,1}=\sum_I\|\Pi_If\|$ and
$\|f\|\leq\|f\|_{\oplus,1}\leq2^{N/2}\|f\|$.
The retained proof supplies every commutator and counting estimate
and the contraction on $\|z\|_*\leq1/16$, with constant at most
$11/36$ in the displayed range. The sector norm equivalence is
explicitly volume dependent. The physical Neumann resolvent excludes
every $0<t<3\kappa(1-512\xi/3)$, transfers under the displayed
domain-preserving similarity, and proves (77) for the
original self-adjoint operator. No form positivity of a nonunitary
similarity or deletion of its norm factors is used.

The retained proof derives numerical constants for the creation method
of Yarotsky~\href{https://arxiv.org/abs/math-ph/0411042v1}{[7]}, Section 2, source equations
\texttt{mmo}, \texttt{c1}, \texttt{vsum}, \texttt{thl},
\texttt{sumj}. Its physical-space reference is Baez~\href{https://arxiv.org/abs/gr-qc/9411007v1}{[8]},
the section ``Gauge Theory on a Graph,'' Lemmas 1, 2 and the source
lemma \texttt{lem2.5}. Baez's link convention
$A_e\mapsto h_{t(e)}A_eh_{s(e)}^{-1}$ is related to the input's
$U_e\mapsto h_{s(e)}U_eh_{t(e)}^{-1}$ by the involution
$\jmath(U)_e=U_e^{-1}$: direct multiplication gives
$\jmath(\alpha_hU)_e=h_{t(e)}\jmath(U)_eh_{s(e)}^{-1}$.
The pullback is a Haar-unitary map with itself as inverse,
preserves $H^1,H^2$ because inversion is a bi-invariant metric
isometry, and intertwines each $E_e$. The raw Wilson products in
our Hamiltonian are unchanged. Both original TeX sources are retained.

The source conditional complement retains individual fine links.
Our complement retains their ordered products. The following proves
the exact relation, including the Hilbert and operator domains.

\paragraph{Proposition 10.1.}
Let $2\leq m\mid2L$, $s=(U_e)_{e\in S_m}$,
$r=(U_e)_{e\in R_m}$, and define
\[
 \mathcal A_m=\SU(2)^{R_m},\quad
 \bar\rho_R(r)=\int\rho(s,r)\dd\lambda_{S_m}(s),\quad
 \mathcal K_R=L^2(\mathcal A_m,\bar\rho_R\lambda_R).
\]
Let $J_S:\mathcal K_R\to\HH$ be $J_SF(s,r)=F(r)$, with
\[
 (J_S^*u)(r)=\frac{1}{\bar\rho_R(r)}
                \int\rho(s,r)u(s,r)\dd\lambda_{S_m}(s),
 \qquad P_S=J_SJ_S^*.
\]
Write $\mathcal K_R^{\rm G}$ for invariance under the endpoint action
of all original vertices. There is a unitary map of physical spaces
\begin{equation}\label{eq:skeletonunitary}
 \iota_m:\KK_m^{\rm G}\longrightarrow\mathcal K_R^{\rm G},
 \qquad (\iota_m f)(r)=f((u_{\gamma,1}\cdots u_{\gamma,m})_\gamma).
\end{equation}
Its inverse on arbitrary $L^2$ classes is
\begin{equation}\label{eq:skeletoninverse}
 (\iota_m^{-1}F)(z)=\int F(\Theta_m^{-1}(v,z))\dd\lambda_v(v),
 \quad
 \Theta_m:\mathcal A_m\longrightarrow
 \left(\prod_\gamma\SU(2)^{m-1}\right)\times\mathcal Q_m,
 \quad r\longmapsto(v,z).
\end{equation}
The inverse coordinates are $u_{\gamma,k}=v_{\gamma,k}$ for $k<m$
and $u_{\gamma,m}=(v_{\gamma,1}\cdots v_{\gamma,m-1})^{-1}z_\gamma$.
Both maps identify the invariant $H^1$ and $H^2$ spaces, and
\begin{equation}\label{eq:physicalprojections}
 J_S\iota_m=J_m,
 \qquad J_S^*|_{\HH^{\rm G}}=\iota_mJ_m^*|_{\HH^{\rm G}},
 \qquad P_S|_{\HH^{\rm G}}=P_m|_{\HH^{\rm G}}.
\end{equation}
Consequently the identity on functions is an isometry from the
source physical conditional complement for $S=S_m$ onto
$\HH_{Q_m}^{\rm G}$. It identifies their full form domains,
their complementary operator domains, their operators and their
resolvents. In the range (77),
\begin{equation}\label{eq:uniformC}
 C\geq\Delta_{\rm sc} I,\qquad
 \|C^{-1}\|\leq\Delta_{\rm sc}^{-1}
                  \leq\frac{96}{287\kappa}.
\end{equation}


\paragraph{Proof.}
The two inverse coordinate compositions and the Haar identity for
$\Theta_m$ follow by the exact prefix cancellations already proved
for $\Phi_m$, with all $S_m$ factors omitted from both sides of
this particular map. Integrating those factors in the density gives
\[
 \bar\rho_R(\Theta_m^{-1}(v,z))=r_m(z),\qquad
 \bar\rho_R\dd\lambda_R=r_m(z)\dd\lambda_v\dd\lambda_m.
\]
To prove the first equality, perform simultaneously the internal
vertex transformations with fixed coarse endpoints as in Section 3;
they send $(v,z)$ to $(1,z)$, preserve the marginal by substitution
on $S_m$, and leave $z$ unchanged. Integrating in $v$ then identifies
the remaining function with the definition of $r_m$.

There is also an $L^2$ proof of the descent, without evaluation on
a measure-zero section. For a chain with endpoints fixed, the
internal gauge action is
$v'_1=u_1h_1$ and $v'_k=h_{k-1}^{-1}u_kh_k$ for $1<k<m$.
For fixed input $r$, Haar integration successively over
$h_1,\ldots,h_{m-1}$ makes the variables $v'_1,\ldots,v'_{m-1}$
product Haar; the chain product remains $z$. Applying this to every
disjoint interior shows that the compact internal gauge average on
$L^2(\bar\rho_R\lambda_R)$ is precisely fibre integration in $v$
followed by constant lift. Hence every invariant $F$ equals that
lift almost everywhere. The remaining coarse endpoint action is
$z_\gamma\mapsto h_x^{-1}z_\gamma h_{x+m\mathbf e_i}$; this proves
the full coarse invariance of (79) and the
fine invariance of (78). Substitution proves
both inverse identities. Moreover
\[
 \|\iota_m f\|_{\mathcal K_R}^2
 =\int |f(z)|^2r_m(z)\dd\lambda_v\dd\lambda_m
 =\|f\|_{\KK_m}^2.
\]
No vector is divided by its mass. Smooth coordinate pullback on
these compact products preserves each Sobolev order. Differentiating
under the compact fibre integral and applying Cauchy--Schwarz bounds
the inverse in $H^1$ and $H^2$. Smooth positive weights give the
same Sobolev sets in each finite box. Approximation, followed by
compact gauge averaging, proves the claimed domain identifications.

The first identity in (80) follows from
the coordinate product. For physical $u$, substitution in the
conditional integral shows that $J_S^*u$ is invariant under every
skeleton endpoint action. In $(v,z)$ coordinates it is therefore
independent of $v$ almost everywhere. Its $v$ integral equals
$r_m(z)^{-1}\int\rho(s,\Theta_m^{-1}(v,z))
u(s,\Theta_m^{-1}(v,z))\dd\lambda_{S_m}\dd\lambda_v$,
which is exactly $J_m^*u$. This proves the second identity and then
the third. It does not assert equality of the projections on all
noninvariant vectors.

Thus the two physical complements are the same closed subspace
of $\HH$, with the same form $q$, including derivatives in every
fine link. Both have form domain $H^1$ intersected with that space.
The weak operator-domain rule is, in both presentations,
\[
\begin{aligned}
 D(C)=\bigl\{u\in H^1\cap\HH_{Q_m}^{\rm G}:{}&
 \text{ there exists }w\in\HH_{Q_m}^{\rm G}\text{ such that}\\
 &q(v,u)=\langle v,w\rangle
 \text{ for every }v\in H^1\cap\HH_{Q_m}^{\rm G}\bigr\},\\
 Cu={}&w.
\end{aligned}
\]
Equality of the tests, inner products and forms proves equality
of the operators and their resolvents, not merely of their formal
differential expressions. Finally every vector in the complement
has $\mu u=0$, and $\mathcal Uu=\psi u$ is physical and perpendicular
to $\psi$. The already-proved (77) applied under
this unchanged unitary gives (81).



\paragraph{Theorem 10.2.}
In the range (77), for every mean-zero
$f\in\VV_m\cap\KK_m^{\rm G}$, put $b_f=B_mf$,
$d_f=\|f\|^2$, $a_f=a_m(f,f)$ and retain the raw measure
$\sigma_f(I)=\|1_I(C)b_f\|^2$.
Then
\begin{equation}\label{eq:uniformzeroenergy}
 s_{m,0}(f,f)\geq\Delta_{\rm sc}
              (d_f+\|C^{-1}b_f\|^2),\qquad
 \lambda_m^{\rm R}\geq\Delta_{\rm sc},\quad c\geq\Delta_{\rm sc}.
\end{equation}
For the unchanged mean-centered simple loop $f_C$, let
$\ell$ be its fine perimeter and $d_C,r_C$ retain their earlier
definitions. The physical zero-shift state $u_0=\psi\mathcal R_m f_C$
has the exact mass and energy
\begin{align}
 \|u_0\|^2&=d_C+\int t^{-2}\dd\sigma_{f_C}(t),\label{eq:uniformloopmass}\\
 \mathfrak q_{H-E_0}[u_0]
 &=\kappa\ell(1-r_C/4)-\int t^{-1}\dd\sigma_{f_C}(t)
 \geq\Delta_{\rm sc}\|u_0\|^2.\label{eq:uniformloopenergy}
\end{align}
Its quantitative upper mass bound is
\begin{equation}\label{eq:uniformloopbound}
 \|u_0\|^2\leq d_C+
 \min\left\{\frac{\kappa\ell(1-r_C/4)}{\Delta_{\rm sc}},
             \frac{16G^2\kappa^2\ell^2}{\Delta_{\rm sc}^2}\right\}
 \leq4+\frac{96\ell}{287}.
\end{equation}
For every $z\geq0$, define $w_z=(C+z)^{-1}b_f$,
$N_f(z)=d_f+\|w_z\|^2$ and
$e_f(z)=a_f-\int(t+z)^{-1}\dd\sigma_f
                    -z\int(t+z)^{-2}\dd\sigma_f$.
These are the full physical mass and excitation energy of
$\psi(J_mf-w_z)$. The exact differences are
\begin{align}
 w_0-w_z&=zC^{-1}(C+z)^{-1}b_f,\label{eq:uniformwdiff}\\
 e_f(z)-e_f(0)&=\int\frac{z^2}{t(t+z)^2}\dd\sigma_f(t),
                                                        \label{eq:uniformediff}\\
 N_f(0)-N_f(z)&=\int\left(\frac1{t^2}-\frac1{(t+z)^2}\right)
                                      \dd\sigma_f(t),    \label{eq:uniformndiff}\\
 s_{m,z}(f,f)-s_{m,0}(f,f)-zd_f
 &=\int\frac{z}{t(t+z)}\dd\sigma_f(t).    \label{eq:uniformmemorydiff}
\end{align}
Writing $\theta=z/(\Delta_{\rm sc}+z)$, their estimates are
\begin{align}
 \|w_0-w_z\|&\leq\theta\sqrt{a_f/\Delta_{\rm sc}},\nonumber\\
 0\leq e_f(z)-e_f(0)&\leq\theta^2a_f,\nonumber\\
 0\leq N_f(0)-N_f(z)&\leq
 \left[1-\left(\frac{\Delta_{\rm sc}}{\Delta_{\rm sc}+z}\right)^2\right]
 \frac{a_f}{\Delta_{\rm sc}},\nonumber\\
 0\leq s_{m,z}(f,f)-s_{m,0}(f,f)-zd_f&\leq\theta a_f.
                                                        \label{eq:uniformcontinuity}
\end{align}
The constants controlling the inverse do not depend on box or block
size. The displayed $a_f$, $\ell$, $G$, $\kappa$, raw norms and
spectral measures keep their full dependence.


\paragraph{Proof.}
The map $\mathcal R_m$ and its exact domain and energy identities
were proved in Proposition~8.1. Its image under
$\mathcal U$ is physical and perpendicular to $\psi$. Applying
(77) proves (82), including
the infimum inequality. The previous exact inverse-moment formulas
give (83)--(84).
Positivity of $s_{m,0}$ gives
$\int t^{-1}\dd\sigma_f\leq a_f$. Since
$\operatorname{supp}\sigma_f\subset[\Delta_{\rm sc},\infty)$,
\[
 \int t^{-2}\dd\sigma_f\leq
 \min\{a_f/\Delta_{\rm sc},\|b_f\|^2/\Delta_{\rm sc}^2\}.
\]
Insert the original loop identities $a_f=\kappa\ell(1-r_C/4)$,
$\|b_f\|\leq4G\kappa\ell$, $d_C\leq4$, $a_f\leq\kappa\ell$
and (77) to prove (85).

The resolvent identity proves (86). Subtraction
of the complete energy integrands gives
$t^{-1}-(t+z)^{-1}-z(t+z)^{-2}=z^2/[t(t+z)^2]$.
The norm subtraction and the definition of $s_{m,z}$ give the
remaining exact differences. On the support of the raw measure,
$z/(t+z)\leq\theta$ and
$1-[t/(t+z)]^2\leq1-[\Delta_{\rm sc}/(\Delta_{\rm sc}+z)]^2$.
Apply these inequalities to the corresponding nonnegative
integrands and use the preceding inverse-moment bounds. Squaring
(86) first proves its norm estimate. This proves
every inequality in (90), with no removal
of the measure or its total mass $\|B_mf\|^2$.



For an elementary coarse face $\ell=4m$, the last mass bound is
$4+384m/287$, and its physical side length remains $ma$.
In original parameters the new lower energy scale and range are
\[
 \Delta_{\rm sc}=\frac{2g^2}{a}\left(3-\frac{128}{g^4}\right)
 \geq\frac{287g^2}{48a},\qquad g^4\geq12288.
\]
The expectation of $H$ is still $E_0N_f(z)+e_f(z)$.
This is a statement for the interacting finite boxes, uniform as
their sizes vary within this explicit coupling range. It takes no
spatial continuum limit and provides no small-$g$ continuation.
Every induced score covariance and the full resolvent and time memory
remain those of Sections 4--6.


\section{Original scales and the exact scope of the spectral conclusion}

For an elementary coarse face $\ell=4m$. The actual state has
\begin{align}
 \|u_C\|^2&=r_C-(\overline W_C)^2\in[\alpha,4],\label{eq:physicalnorm}\\
 \mathfrak q_{H-E_0}[u_C]
 &=\frac{8g^2m}{a}\left(1-\frac{r_C}{4}\right)
       \in\left[\frac{6\alpha g^2m}{a},\frac{8g^2m}{a}\right],\label{eq:physicalenergy}\\
 \frac{\mathfrak q_{H-E_0}[u_C]}{\|u_C\|^2}
 &\in\left[\frac{3\alpha g^2m}{2a},\frac{8g^2m}{a\alpha}\right],
 \qquad \Omega_C=\frac{3\alpha g^2m}{4a}.\label{eq:physicalquotient}
\end{align}
Here
$\alpha=1/(9\,12^{16/g^4})$, and $G,C_*,c_H$ are the explicit functions
(7), (53) of $d=4/g^4$.
The physical length is $ma$, so $g^2m/a=g^2(ma)/a^2$ as an exact
rewriting with no coordinate change. The expectation of the unsubtracted
Hamiltonian is
\[
 \langle u_C,Hu_C\rangle
 =E_0(r_C-(\overline W_C)^2)
       +\frac{8g^2m}{a}(1-r_C/4).
\]
The same scalar $E_0\|u_{z_C}\|^2$ is added to the dressed excitation
energy. Also $0\leq E_0\leq2bM$ follows from $H\geq0$ and the constant
unit Haar trial vector, whose plaquette means vanish. The original
$2bM$ term has not been omitted from $H$.

The exact instantaneous kinetic term has the two metric expressions
\[
 -\kappa m\Delta^T=-(\kappa m/4)\Delta^c,\qquad \kappa m=2g^2m/a.
\]
If its coefficient is written solely as
$2g_m^2/(ma)$, algebra forces $g_m^2=m^2g^2$. This is only the
coefficient dictionary for this marginal kinetic form. It does not make
(30) and the energy-dependent memory into the
original one-plaquette Hamiltonian at a claimed running coupling.

The map from the spatial coarse observable to the quantum theory is
exactly $f\mapsto\psi J_mf$, followed by (51); the retained
resolvent is (35), and sequential eliminations obey
(45). Theorems~7.2 and
Corollary~7.3 are statements about that full quantum
spectral measure, with no omission of generated interactions. For fixed
$g>0$, these particular large-loop channels retain a positive fraction
of weight at energies proportional to $g^2m/a$, including the specified
fully dressed channels. Their unnormalized spectral mass above the
displayed threshold is bounded below by the stated positive multiple
of their unchanged squared norm, at every increasing $m$ and fixed $a,g$.

This assertion does not give a lower bound on the support of the whole
measure: some weight can lie below $\Omega_C$. Nor does a lower bound
on the quotient of this family bound the infimum over all physical
states. At variable $g=g(a)$ the constants $\alpha$ and $c_H$ have the
explicit coupling dependence above and can tend to zero. The calculation
therefore proves neither a continuum mass gap nor a four-dimensional
mass-gap counterexample. It supplies an exact extensive block map and
an all-positive-coupling, volume-independent quantitative channel result;
Theorem~9.1 additionally relates every physical
low-energy state to the zero-shift reconstructed and complementary
sectors, with parameter-independent comparison constants and explicit
state maps. Section~10 now supplies the proved
physical lower bound $\Delta_{\rm sc}=\kappa(3-512\xi)$, uniform
in box and block size for $0<\xi\leq1/49152$, together with the
exact skeleton/complement morphism and full-memory zero-shift bounds.
The all-positive-coupling finite bound (58) remains
valid with its displayed volume dependence. Neither the comparison
nor the exhaustive affine calculation
in Theorem~8.2 presumes which sector becomes soft or an
answer for the joint spatial continuum limit.

\section{Primary-source map and reproducibility}

The general Feshbach--Schur construction is attributed to its established
literature. The locally retained source of Dusson, Sigal and Stamm was
read at Theorem~1.2, its reconstruction map, and Eqs.~(1.10)--(1.16).
Their Corollary~1.3 is the energy-dependent fixed-point formulation.
Its full eigenvalue condition tests all retained directions; the
scalar affine root (66) tests the fixed retained
direction only, as proved after Theorem~8.2.
Their operator $H-\lambda$ corresponds here to $\LL+z$, with
$\lambda=-z$, their complementary inverse to $(C_m+z)^{-1}$, and
their reconstruction to (36). Their effective
interaction has the negative sign in (34). The infinite-rank
form argument and the conditional coefficients are proved in this paper.

Bakry and Ledoux's reverse Poincare estimate is their Eq.~(1.7), with
$d(t)=2t$ at curvature bound zero, on pp.~686--687. Our single-edge
Markov generator is $\kappa\Delta^T$ and its carré du champ is
$\kappa\sum_a(X_af)^2$, giving precisely (8). Its full
compact-group proof was included. Simon's locally retained Theorem~1.1
uses the Euclidean generator $-\Delta/2$ and its Brownian-bridge formula.
It is a primary reference for positive potential weights, not a theorem
identifying Euclidean space with this compact configuration product;
the compact parabolic comparison needed here was proved in
(11). Mori's projection-memory method provides the
historical comparison: its projection is Eq.~(2.13) and its exact
memory equations are Eqs.~(3.1)--(3.10), pp.~430--431.
Our evolution uses $-\LL$ on functions with the vacuum inner product;
Mori's evolution uses a Liouville operator on dynamical variables.
Equation~(38) derives the corresponding Euclidean block
identity here, including its signs and domains.

The companion exact checker verifies the entire open-box skeleton and
path counts on several finite boxes, chain disjointness, nested product
paths, every retained loop orientation, exact Pauli rotations and the
mixed derivative tensor in explicit noncommuting samples, and exact
rational Schur associativity and spectral constants.
The zero-shift checker separately verifies physical norm and quotient
derivatives, all six interior/boundary and coupled/uncoupled affine
examples, the complete energy and norm decomposition, and the harmonic
gap comparison in 81 rational interacting block examples.
The physical-complement checker additionally verifies noncommuting
chain gauge actions, their inverse compositions, the exact original
strong-coupling constants, and the full inverse-moment continuity
identities and inequalities. These finite checks
diagnose algebra; they do not stand in for the all-box proofs above or
compute a substitute vacuum. No Lean run or large parallel job is used.

\section{Material-tensor transfer into interacting local-energy states}
\label{sec:fabel-tensor}
This section retains the original Wilson Hamiltonian, vacuum, gauge action,
spacing and coupling from Section~1. It constructs a particular observable
map from the full Fabel material tensor. No equality between fluid evolution
and quantum Hamiltonian evolution is asserted. The map and its exact
spectral measures are calculated below; the remaining infinite-volume
spectral question is not replaced by the value of a coordinate.

\subsection{Transport conventions and the vacuum-moment correction}
For the column-section connection $d+\mathcal A$, ordinary parallel transport
along an oriented edge $e:[0,a]\to\mathbb R^3$ solves
\[
 P_e'(s)=-\mathcal A(e(s))(\dot e(s))P_e(s),\qquad P_e(0)=I.
\]
Its gauge transform is $P_e^h(s)=h(e(s))^{-1}P_e(s)h(e(0))$.
Differentiating this formula verifies the transformed ODE and its initial
value; uniqueness gives the assertion. The links in Section~1 have the
opposite endpoint convention and are therefore \emph{inverse} transports:
\begin{equation}\label{eq:fabel-inverse-transport}
 U_e=P_e(a)^{-1},\quad
 U_e^h=h(e(0))^{-1}U_eh(e(a)),\quad U'(s)=U(s)\mathcal A(e(s))(\dot e(s)).
\end{equation}
Inverses and concatenations follow by multiplication of the corresponding
transport maps. Thus this is an exact gauge-equivariant map for arbitrary
noncommuting connection components. It is not a global inverse from finite
links to connections.

The Cartan specialization $\mathcal A_i=\lambda u_iT$, $T=-i\sigma_3/2$,
has
\[
 U_e=\exp\left(\lambda T\int_e u\cdot dx\right),\qquad
 U_p=\exp(\vartheta_pT),\quad
 \vartheta_p=\lambda\oint_{\partial p}u\cdot dx .
\]
With positively oriented $(i,j)$ faces, Stokes' theorem gives
$\vartheta_p=\lambda\int_p(\partial_i u_j-\partial_j u_i)\,dx^i dx^j$.
For a fixed smooth field on a compact region it equals
$a^2\lambda(\partial_i u_j-\partial_j u_i)+O(a^3)$.
In the general non-Abelian case the four edge expansions to second order
give $U_p=I+a^2F_{ij}+O(a^3)$: the differentiated terms are
$\partial_i\mathcal A_j-\partial_j\mathcal A_i$ and the four products
give $[\mathcal A_i,\mathcal A_j]$, with this positive sign.
For uniform $C^2$ bounds the remainder is uniform on that compact region.
As $\tr e^{\vartheta T}=2\cos(\vartheta/2)$, the exact magnetic value is
\begin{equation}\label{eq:fabel-magnetic-value}
 b(2-W_p)=\frac{2-2\cos(\vartheta_p/2)}{2g^2a}
       =\frac{\vartheta_p^2}{8g^2a}
          +O_g(\vartheta_p^4/a).
\end{equation}
For fixed $g,\lambda$ and a fixed smooth field, summation on a regular
mesh of a bounded region yields
$\lambda^2(8g^2)^{-1}\int|\operatorname{curl}u|^2dx$.
This follows by a Riemann sum; each leading cell term is $a^3$ times the
displayed density and each uniform cell error is $O(a^4)$.
Pointwise curvature growth alone does not imply growth of this integral:
the volume of the concentrating region must also be accounted for.

In Lorentzian coordinates $X^0=ct,X^i=x^i$, the same Cartan field with
$\mathcal A_0=0$ has the complete conserved current
\begin{align*}
 j_0&=0,\\
 j_i&=\frac{\lambda}{g^2}
       (\Delta u_i-c^{-2}\partial_t^2u_i)T\\
 &=\frac{\lambda}{g^2}\left[
 \Delta u_i-c^{-2}\left(\nu\Delta w_i
 -\sum_jw_j\partial_ju_i-\sum_ju_j\partial_jw_i
 -\partial_i\partial_tp+\partial_tf_i\right)\right]T,\qquad w_i=\partial_tu_i.
\end{align*}
This is obtained by differentiating the original Navier--Stokes equation,
not by changing its viscosity or force. $D^\nu j_\nu=0$ follows from
incompressibility and the common Cartan factor. The field solves that
sourced classical equation. It can be used as data for a quantum trial
state without claiming it solves source-free Yang--Mills. Proving
source-free classical dynamics is not a prerequisite for defining a trial
state; its physical Hilbert and spectral properties must instead be proved.

There is a separate important correction concerning Wilson observables.
Evaluating $W_C$ at the links of a background gives a number $W_C(U[\mathcal A])$.
The functional $W_C(U)$ on the full configuration space has no dependence
on the chosen background until an additional prescription supplies it.
The generic centered state $\psi(W_C-\mu W_C)$ is physical and nonzero,
but naming a background does not make it a background-dependent state map.
Moreover $r_C=\mu(W_C^2)$ is the actual quantum second moment. It is not
$W_C(U[\mathcal A])^2$.

Keep the exact bound from Section~7,
$\alpha\le r_C\le4-3\alpha$, with
$\alpha=(9\,12^{16/g^4})^{-1}>0$. Then
\begin{equation}\label{eq:fabel-loop-scale-corrected}
 \frac{3\alpha}{4}\kappa\ell\le
 a_C=\kappa\ell(1-r_C/4)
 \le(1-\alpha/4)\kappa\ell .
\end{equation}
For physical perimeter $P_C=a\ell$ this bounds the \emph{bare} loop
numerator between the displayed constants times $2g^2P_C/a^2$ at fixed $g$;
it does not give asymptotic ratio one. The reconstructed numerator and norm
remain $a_C-M_1(0)$ and $d_C+M_2(0)$.
The lower bound $\Delta_{\rm sc}$ for these reconstructed states is used
only for $\xi\le1/49152$. In a spatial continuum path $a\downarrow0$,
$\kappa\to0$ would imply $g^2=(a/2)\kappa\to0$ and thus $\xi\to\infty$.
Without $a\downarrow0$ the latter implication is not valid.

\subsection{The full weighted local-energy operator}
For arbitrary real weights $f=(f_e)_{e\in\mathsf E_L}$ define
\begin{equation}\label{eq:fabel-local-D}
 f_p=\frac14\sum_{e\in\partial p}f_e,\quad
 V_p=b(2-W_p),\quad
 D_f=\kappa\sum_ef_eE_e+\sum_pf_pV_p,\quad
 \Xi_f=(D_f-\langle\psi,D_f\psi\rangle)\psi .
\end{equation}
All sums include the boundary edges and faces.
The electric-only operator $\sum f_eE_e$ is a different operator.
On a Peter--Weyl product sector with edge spins $j_e$, the signed electric
sum in $D_f$ has eigenvalue $\kappa\sum_ef_ej_e(j_e+1)$.
Its domain consists of square-summable coefficients with the squares of
these eigenvalues as weights. Addition of the bounded real magnetic
multiplication preserves that self-adjoint domain. Smooth functions lie in
the domain, and $D_f$ maps smooth functions to smooth functions.
In particular $\Xi_f$ is smooth and belongs to every power domain of $H$.
Gauge invariance of each link Casimir and each face trace gives
gauge invariance of $\Xi_f$. Its definition proves
$\langle\psi,\Xi_f\rangle=0$ with no rescaling.

Put $\mathcal E=E_0$ and $\mathcal A=H-\mathcal E$.
The product rule, including both derivative terms, gives
\begin{align}
 [E_e,V_p]F&=-\sum_a\bigl[(X_{e,a}^2V_p)F+
                           2(X_{e,a}V_p)X_{e,a}F\bigr],\label{eq:fabel-local-comm}\\
 [H,D_f]&=\kappa\sum_p\sum_{e\in\partial p}(f_p-f_e)[E_e,V_p].
 \label{eq:fabel-current-full}
\end{align}
The first expression vanishes when $e\notin\partial p$.
The second identity follows by expanding the commutator: electric terms
commute with electric terms, magnetic multiplications commute, and the
two remaining electric--magnetic sums have the stated coefficient.
It proves
\[
 \mathcal A\Xi_f=[H,D_f]\psi,\qquad
 q_{\mathcal A}[\Xi_f]=\tfrac12
 \langle\psi,[D_f,[H,D_f]]\psi\rangle .
\]
For the raw positive measure
$\nu_f(B)=\|1_B(\mathcal A)\Xi_f\|^2$, $B\subset(0,\infty)$, spectral
integration yields
\begin{equation}\label{eq:fabel-three-moments}
 \|\Xi_f\|^2=\int d\nu_f,\quad
 q_{\mathcal A}[\Xi_f]=\int E\,d\nu_f,\quad
 \|[H,D_f]\psi\|^2=\int E^2d\nu_f .
\end{equation}
All moments are finite by smoothness. For a constant $f_e=c_0$,
$D_f=c_0H$, so its expectation is $c_0\mathcal E$ and $\Xi_f=0$.
This is an exact cancellation, not a zero-energy nonzero state.

\subsection{Coordinate-level cubic symmetries and the full covector map}
Write $m(e)=n+\mathbf e_i/2$ for edge midpoints. A signed permutation
matrix $R$ maps the vertex box to itself. If $R\mathbf e_i=\epsilon\mathbf e_j$,
then the positive image edge of $e=(n,i)$ is $(Rn,j)$ for $\epsilon=1$,
and $(Rn-\mathbf e_j,j)$ for $\epsilon=-1$.
Define $\Gamma_R(U)$ at this image edge to be $U_e^\epsilon$.
The inverse is $\Gamma_{R^{-1}}$; composing the signs and endpoints
verifies both inverse compositions. Every midpoint is mapped to $Rm(e)$.
Each face becomes a face, possibly with its orientation reversed.
Its SU(2) trace is unchanged under reversal.
Inversion of a group variable is Haar preserving and takes the
left Casimir to the equal right Casimir, since the metric is bi-invariant.
Thus the pullback unitary
$\mathsf T_RF=F\circ\Gamma_R^{-1}$ preserves $H$, $H^1$ and $H^2$ and
intertwines the full gauge actions by $h'(Rn)=h(n)$.
Uniqueness of the positive unit vacuum gives $\mathsf T_R\psi=\psi$.
For weights transported by $f_R(e')=f(e)$ it follows that
$\mathsf T_RD_f\mathsf T_R^{-1}=D_{f_R}$ and
$\mathsf T_R\Xi_f=\Xi_{f_R}$. Spectral projections of $\mathcal A$ commute
with $\mathsf T_R$.

Let $\Xi_r=\Xi_{m_r}$. Reflection in coordinate $r$ changes $\Xi_r$ to
$-\Xi_r$ and fixes $\Xi_s$ for $s\ne r$; permutations exchange the three.
Consequently, for every real $t\in\mathbb R^3$ and $c_0\in\mathbb R$,
\begin{equation}\label{eq:fabel-affine-raw}
 \Xi_{c_0+t\cdot m}=\sum_rt_r\Xi_r,\quad
 \langle D_{c_0+t\cdot m}\rangle=c_0\mathcal E,\quad
 \nu_{c_0+t\cdot m}(B)=|t|^2\nu_{m_1}(B).
\end{equation}
Indeed a reflection negates each off-diagonal matrix element
$\langle\Xi_r,1_B(\mathcal A)\Xi_s\rangle$ while leaving it invariant;
it is therefore zero. Permutations equate the diagonal elements.
This proof applies at every positive coupling, including when a state
or a measure vanishes.

For the retained Fabel time parameter $z=\sqrt{1-8\tau}>0$,
$\tau=(1-z^2)/8$, the full inverse material derivative is
computed from the unchanged real polynomial map
\begin{align*}
 F_1(x,y,w)&=(1+xy)^3w+y^2(1+xy)(4+3xy),\\
 F_2(x,y,w)&=y+3x(1+xy)^2w+3xy^2(4+3xy),\\
 F_3(x,y,w)&=2x-3x^2y-x^3w .
\end{align*}
Here $w$ denotes the original third spatial coordinate, not time.
Expanding the three-by-three determinant of its differentiated matrix
gives $\det DF=-2$. At the retained points the full matrices are
\[
 J_0=\begin{pmatrix}-9/16&-3/8&-1/8\\3/8&25/4&3/4\\-17/2&-3&-1\end{pmatrix},
 \quad
 J_\gamma=\begin{pmatrix}
 -9z^3/16&-3z/8&-1/8\\3z^2/8&25/4&3/(4z)\\-17/2&-3/z^2&-1/z^3
 \end{pmatrix}.
\]
Direct matrix multiplication gives
\begin{equation}\label{eq:fabel-C-full}
 C_\tau=
 \begin{pmatrix}
 (17-9z^3)/8&(3-3z^3)/(4z^2)&(1-z^3)/(4z^3)\\
 (51-24z^2-27z^3)/16&(9+8z^2-9z^3)/(8z^2)&(3-3z^3)/(8z^3)\\
 (-153+36z^2+117z^3)/8&(-27-12z^2+39z^3)/(4z^2)&(-9+13z^3)/(4z^3)
 \end{pmatrix}.
\end{equation}
It is $J_0^{-1}B_\tau^{-1}J_\gamma$, where $J=DF$ for the original
polynomials and
\[
q_0=(1,-3/2,13/2),\quad
\gamma=(z^{-1},-3z/2,13z^2/2),\quad B_\tau=I+6\tau E_{23}.
\]
The factorization gives the inverse $A_\tau=J_\gamma^{-1}B_\tau J_0$
and determinant one.
The source matrix entries are reproduced and independently rederived
by the companion checker from $F$.

The covector $\zeta_\tau=C_\tau^{\mathsf T}e_3$ has the three entries
in the last row of (102). Choose the explicit affine
weight $f_\tau(e)=\zeta_\tau\cdot m(e)$. Then
\[
 \nu_{f_\tau}(B)=Q(z)\nu_{m_1}(B),
\]
where the entire, unmodified scalar is
\begin{align*}
 Q(z)={}&\frac{(-153+36z^2+117z^3)^2}{64}
 +\frac{(-27-12z^2+39z^3)^2}{16z^4}
 +\frac{(-9+13z^3)^2}{16z^6},\\
 &\lim_{z\downarrow0}z^6Q(z)=81/16 .
\end{align*}
Thus the full covector, with its changing direction, is realized by this
specified parameter-to-state map. Its raw norm, energy, and second spectral
moment all have the same factor $Q(z)$; their ratio does not depend on
$\tau$ whenever the state is nonzero. This is not an identification of
the Fabel cotangent space with the whole quantum state space.

\subsection{A non-affine map using every entry of the material tensor}
Set $K_\tau=C_\tau C_\tau^{\mathsf T}$ and keep all six independent entries.
For any real symmetric matrix $S$ define
\begin{equation}\label{eq:fabel-quadratic-map}
 f_S(e)=m(e)^{\mathsf T}S\,m(e),\qquad
 \mathfrak M_{L,a,g}(S)=\Xi_{f_S}\in
 C^\infty(\mathcal Q)^{\mathcal G}\cap\{\psi\}^{\perp}.
\end{equation}
This is a linear, typed map on the six-dimensional real matrix space.
It does not change the Hamiltonian metric: $S$ is the spatial weight of
the observable $D_f$, not a replacement kinetic tensor in $H$.
Using physical midpoints $a\,m(e)$ multiplies $f_S$ and $\Xi_{f_S}$ by
$a^2$, and every raw spectral measure by $a^4$.
Those factors remain present if physical midpoint units are chosen.

Use the three unchanged seed functions
\[
 f_0=m_1^2+m_2^2+m_3^2,\quad
 f_{\rm d}=m_1^2-m_2^2,\quad f_{\rm o}=2m_1m_2,
\]
and write $\nu_0,\nu_{\rm d},\nu_{\rm o}$ for their actual centered-state
spectral measures. These are positive measures of the full interacting
Hamiltonian, not free-channel approximations.

\paragraph{Theorem 13.1.}
For every fixed $L\ge2$, $a,g>0$, every real symmetric $S$ and every Borel
$B\subset(0,\infty)$,
\begin{align}
 \nu_{f_S}(B)&=w_0(S)\nu_0(B)+w_{\rm d}(S)\nu_{\rm d}(B)
                         +w_{\rm o}(S)\nu_{\rm o}(B),\label{eq:fabel-sector-measure}\\
 w_0(S)&=(\tr S)^2/9,\nonumber\\
 w_{\rm d}(S)&=\tfrac12\sum_{i=1}^3(S_{ii}-\tr S/3)^2,\quad
 w_{\rm o}(S)=\sum_{i<j}S_{ij}^2.\nonumber
\end{align}
The same identity holds after integration against $1,E,E^2$ and
$e^{-sE}$ for $s\ge0$, retaining the corresponding raw norms,
excitation energies, current norms, and time correlations.

\paragraph{Proof.}
Put $u_i=\Xi_{m_i^2}$ and $v_{ij}=\Xi_{2m_im_j}$.
A coordinate reflection fixes each $u_i$ and negates $v_{ij}$ if it
reflects just one of $i,j$. Thus every spectral inner product between
a $u$ and a $v$ vanishes. For two different off-diagonal pairs there
is a reflection negating exactly one, so their cross term also vanishes.
Permutations give the same diagonal spectral norm, equal to $\nu_{\rm o}(B)$,
for the three $v_{ij}$.

For the three $u_i$, permutations equate diagonal elements, denoted
$d(B)$, and equate off-diagonal elements, denoted $h(B)$.
A permutation exchanging two indices shows that $h(B)$ equals its
complex conjugate; it is real. The diagonal contribution for $s_i=S_{ii}$
is
\[
 d\sum_i s_i^2+2h\sum_{i<j}s_is_j
 =(d-h)\sum_i(s_i-\bar s)^2+(d+2h)(\tr S)^2/3,
 \qquad\bar s=\tr S/3.
\]
But $\nu_0=3d+6h$ and $\nu_{\rm d}=2(d-h)$ by the definitions of the
unchanged seed states. Substitution gives (104).
All identities are identities of finite positive measures. Integration
against the listed nonnegative functions is justified by monotone
convergence; smoothness makes the moments finite.



Let $\eta=(1/4,3/8,-9/4)^{\mathsf T}$.
Every entry of (102) yields
$z^3C_\tau\to\eta e_3^{\mathsf T}$ and hence
\[
 z^6K_\tau\longrightarrow S_*=\eta\eta^{\mathsf T}
 =\frac1{64}\begin{pmatrix}4&6&-36\\6&9&-54\\-36&-54&324\end{pmatrix}.
\]
Consequently the full nonlinear spatial profile
$m^{\mathsf T}K_\tau m$ has the following exact endpoint:
\begin{equation}\label{eq:fabel-tensor-limit-measure}
 z^{12}\nu_{f_{K_\tau}}\longrightarrow
 \frac{113569}{36864}\nu_0+
 \frac{100825}{12288}\nu_{\rm d}+
 \frac{531}{512}\nu_{\rm o}.
\end{equation}
This is convergence in total variation at each fixed regulator and
convergence of the first two moments. To prove it, note that each
$w_j$ is homogeneous of degree two and continuous in $S$, and apply
(104). The total variation of the difference
is bounded by the sum of the three coefficient errors times the finite
masses of the seed measures; the moment versions use their finite moments.
The constants follow by substituting $S_{*,ii}=4/64,9/64,324/64$ and
$S_{*,12}=6/64,S_{*,13}=-36/64,S_{*,23}=-54/64$.
Equivalently $z^6\Xi_{f_{K_\tau}}\to\Xi_{f_{S_*}}$ in every Sobolev norm
at a fixed regulator, by the linear six-state expression.
These displayed rescaled limits record the growth of the raw state;
the state itself has not been divided by its norm.

All three coefficients in (105) are positive.
For sufficiently small $z>0$, each $w_j(K_\tau)$ is positive as well.
Then the positive-energy support of $\nu_{f_{K_\tau}}$ equals the union
of the supports of the nonzero seed measures: a Borel neighborhood
has zero mass in the sum precisely when it has zero mass in each seed.
Write $d_j=\int d\nu_j$, $n_j=\int E\,d\nu_j$.
If at least one $d_j$ is nonzero, the limiting quotient is exactly
\begin{equation}\label{eq:fabel-tensor-limit-quotient}
 \lim_{\tau\uparrow1/8}
 \frac{q_{\mathcal A}[\Xi_{f_{K_\tau}}]}{\|\Xi_{f_{K_\tau}}\|^2}
 =
 \frac{\frac{113569}{36864}n_0+
       \frac{100825}{12288}n_{\rm d}+\frac{531}{512}n_{\rm o}}
      {\frac{113569}{36864}d_0+
       \frac{100825}{12288}d_{\rm d}+\frac{531}{512}d_{\rm o}} .
\end{equation}
This is an evaluated relation between actual moments, not an estimate
that any unknown moment is small. If all $d_j=0$, every state in this
quadratic family is zero. The next calculation proves that this latter
case does not occur for sufficiently small positive $\xi$ on any fixed box.

\subsection{The first surviving state and exact open-boundary coefficients}
The full local-energy source proof is retained unchanged as
\path{references/LOCAL_ENERGY_CURRENT_INPUT.md}. We reproduce the
coefficient argument needed here. Set $H_0=\sum_eE_e$ and
$\mathcal W=\sum_pW_p$. On a fixed box, $\|\mathcal W\|_\infty=2M$ and
the nonconstant free gap is $3/4$.
The circle $|\zeta|=3/8$ has free resolvent norm at most $8/3$, so the
Neumann series for $H_0-\xi\mathcal W$ converges for $|\xi|<3/(16M)$.
The simple ground projection persists there. For real $\xi$ near zero
its positive unit vacuum is real analytic in each Sobolev space, by the
eigen-equation and elliptic regularity. Keep its unit-vacuum scalar:
\begin{align*}
 \psi_\xi&=1+\xi\mathcal W/3+\xi^2 B+O_{H^k,L}(\xi^3),\\
 H_0B&=(\mathcal W^2-M)/3,\qquad \int B=-M/18.
\end{align*}
For two faces sharing one edge, decompose
$W_pW_q=P_{pq,0}+P_{pq,1}$ by Haar averaging that edge.
The first term is one half the six-edge boundary trace, by the
fundamental matrix-element identity $\int U_{ij}\overline U_{kl}
=\delta_{ik}\delta_{jl}/2$.
The shared-edge spins are zero and one, and all six outer spins
are $1/2$. Hence their free energies are $9/2,13/2$ and their squared
norms are $1/4,3/4$.
Repeated faces give $W_p^2-1$ of energy $8$; edge-disjoint distinct faces
give $W_pW_q$ of energy $6$. All unordered pairs are thus retained in
\[
 B=-M/18+\sum_p\frac{W_p^2-1}{24}
  +\sum_{\partial p\cap\partial q=\varnothing}\frac{W_pW_q}{9}
  +\sum_{|\partial p\cap\partial q|=1}
       \left(\frac4{27}P_{pq,0}+\frac4{39}P_{pq,1}\right).
\]
Inserting this expression in (97), the first-order
electric term equals the first-order magnetic term because
$E_f(\mathcal W/3)=\sum_pf_pW_p$, where $E_f=\sum_ef_eE_e$.
Centering subtracts the expectation of the same vacuum.
Repeated and edge-disjoint face pairs cancel at order two.
For an adjacent pair with shared edge $e$, the weighted electric
eigenvalues are $3(f_p+f_q)-3f_e/2$ and $3(f_p+f_q)+f_e/2$.
The remaining second-order state is therefore
\begin{align}
 \Xi_f&=\kappa\xi^2 Z_f+O_{H^k,L,f}(\kappa\xi^3),\nonumber\\
 Z_f&=\sum_{\{p,q\}\ {\rm adjacent}}\delta_{pq}(f)
         (P_{pq,0}/9-P_{pq,1}/39),\quad
 \delta_{pq}(f)=f_p+f_q-2f_e. \label{eq:fabel-adjacent-state}
\end{align}
Different adjacent pairs have different six-edge outer boundaries.
Otherwise their symmetric difference would be a nonempty mod-two
two-cycle with at most four elementary faces. A face in that cycle
needs four distinct other faces to cancel its four edges (two distinct
square faces share at most one edge), a contradiction.
Central sign reversal of a boundary link present in only one pair
therefore kills the cross integral. It also kills the cross integral
with $H_0$ inserted. Both spin channels remain orthogonal. Thus
\begin{align}
 \|Z_f\|^2&=\frac{49}{13689}\mathcal D_f,&
 \langle Z_f,H_0Z_f\rangle&=\frac2{117}\mathcal D_f,&
 \mathcal D_f&=\sum_{\{p,q\}\ {\rm adjacent}}\delta_{pq}(f)^2 .
 \label{eq:fabel-D-moments}
\end{align}
The constants are respectively $1/(4\cdot9^2)+3/(4\cdot39^2)$ and
$(9/2)/(4\cdot9^2)+3(13/2)/(4\cdot39^2)$.
The central transformation
$U_i(n)\mapsto(-1)^{\sum_{j<i}n_j}U_i(n)$ negates every $W_p$ and
preserves the link Casimirs. It identifies the centered problems at
$\xi$ and $-\xi$, including the centered weighted state.
Consequently norm and energy are even:
\begin{equation}\label{eq:fabel-first-moments}
 \|\Xi_f\|^2=\frac{49}{13689}\kappa^2\xi^4\mathcal D_f+
 O_{L,f}(\kappa^2\xi^6),\quad
 q_{\mathcal A}[\Xi_f]=\frac2{117}\kappa^3\xi^4\mathcal D_f+
 O_{L,f}(\kappa^3\xi^6).
\end{equation}

Here are complete all-box counts for the new quadratic profiles.
For $e=(n,i)$, put $m=n+\mathbf e_i/2$.
A face in transverse direction $j$ on side $\epsilon$ has center
$m+\epsilon\mathbf e_j/2$. Averaging its four edge-midpoint quadratic
values gives
\[
 f_p=f_S(e)+\epsilon(Sm)_j+(S_{ii}+3S_{jj})/8.
\]
Indeed its midpoint displacements from its center are
$\pm\mathbf e_i/2,\pm\mathbf e_j/2$; the average correction to the
center value is $(S_{ii}+S_{jj})/8$.
Consequently opposite coplanar faces give
$\delta=(S_{ii}+3S_{jj})/4$.
Two hinged faces in directions $j,k$ on sides $\epsilon,\eta$ give
\[
 \delta=\epsilon(Sm)_j+\eta(Sm)_k+c_i,\qquad
 c_i=(2S_{ii}+3S_{jj}+3S_{kk})/8 .
\]
For each longitudinal coordinate the count and second midpoint moment
are $N_i=2L$ and $I_2=L(4L^2-1)/6$.
For each oriented transverse pair $(n,\epsilon)$ satisfying
$-L\le n,n+\epsilon\le L$, the sums are
\[
 N=4L,\quad \sum n=\sum\epsilon=0,\quad
 N_2=\sum n^2=\frac{2L(2L^2+1)}3,\quad
 E_n=\sum\epsilon n=-2L .
\]
These follow from the sums over $n=-L,\ldots,L-1$ at $\epsilon=1$
and $n=-L+1,\ldots,L$ at $\epsilon=-1$.
Expansion of every square now yields the exact finite formula
\begin{align*}
 \mathcal D_S^{\rm cop}&=N_i(2L-1)(2L+1)
                  \sum_{i\ne j}(S_{ii}+3S_{jj})^2/16,\\
 \mathcal D_S^{\rm hinge}
 &=\sum_i\bigl[
 I_2N^2(S_{ij}^2+S_{ik}^2)
 +N_iNN_2(S_{jj}^2+S_{kk}^2+2S_{jk}^2)\\
 &\hspace{12mm}+2N_iE_n^2(S_{jj}S_{kk}+S_{jk}^2)
 +2c_iN_iNE_n(S_{jj}+S_{kk})+c_i^2N_iN^2\bigr].
\end{align*}
In each summand $\{j,k\}=\{1,2,3\}\setminus\{i\}$ and
$\mathcal D_S=\mathcal D_S^{\rm cop}+\mathcal D_S^{\rm hinge}$.
There is no periodic boundary replacement.

Substitution of the three seed matrices gives
\begin{align}
 \mathcal D_0&=4L(2L-1)(2L+1)(4L^2+3),\nonumber\\
 \mathcal D_{\rm d}&=L(256L^4+74L^2-21)/6,\nonumber\\
 \mathcal D_{\rm o}&=64L^3(2L^2+1)/3.\label{eq:fabel-three-D}
\end{align}
All are positive for $L\ge2$.
The same substitution for $S_*=\eta\eta^{\mathsf T}$, or the already
proved cubic decomposition of the coefficient Gram matrix, gives
\begin{equation}\label{eq:fabel-rank-one-D}
 \mathcal D_{S_*}
 =\frac{L(14536832L^4+5453566L^2-1614327)}{24576}>0 .
\end{equation}
Equation (109) proves nonzero seed states for
sufficiently small positive $\xi$ on each fixed box and proves
\[
 \frac{q_{\mathcal A}[\Xi_{f_{S_*}}]}{\|\Xi_{f_{S_*}}\|^2}
       =\frac{234}{49}\kappa+O_L(\kappa\xi^2).
\]
This is an ordered fixed-box coefficient result, not a volume-uniform
remainder bound. In original variables
$\kappa^2\xi^4=1/(64g^{12}a^2)$ and
$\kappa^3\xi^4=1/(32g^{10}a^3)$.
The $L^5$ growth in (111) belongs to the computed
Taylor coefficient of both raw moments. Their common spatial and
material growth does not by itself reduce the quotient.

\subsection{Full Schur map and the remaining collective spectral question}
For a physical state $\Xi_f$ in the original Haar space set
$h_f=\Xi_f/\psi$, which is smooth at each finite regulator, and retain
\[
 F_f=J_m^*h_f,\quad
 v_f=\mathcal R_mF_f,\quad
 k_f=Q_mh_f+C^{-1}B_mF_f.
\]
These are the exact maps and domains of Theorem~9.1.
The identity $h_f=v_f+k_f$ has
\begin{align*}
 q(h_f,h_f)&=s_{m,0}(F_f,F_f)+c_m(k_f,k_f),\\
 \|\Xi_f\|^2&=\|v_f\|^2+\|k_f\|^2
                         +2\operatorname{Re}\langle v_f,k_f\rangle,\\
 \|v_f\|^2&=\|F_f\|^2+\int t^{-2}\,d\sigma_{F_f}(t),\\
 s_{m,0}(F_f,F_f)&=a_m(F_f,F_f)-\int t^{-1}\,d\sigma_{F_f}(t).
\end{align*}
No interaction or norm term has been removed.
The signed permutations above preserve the skeleton with $m\mid2L$:
the transverse coarse coordinates $-L+m\{0,\ldots,2L/m\}$ are invariant
under sign reversal, and chain products reverse by group inversion.
Vacuum conditional averaging therefore intertwines these symmetries by
change of variables in its defining integral. So do $B_m,C$ and the
full memory kernels; for the latter use functional calculus for
$e^{-sC}$ or $(C+z)^{-1}$. Their spectral matrix elements have the same
symmetry-forced vanishing cross terms. This does not evaluate the
remaining same-sector measures by a free approximation.

In the proved regime $\xi\le1/49152$, every nonzero constructed physical
state obeys $q/\|\Xi\|^2\ge\kappa(3-512\xi)\ge287\kappa/96$.
Outside that regime the exact three-measure formula remains valid but
does not compute the unknown seed measures. On any regulator sequence,
at least one of the nonzero three seed states has quotient no larger
than the quadratic combination, since (104)
makes its quotient a weighted mean with weights $w_jd_j\ge0$.
Thus a vanishing quotient in this full-tensor quadratic family would
also give a vanishing quotient in a selected seed sector along that
sequence. At fixed regulator (106)
instead gives the precise endpoint measure and ratio. No exchange of
the material endpoint, infinite-volume limit, or spatial continuum
limit has been made. The calculation leaves those three interacting
seed-sector limits as substantive research, not as assumed estimates.
Section~14 now calculates their complete fixed-box
weak-coupling pair measures, proves an injective lowest-band tensor map,
and constructs an actual positive-coupling projected Fabel sequence with
a vanishing energy quotient. It retains the whole-state raw norm and
band fraction. The common interacting continuum identification remains
unresolved.

\section{The exact lowest-energy image of the full Fabel tensor}
\label{sec:fabel-low}

This section continues the actual tensor map of Section~13 through a
spectral projection of the original interacting Hamiltonian. The resulting
states are explicitly derived from the full material tensor. Their
positive-coupling finite-regulator energy quotients tend to zero.
The construction does not identify their changing Hilbert spaces with an
interacting four-dimensional continuum theory.

\subsection{Full weighted weak-coupling state map}

The source proofs retained with this manuscript are \emph{The original
finite-box SU(2) Hamiltonian at small positive coupling},
\emph{Actual weak-coupling vacuum, observable and spectral maps}, and
\emph{Full-vacuum potential moments and the electric-state graph limit}.
We retain their full proofs as authored reference files. Here is the
complete dictionary and the additional weighted magnetic calculation.

The original vertex gauge law is $U_e\mapsto h_s^{-1}U_eh_t$. Setting
$k_v=h_v^{-1}$ gives exactly the same transformation written
$U_e\mapsto k_sU_ek_t^{-1}$ in those sources. Thus their invariant function
space is identical, with the identity map on link configurations and
wavefunctions. Root the specified tree at $(-L,-L,-L)$, using the first
coordinate, in order $1,2,3$, that exceeds $-L$ to select each vertex's
parent. If $t_v$ is its ordered root-to-vertex holonomy, put
$Z_c=t_{s(c)}U_ct_{t(c)}^{-1}$ for each non-tree chord $c$.
The inverse keeps the tree variables and sets $U_c=t_s^{-1}Z_ct_t$.
For each fixed tree configuration these chord changes are independent
left/right Haar translations. Fubini gives the exact product measure
identity. Based gauge invariance removes the tree variables; the
remaining gauge action is simultaneous conjugation of every $Z_c$.
This constructs the same physical Hilbert space, including the residual
constraint.

On real cochains with counting inner products let $d_0$ be the original
vertex difference and $d_1$ the oriented four-edge curl. Let $p_v(x)$
be the additive tree integral and
\[
 (Tx)_c=p_{s(c)}(x)+x_c-p_{t(c)}(x),\quad
 G=TT^{\mathsf T},\quad C=d_1\iota,\quad
 R=T^{\mathsf T}G^{-1/2},
\]
where $\iota$ inserts chords with zero tree entries.
Square interchanges of monotone paths prove
$\ker T=\ker d_1=\operatorname{im}d_0$ and $T\iota=I$.
Consequently $R^{\mathsf T}R=I$, its range is
$\ker d_0^{\mathsf T}$, and $d_1R=CG^{1/2}$.
Choose an orthogonal matrix $O$ with
\[
 O^{\mathsf T}G^{1/2}C^{\mathsf T}CG^{1/2}O=\Sigma^2,\qquad
 \Sigma=\operatorname{diag}(\sigma_\nu),\quad \sigma_\nu>0.
\]
Thus the columns of $RO$ are an orthonormal transverse edge basis
with curl frequencies $\sigma_\nu$. None of $T,G,C,R,O$ changes
the original edge or face inner products.

Write $Z_c=\exp(y_c^\alpha T_\alpha)$, $|y_c|<2\pi$.
Outside a Haar-null set these logarithms are unique. For
$\Omega_g=\{x:|x_c|<2\pi/g\}$ the exact Hilbert map is
\begin{equation}\label{eq:low-chart}
 (\mathcal B_gF)(x)=
 g^{3r/2}\mathcal J(gx)^{1/2}F(\exp(gx)),\qquad
 \mathcal J(y)=(16\pi^2)^{-r}
 \prod_c\left(\frac{\sin(|y_c|/2)}{|y_c|/2}\right)^2,
\end{equation}
extended by zero off $\Omega_g$, where $r=|\mathsf E_L|-|\mathsf V_L|+1$.
Change of variables proves that it is an isometry onto the supported
subspace; its adjoint includes the reciprocal density and dilation.
Both commute with simultaneous colour rotations.

Fix $L,a,f$ before sending $g$ to zero. Put $W_p= \operatorname{tr}U_p$
as before, $w_p=2-W_p$, $W=\sum_pw_p$, and $W_f=\sum_pf_pw_p$.
The vacuum and spectral source proofs give
\[
 \mathcal B_g\psi_g\longrightarrow\Phi_0,\qquad
 \mathcal B_gP_g(I)\mathcal B_g^*\longrightarrow P_0(I)
 \quad\hbox{in operator norm}
\]
for bounded excitation intervals $I$ whose endpoints avoid the comparison
spectrum. Here $P_g(I)=\mathbf1_I(H_g-\mathcal E_g)$ on the physical
space and
\[
 H_{\mathrm{osc}}=\frac1a\sum_{\nu,\alpha}
 \left(-2\partial_{z_{\nu,\alpha}}^2+
             \frac{\sigma_\nu^2z_{\nu,\alpha}^2}{8}\right),
 \qquad z^\alpha=O^{\mathsf T}G^{-1/2}x^\alpha .
\]
The density and constant Jacobian in this linear coordinate map remain
part of the unitary representation. In $z$ coordinates the unit
Gaussian has variance $2/\sigma_\nu$ in each real coordinate and
ground energy $\mu_0=3\sum_\nu\sigma_\nu/(2a)$.

For specificity, these convergence statements follow from the full
operator, as follows. The full potential has a unique zero at all
$Z_c=I$, by the contained-square path argument, and its Taylor term is
$\frac14\sum_\alpha\|Cy^\alpha\|^2$. It bounds $c|y|^2$ from below in a
fixed small chart and has positive minimum $c_\rho$ outside that chart.
A vector with bounded original energy has chart-tail mass bounded by
$2aK/(cR_0^2)+2aKg^2/c_\rho$ outside $|x|\le R_0$.
Local ellipticity gives a local $H^1$ bound, so Rellich compactness
and that tail estimate give global $L^2$ compactness.
Cut-off invariant oscillator eigenfunctions give the min--max upper
bound. The exact partition-of-unity identity for the full kinetic
form has error $O_{L,\rho}(g^2)$; the exterior potential and interior
quadratic comparison give the lower bound on each finite eigenvalue.
The bounded-energy compactness identifies every subsequential
eigenvector limit. Equality of finite ranks then proves the stated
finite spectral projection convergence. Positivity fixes the vacuum
phase. The retained source files contain the complete form estimates,
domains, cutoffs and min--max calculation.

The electric graph result for arbitrary real $f_e$ is
\[
 \mathcal B_g\left(g^2\sum_ef_eE_e\psi_g\right)
 \longrightarrow
 -\sum_{\alpha,c,d}(T\operatorname{diag}(f_e)T^{\mathsf T})_{cd}
              \partial_{x_c^\alpha}\partial_{x_d^\alpha}\Phi_0 .
\]
Its derivative control comes from the vacuum equation and all potential
moments, not vacuum $L^2$ convergence alone. Indeed
$\sum_{e,\alpha}|X_{e,\alpha}W|^2\le16W$ follows from
$\sum_\alpha|X_{e,\alpha}w_p|^2=w_p-w_p^2/4$ and at most four incident
faces. The ground-state transform applied to $W^{k/2}$, with positive
regularization, gives
\[
 b\langle W^{k+1}\rangle
 \le\mathcal E_g\langle W^k\rangle+4\kappa k^2\langle W^{k-1}\rangle.
\]
At fixed $L,a$, $\mathcal E_g$ is bounded. Hence every
$\langle(W/g^2)^k\rangle$ is bounded. This proves strong convergence
of $(W/g^2)\psi_g$ by local Taylor expansion and its third-moment
tail bound. The vacuum equation gives the strong graph limit for
$g^2\sum_eE_e\psi_g$.
The strongly commuting nonnegative link Casimirs obey
$\|\sum_ef_eE_eF\|\le\|f\|_\infty\|\sum_eE_eF\|$.
Apply this inequality to $\psi_g-\mathcal B_g^*(\chi_{R_0}\Phi_0)$,
use convergence of all differentiated coefficients on the fixed
compact support, and then let $R_0$ grow. Gaussian graph decay proves
the displayed arbitrary-weight limit.

The missing magnetic-weight extension is also strong:
\begin{equation}\label{eq:low-mag}
 \mathcal B_g((W_f/g^2)\psi_g)
 \longrightarrow \frac14\sum_\alpha
 (Cx^\alpha)^{\mathsf T}\operatorname{diag}(f_p)Cx^\alpha\Phi_0.
\end{equation}
To prove it, use the exact inequality
$|W_f|\le\max_p|f_p|\,W$. The same third-moment tail estimate controls
its squared norm outside every chart ball. On a fixed ball every face
word has its full Taylor limit; finite summation, strong vacuum
convergence and bounded local multipliers give the limit there.
The limiting quadratic Gaussian has an integrable squared tail.
These statements prove (113) globally, for signed weights
as well as positive weights.

Define the full matrices
\begin{align}
 A_f&=O^{\mathsf T}R^{\mathsf T}\operatorname{diag}(f_e)RO,
 \qquad
 M_f=O^{\mathsf T}R^{\mathsf T}d_1^{\mathsf T}
                  \operatorname{diag}(f_p)d_1RO, \nonumber\\
 \mathsf R_f&=\Sigma^{-1/2}M_f\Sigma^{-1/2}
                   -\Sigma^{1/2}A_f\Sigma^{1/2}. \label{eq:low-pair}
\end{align}
All magnetic edge averages are the original $f_p=\frac14\sum_{\partial p}f_e$.
Combining both graph limits with $\kappa=2g^2/a$ and $b=1/(2g^2a)$
proves the actual raw centered-vector limit
\begin{equation}\label{eq:low-vector}
 \mathcal B_g\Xi_f\longrightarrow
 V_f=\frac1{4a}\sum_{\nu,\eta}(\mathsf R_f)_{\nu\eta}
       \sum_{\alpha=1}^3
       a_{\nu,\alpha}^\dagger a_{\eta,\alpha}^\dagger\Phi_0 .
\end{equation}
Here $z_{\nu,\alpha}=\sqrt{2/\sigma_\nu}
(a_{\nu,\alpha}+a_{\nu,\alpha}^\dagger)$, with the usual commutator
equal to $\delta_{\nu\eta}\delta_{\alpha\beta}$.
To check the coefficient, differentiating the Gaussian makes the
electric centered creation coefficient $-\Sigma^{1/2}A_f\Sigma^{1/2}/(4a)$.
The magnetic quadratic is $\sum_\alpha z_\alpha^{\mathsf T}M_fz_\alpha/(8a)$;
its centered creation coefficient is
$\Sigma^{-1/2}M_f\Sigma^{-1/2}/(4a)$.
The scalar expectations converge by inner products with the converging
unit vacuum. This retains both terms and their cancellation; for constant
$f$ the matrices give $\mathsf R_f=0$, agreeing with $D_f=fH$.

A diagonal colour contraction has squared norm $6$, and an
off-diagonal contraction has squared norm $3$. Commuting its two
annihilators through its two creators proves these values, with
different unordered pairs orthogonal. Thus the full limiting raw measure is
\begin{equation}\label{eq:low-full-measure}
 \nu_{f,0}=\frac3{8a^2}\sum_\nu(\mathsf R_f)_{\nu\nu}^2
                       \delta_{2\sigma_\nu/a}
 +\frac3{4a^2}\sum_{\nu<\eta}(\mathsf R_f)_{\nu\eta}^2
                       \delta_{(\sigma_\nu+\sigma_\eta)/a}.
\end{equation}
Its total mass is $3\operatorname{tr}(\mathsf R_f^2)/(8a^2)$ and its
first moment is $3\operatorname{tr}(\Sigma\mathsf R_f^2)/(4a^3)$.
Strong vector convergence and finite spectral projection convergence,
followed by a finite comparison-energy cutoff, prove weak convergence
of the actual raw measures with total-mass convergence. We do not
infer convergence of an unbounded energy test from weak convergence.
For the bounded spectral intervals used below, the projected
energy moments do converge.

\subsection{All open-boundary coefficients in the lowest mode block}

Put $N=2L+1$, $h=\pi/(2N)$, $s=2\sin h$, and $\sigma=\sqrt2s$.
The one-dimensional vertex and edge functions, in the original
coordinates, are
\begin{align*}
 v_0(n)&=N^{-1/2},\quad
 v_1(n)=-\sqrt{2/N}\sin(\pi n/N), &&-L\le n\le L,\\
 w_1(n)&=-\sqrt{2/N}\cos(\pi(n+\tfrac12)/N), &&-L\le n<L.
\end{align*}
Finite geometric sums prove their unit norms, orthogonality of
$v_0,v_1$, and $dv_1=sw_1$, $d^*w_1=sv_1$ with the open endpoints.
For each cyclic triple $(i,p,q)=(1,2,3),(2,3,1),(3,1,2)$ define the
edge cochain $V_i$ by its only two nonzero directional components:
\begin{equation}\label{eq:low-modes}
 (V_i)_p(n)=\frac{v_0(n_i)w_1(n_p)v_1(n_q)}{\sqrt2},\qquad
 (V_i)_q(n)=-\frac{v_0(n_i)v_1(n_p)w_1(n_q)}{\sqrt2}.
\end{equation}
The exact vertex divergence cancels between the two components.
Tensor products give $\langle V_i,V_j\rangle=\delta_{ij}$.
Their curls have disjoint face-plane supports and squared norms $\sigma^2$.
For example $d_1V_1$ on plane $23$ is
$-\sigma v_0(n_1)w_1(n_2)w_1(n_3)$; the other two signs follow from
the same original positive face ordering. These are precisely the three
lowest transverse modes: the full frequency list is
$\sigma_{\mathbf j}^2=\sum_i4\sin^2(\pi j_i/(2N))$,
$0\le j_i\le N-1$, with multiplicity $k-1$ for $k\ge2$ positive indices.
It follows by tensoring the one-dimensional difference bases; curl on a
fixed index triple has squared form $|s|^2|v|^2-|s\cdot v|^2$.
This also proves completeness and all endpoint counts.

Retain $f_S(e)=m(e)^{\mathsf T}S m(e)$ for real symmetric $S$,
$m(e)=n+\mathbf e_i/2$. Define
\begin{equation}\label{eq:low-delta}
 \delta_L=\frac{1-\cot^2h}{8},\qquad
 \mathfrak m_L^2=\frac{\cos^2h}{2N^2\sin^4h}.
\end{equation}

\paragraph{Theorem 14.1.}
On the three-mode basis (117), the exact block
$\mathsf R_S^{\min}$ of (114) is
\begin{equation}\label{eq:low-block}
 (\mathsf R_S^{\min})_{ii}=\sigma\delta_L(\operatorname{tr}S-S_{ii}),
 \qquad
 (\mathsf R_S^{\min})_{ij}=\sigma\mathfrak m_L^2S_{ij}\quad(i\ne j).
\end{equation}
The linear map from real symmetric $S$ to this real symmetric block is
injective. Its full raw first-physical-band mass is
\begin{equation}\label{eq:low-d}
 d_{L,a}(S)=\frac{3\sigma^2}{8a^2}
 \left[
 \delta_L^2\sum_i(\operatorname{tr}S-S_{ii})^2+
 2\mathfrak m_L^4\sum_{i<j}S_{ij}^2
 \right].
\end{equation}
It is strictly positive for every $S\ne0$.


\paragraph{Proof.}
Let
\[
 A=\sum_{n=-L}^Ln^2v_1(n)^2,\quad
 B=\sum_{n=-L}^{L-1}(n+\tfrac12)^2w_1(n)^2,\quad
 M=\sum_{n=-L}^Ln\,v_1(n)v_0(n).
\]
These exact finite sums obey
\begin{equation}\label{eq:low-moments}
 B-A=-\tfrac14\cot^2h,\qquad
 M=-\frac{\sqrt2\cos h}{2N\sin^2h},\qquad M^2=\mathfrak m_L^2.
\end{equation}
Here is a derivation including the endpoint terms. Put
\[
 D_N(t)=\sum_{n=-L}^Le^{int}=\frac{\sin(Nt/2)}{\sin(t/2)},
 \qquad D_{N-1}(t)=\sum_{n=-L}^{L-1}e^{i(n+1/2)t}.
\]
The second sum equals
$D_N(t)\cos(t/2)-\cos(Nt/2)$.
Using $\sin^2u=(1-\cos2u)/2$ and $\cos^2u=(1+\cos2u)/2$ gives
\[
 B-A=\frac1N\left[
 -\frac{N(N-1)}4-D_N''(2\pi/N)-D_{N-1}''(2\pi/N)\right].
\]
Direct differentiation of the displayed quotient and identity yields
\[
 D_N''(2\pi/N)+D_{N-1}''(2\pi/N)
 =\frac{N}{4\sin^2h}-\frac{N^2}4.
\]
Substitution gives the first equality in (121).
The first derivative at $\pi/N$ is
\[
 D_N'(\pi/N)=-\frac{\cos h}{2\sin^2h}.
\]
This gives the displayed $M$.
Thus neither endpoint lattice has been replaced by an integral.

For a face in plane $pq$ its four edge midpoints have mean
$c=n+(\mathbf e_p+\mathbf e_q)/2$, and averaging the four full
quadratics gives
\[
 f_p=c^{\mathsf T}S c+(S_{pp}+S_{qq})/8 .
\]
For the mode normal to $i$, parity removes the mixed-coordinate terms
in both diagonal expectations. The electric expectation is
$(S_{pp}+S_{qq})(A+B)/2+S_{ii}C_0$ and the unit-curl face expectation
is $(S_{pp}+S_{qq})(B+1/8)+S_{ii}C_0$, where
$C_0=\sum n^2v_0(n)^2$ is retained in both and cancels exactly.
Their difference is
$(S_{pp}+S_{qq})((B-A)/2+1/8)$, which is
$\delta_L(\operatorname{tr}S-S_{ii})$.
The common frequency factors in (114) give the diagonal
part of (119).

For distinct normal indices $i,j$, the two cochains share one edge
direction. The product of their two signed components is minus one half
of the squared $w_1$ factor times one $v_0v_1$ factor in each of
the other directions. Only the term $2S_{ij}n_in_j$ in the full
edge quadratic survives the two odd sums. The electric matrix entry
is therefore $-S_{ij}M^2$. Their curls have disjoint face-plane
supports, so the magnetic matrix entry is zero even with arbitrary
face weights. The negative electric sign in (114)
gives the positive off-diagonal entry in (119).

Since $L\ge2$, $0<h\le\pi/10$; hence $\delta_L<0$ and
$\mathfrak m_L^2>0$. A zero block first forces every off-diagonal
$S_{ij}$ to vanish and then $S_{ii}=\operatorname{tr}S$ for all $i$.
Summing gives $\operatorname{tr}S=3\operatorname{tr}S$, so all entries
vanish. For an explicit inverse, if $B=\mathsf R_S^{\min}/\sigma$,
then
$\operatorname{tr}S=\operatorname{tr}B/(2\delta_L)$,
$S_{ii}=\operatorname{tr}B/(2\delta_L)-B_{ii}/\delta_L$, and
$S_{ij}=B_{ij}/\mathfrak m_L^2$.
Finally the six colour contractions at energy $2\sigma/a$ and their
squared norms in (116) give (120).



The first physical energy is $\Delta=2\sigma/a$ with multiplicity six.
Every one-quantum state transforms as a colour vector and has no
invariant component. Among two-quanta states the only energy $\Delta$
comes from the three lowest modes, and rotation invariance of a
bilinear tensor forces the scalar colour contraction.
The next transverse frequency is $\sqrt3s$: triples $(1,1,1)$
give it, whereas a triple including $2$ has squared frequency at least
$(1+4\cos^2h)s^2>3s^2$.
All other two-quanta energies are at least
$(\sqrt2+\sqrt3)s/a>(11/10)\Delta$.
Every state with three or more quanta has energy at least
$3\sigma/a>(11/10)\Delta$.
Consequently
\[
 I_{L,a}=(\Delta/2,11\Delta/10)
\]
isolates precisely the first six-dimensional physical comparison space.
For $S\ne0$, (115) and spectral projection convergence
prove
\begin{equation}\label{eq:low-positive-g}
 \|P_g(I_{L,a})\Xi_{f_S}\|^2\longrightarrow d_{L,a}(S)>0.
\end{equation}
This is a conclusion for the actual positive-coupling vector, not
an assignment of the comparison Gaussian as its vacuum.

\subsection{The entire raw state and its lowest-band fraction}

Set
\[
 \mathcal N_{L,a}(S)=\frac3{8a^2}\operatorname{tr}(\mathsf R_{f_S}^2).
\]
This is the limiting squared norm of the entire unchanged tensor state,
including all higher mode pairs in (116).
It is positive for $S\ne0$, since its lowest block is nonzero.
The exact limiting fraction in the first physical band is therefore
\begin{equation}\label{eq:low-fraction}
 \rho_{L,a}(S)=
 \frac{d_{L,a}(S)}{\mathcal N_{L,a}(S)}
 =\frac{\operatorname{tr}((\mathsf R_S^{\min})^2)}
        {\operatorname{tr}(\mathsf R_{f_S}^2)}.
\end{equation}
This ratio records two separately retained raw masses; it does not
change the vector. Both strong vector and projection convergence prove
that the corresponding fraction for $\Xi_{f_S}$ at positive $g$
tends to (123).

There is an explicit lower bound retaining all original modes:
\begin{equation}\label{eq:low-fraction-bound}
 1\ge\rho_{L,a}(S)\ge
 \frac{\sigma^2\delta_L^2}{432rL^4}>0,\qquad
 r=4L^2(4L+3).
\end{equation}
To prove it, let $V=RO$ and $Y=d_1V\Sigma^{-1}$. Both are
isometries into the original edge and face spaces, respectively.
Equation (114) becomes
\[
 \mathsf R_{f_S}=\Sigma^{1/2}
  [Y^{\mathsf T}\operatorname{diag}(f_p)Y
             -V^{\mathsf T}\operatorname{diag}(f_e)V]\Sigma^{1/2}.
\]
Every original midpoint has squared length at most
$B_L=2L^2+(L-\tfrac12)^2<3L^2$. Thus
$|f_e|,|f_p|\le B_L\|S\|_{\mathrm{op}}$.
Since $\sigma_{\max}^2\le12$, compression and finite-rank
Hilbert--Schmidt estimates give
\[
 \operatorname{tr}(\mathsf R_{f_S}^2)
 \le 4\sigma_{\max}^2rB_L^2\|S\|_{\mathrm{op}}^2
 \le432rL^4\|S\|_{\mathrm{op}}^2.
\]
On the lowest block,
$\sum_i(\operatorname{tr}S-S_{ii})^2
=\sum_iS_{ii}^2+(\operatorname{tr}S)^2$.
Moreover
\[
 \frac{\mathfrak m_L^2}{|\delta_L|}
 =\frac{4\cos^2h}{N^2\sin^2h\cos(2h)}
 >\frac{16}{\pi^2}>1.
\]
Here $\cos(2h)>0$, $\cos^2h>\cos(2h)$ and $\sin h<h$
give the strict inequalities. Consequently
$\operatorname{tr}((\mathsf R_S^{\min})^2)
\ge\sigma^2\delta_L^2\|S\|_{\mathrm F}^2$.
Use $\|S\|_{\mathrm F}\ge\|S\|_{\mathrm{op}}$ to obtain
(124). The upper bound follows because
the full Hilbert--Schmidt sum contains the lowest block.
The displayed lower bound tends to zero at large $L$; it is not
an asserted positive volume-uniform spectral fraction.

\subsection{The original material endpoint and an actual diagonal}

For integers $j\ge2$ retain the original full matrices from Section~13 and set
\begin{equation}\label{eq:low-diagonal}
 L_j=j^2,\quad a_j=\frac1{100j},\quad z_j=\frac1j,\quad
 \tau_j=\frac{1-j^{-2}}8,\quad S_j=K_{\tau_j}=C_{\tau_j}C_{\tau_j}^{\mathsf T}.
\end{equation}
These times stay on the regular material orbit; no singular matrix is
inserted at its endpoint. Each $S_j$ is positive definite because
$C_{\tau_j}$ is invertible. In particular the preceding theorem applies.
Section~13 proves $j^{-6}S_j\to S_*=\eta\eta^{\mathsf T}$,
$\eta=(1/4,3/8,-9/4)^{\mathsf T}$, with every matrix entry retained.
Define $d_j=d_{L_j,a_j}(S_j)$, $\Delta_j=2\sigma_j/a_j$, and
\begin{equation}\label{eq:low-constant}
 \mathcal C_*=
 \frac{204976875}{512\pi^2}+\frac{3982500}{\pi^6}>0 .
\end{equation}

\paragraph{Theorem 14.2.}
There is a definite sequence of strictly positive dyadic couplings
$g_j<1/j$ and actual smooth physical vectors
\begin{equation}\label{eq:low-actual-vector}
 w_j=\mathbf1_{(\Delta_j/2,\,11\Delta_j/10)}
        (H_{g_j,L_j,a_j}-\mathcal E_{g_j,L_j,a_j})
        \,\Xi_{f_{S_j}},
\end{equation}
orthogonal to the actual vacuum, for which
\begin{align}
 j^{-18}\|w_j\|^2&\longrightarrow\mathcal C_*,
 \label{eq:low-raw-norm}\\
 j^{-17}\langle w_j,(H_{g_j,L_j,a_j}-\mathcal E_{g_j,L_j,a_j})w_j\rangle
 &\longrightarrow100\sqrt2\pi\,\mathcal C_*,
 \label{eq:low-raw-energy}\\
 j\,\frac{\langle w_j,(H_{g_j,L_j,a_j}-\mathcal E_{g_j,L_j,a_j})w_j\rangle}
               {\|w_j\|^2}
 &\longrightarrow100\sqrt2\pi .
 \label{eq:low-quotient}
\end{align}
No vector is divided by its norm in this construction.


\paragraph{Proof.}
For each fixed $j$, fixed-box eigenvalue convergence and
(122) imply that every sufficiently small $g>0$
satisfies the following finite set of strict inequalities.
The first six physical excitation energies lie within
$\Delta_j/(10j)$ of $\Delta_j$, the seventh exceeds
$11\Delta_j/10$, and
\[
 \left|\frac{\|P_g(I_{L_j,a_j})\Xi_{f_{S_j}}\|^2}{d_j}-1\right|<\frac1j.
\]
Also require
$|\|\Xi_{f_{S_j}}\|^2/\mathcal N_{L_j,a_j}(S_j)-1|<1/j$,
which follows from the proved full-vector norm convergence.
Let $n_j$ be the least positive integer for which $g_j=2^{-n_j}<1/j$
and all these inequalities hold. Their established fixed-box limits
prove existence of this integer. This definition asserts no unproved
volume-uniform convergence rate or prescribed renormalized coupling law.

The projection in (127) has rank six.
Its range consists of smooth physical eigenvectors; it excludes the
vacuum, since its energy interval is strictly positive.
The norm inequality gives $\|w_j\|^2>(1-1/j)d_j>0$.
On that range every excitation eigenvalue differs from $\Delta_j$
by at most $\Delta_j/(10j)$. Expansion in its actual orthonormal
eigenbasis therefore gives the exact two-sided bound
\[
 \left|
 \frac{\langle w_j,(H-\mathcal E)w_j\rangle}{\Delta_j\|w_j\|^2}-1
 \right|<\frac1{10j}.
\]
The basis is used to evaluate the original vector, not to rescale it.

It remains to compute the unchanged raw scale $d_j$.
For $N_j=2j^2+1$ and $h_j=\pi/(2N_j)$, the limit
$\sin h_j/h_j\to1$ gives
\[
 \frac{\sigma_j^2}{2\pi^2/N_j^2}\to1,\qquad
 \frac{\delta_{L_j}}{N_j^2}\to-\frac1{2\pi^2},\qquad
 \frac{\mathfrak m_{L_j}^2}{N_j^2}\to\frac8{\pi^4},
 \qquad \frac{N_j^2}{a_j^2j^6}\to40000.
\]
The exact tensor entries give
\[
 \sum_i(\operatorname{tr}S_*-S_{*,ii})^2
 =\frac{333^2+328^2+13^2}{64^2}=\frac{109321}{2048},
 \qquad
 \sum_{i<j}S_{*,ij}^2=\frac{531}{512}.
\]
Insert these into the full expression (120), using
$j^{-6}S_j\to S_*$, to obtain
\[
 j^{-18}d_j\to
 \frac{7500}{\pi^2}\frac{109321}{2048}
 +\frac{3840000}{\pi^6}\frac{531}{512}=\mathcal C_* .
\]
Moreover
$j\Delta_j=400\sqrt2\,j^2\sin(\pi/(4j^2+2))
\to100\sqrt2\pi$.
The norm and energy inequalities above now prove all three limits.
The additionally retained whole-state norm gives the actual
positive-coupling fraction bound
\[
 \frac{\|w_j\|^2}{\|\Xi_{f_{S_j}}\|^2}
 \ge\frac{j-1}{j+1}\,
 \frac{\sigma_j^2\delta_{L_j}^2}{432r_jL_j^4}.
\]
Its right side is asymptotic to
$j^{-10}/(3456\pi^2)$. No limit of the full fraction is inferred
from this lower estimate.
Every finite Hamiltonian still has
$\kappa_j=200jg_j^2$, $b_j=50j/g_j^2$ and
$\xi_j=1/(4g_j^4)$, including the full scalar $2b_j|\mathsf P_{L_j}|$.
No strong-coupling inverse estimate was used.



If the chosen observable uses physical midpoints $am(e)$ instead,
its exact state is $a_j^2w_j$ by linearity of $D_f$, centering and
spectral projection. Its raw norm and energy are multiplied by $a_j^4$.
Thus their powers become $j^{14}$ and $j^{13}$ with constants
$\mathcal C_*/10^8$ and $100\sqrt2\pi\,\mathcal C_*/10^8$.
The quotient is identical. This is the full coordinate/unit dictionary,
not a silent change of weight.

The construction proves the full morphism from each material tensor
through the energy-current state and the actual physical spectral
projection. The six-dimensional coefficient map has the explicit
inverse in Theorem~14.1.
The spectral projection itself is global in space and depends on the
full Hamiltonian. It has not been proved to be a local observable map
between the changing regulators, nor a dynamical intertwiner of the
forced Navier--Stokes solution and source-free Yang--Mills.
The compactly supported Navier--Stokes profile is not replaced here by
the stationary Fabel material orbit. The next unresolved part of the
original objective is a common interacting continuum state/observable
construction preserving this low-energy image. The finite-regulator
sequence alone does not supply that identification.
Section~15 constructs the next common
vacuum-plus-band frame with its exact Gram tensor, calculates compact
local-energy and Navier--Stokes curvature-state matrix elements, and
retains the explicit higher-mode leakage in products. The common
full local interacting continuum representation remains unresolved.

\section{Local energy, curvature profiles, and the common lowest-band map}
\label{sec:local-band}

The global spectral projection of Section~14 now receives local
weighted-energy states and observables. We calculate its entire
vacuum-plus-first-band operator, its raw matrix elements against the
full Fabel state, and its two physical time scales. We also retain the
exact terms that leave this band. A compressed product is not substituted
for a product of the original local operators.

\subsection{Graph convergence on the actual excited space}

All definitions and coefficients of Sections~1, 13 and 14 remain in
force. We first extend the vacuum graph calculation to the actual finite
physical eigenspaces. At fixed $L,a$, let $u_g$ be an original physical
eigenvector, with eigenvalue $e_g$ in a bounded interval and bounded
squared norm. Positivity of $u_g$ is not required. For a real smooth
multiplier $F$, expand the full kinetic product and use the eigenvalue
equation, taking its real part. The exact identity is
\begin{equation}\label{eq:local-excited-transform}
 q_H[Fu_g]-e_g\|Fu_g\|^2
 =\kappa\sum_{e,\alpha}\int |u_g|^2|X_{e,\alpha}F|^2\,dU.
\end{equation}
The terms $F^2|Xu_g|^2$ and $2F\operatorname{Re}(\overline{u_g}Xu_g)XF$
occur on both sides of the integration-by-parts eigenvalue identity;
their cancellation proves the formula also for complex $u_g$.

Put $M_k(u_g)=\int W^k|u_g|^2\,dU$. The same regularized powers and
the exact bound $\sum|XW|^2\le16W$ used in Section~14 give
\[
 bM_{k+1}(u_g)\le e_gM_k(u_g)+4\kappa k^2M_{k-1}(u_g),
 \qquad bM_1(u_g)\le e_gM_0(u_g).
\]
Thus for every fixed $k$, $M_k(u_g)/g^{2k}$ is bounded at fixed $L,a$.
The constants are uniform over a bounded-energy unit eigenspace.
If $\mathcal B_gu_g\to u_0$ strongly and $e_g\to e_0$,
the third moment gives the same chart-tail estimate for
$(W/g^2)u_g$ as for the vacuum. Local Taylor convergence proves
\[
 \mathcal B_g((W/g^2)u_g)\to V_2u_0,\qquad
 \mathcal B_g(g^2H_0u_g)\to
          (ae_0/2-V_2/4)u_0=Du_0 .
\]
The second identity is the full eigenvalue equation
$g^2H_0u_g=(ae_g/2)u_g-(W/(4g^2))u_g$.
For every fixed real edge weight $f$, the strongly commuting Casimir
inequality
$\|\sum_ef_eE_eF\|\le\|f\|_\infty\|H_0F\|$
and compactly supported Gaussian-polynomial graph approximations prove
the arbitrary electric-weight graph limit, exactly as in Section~14.
The weighted magnetic limit follows from
$|W_f|\le\max_p|f_p|W$. These arguments prove
\begin{equation}\label{eq:local-graph}
 \mathcal B_gD_fu_g\longrightarrow D_f^{\mathrm{osc}}u_0
\end{equation}
in full $L^2$, with the original factors $2/a$ and $1/(8a)$,
for every eigenvector limit in any fixed finite physical spectral range.

To obtain this result for a finite spectral frame, choose an eigenbasis
of that range along an arbitrary sequence $g\downarrow0$.
The retained compactness and projection theorem gives a subsequence
on which every basis vector converges. Apply the preceding moment and
graph estimates to each vector; bounded finite linear combinations
preserve them. The limit identifies the same operator matrix
independently of the chosen subsequence and eigenbasis. Therefore the
graph conclusion holds for the original finite spectral projections,
including degenerate bands, not merely a chosen eigenvector subsequence.

\subsection{Exact common coordinates and the whole compressed operator}

Let
\[
 \mathcal K=\mathbb C\oplus\operatorname{Sym}_3(\mathbb C),\qquad
 \langle(c,B),(d,E)\rangle_{\mathcal K}
 =\overline c\,d+6\operatorname{tr}(B^*E).
\]
The factor $6$ is the unchanged creation-vector norm:
\[
 \Phi(B)=\sum_{i,k=1}^3B_{ik}
            \sum_{\alpha=1}^3a_{i,\alpha}^\dagger
                            a_{k,\alpha}^\dagger\Phi_0,
 \qquad \langle\Phi(B),\Phi(E)\rangle=6\operatorname{tr}(B^*E).
\]
Commuting two annihilators through two creators proves this formula
for complex symmetric matrices; off-diagonal entries occur twice in
the displayed sum and twice in the matrix trace. Define the comparison
isometry $I_{L,a}(c,B)=c\Phi_0+\Phi(B)$ using the three exact lowest
modes of Section~14.

Write $P_g$ for the first-band projection on
$(\Delta/2,11\Delta/10)$ and $\Pi_g=|\psi_g\rangle\langle\psi_g|+P_g$.
For sufficiently small positive $g$ define
\begin{equation}\label{eq:local-frame}
 M_g(c,B)=c\psi_g+P_g\mathcal B_g^*\Phi(B),\qquad G_g=M_g^*M_g.
\end{equation}
Here every adjoint on $\mathcal K$ uses its displayed factor-$6$
inner product. Projection convergence proves $G_g\to I$ in
operator norm. In particular $G_g$ is positive definite for small
$g$, $M_g$ is a bijection onto the actual seven-dimensional range
of $\Pi_g$, and its exact inverse there is
\[
 M_g^{-1}=G_g^{-1}M_g^*,\qquad
 M_g^{-1}M_g=I,\qquad M_gM_g^{-1}=\Pi_g.
\]
The last equality on the full Hilbert space means the indicated
orthogonal projection. These identities follow by multiplication and
the equality of the range dimensions. No original vector or frame
column is divided by its norm. The complete finite-$g$ Gram tensor
$G_g$ is retained.

Put $C_{g,f}=D_f-\langle\psi_g,D_f\psi_g\rangle$.
Let $J_f=\Sigma^{1/2}A_f\Sigma^{1/2}$ and
$K_f=\Sigma^{-1/2}M_f\Sigma^{-1/2}$ in the notation of
(114); here the face matrix $M_f$ is not the Hilbert
frame $M_g$. Set
\[
 \mathsf R_f=K_f-J_f,\qquad \mathsf Q_f=K_f+J_f.
\]
The full centered oscillator operator is
\begin{equation}\label{eq:local-normal-order}
 C_{0,f}=\frac1{4a}\sum_{\nu,\eta,\alpha}
 \left[(\mathsf R_f)_{\nu\eta}
 (a_{\nu,\alpha}a_{\eta,\alpha}
       +a_{\nu,\alpha}^\dagger a_{\eta,\alpha}^\dagger)
       +2(\mathsf Q_f)_{\nu\eta}
                      a_{\nu,\alpha}^\dagger a_{\eta,\alpha}\right].
\end{equation}
Expand $\partial_z=\sqrt{\sigma}/(2\sqrt2)(a-a^\dagger)$
and $z=\sqrt{2/\sigma}(a+a^\dagger)$ in the full electric and
magnetic quadratic forms to obtain this identity.
Their vacuum scalar is $3\operatorname{tr}(\mathsf Q_f)/(4a)$;
subtracting that scalar gives exactly
(134). No magnetic or number-preserving term
is discarded.

Use superscript $\min$ for the three-lowest-mode blocks. The resulting
exact compressed operator on $\mathcal K$ is
\begin{equation}\label{eq:local-compression}
 \mathcal D_{0,f}(c,B)=
 \left(\frac{3}{2a}\operatorname{tr}(\mathsf R_f^{\min}B),
 \ \frac{c}{4a}\mathsf R_f^{\min}
       +\frac1{2a}(\mathsf Q_f^{\min}B+B\mathsf Q_f^{\min})\right).
\end{equation}
Indeed the pair annihilation gives
$6\operatorname{tr}(\mathsf R_f^{\min}B)/(4a)$ in the vacuum,
pair creation from the vacuum gives the second component, and the
commutator of
$\sum Q_{\nu\eta}a_\nu^\dagger a_\eta$ with $\Phi(B)$ is
$\Phi(QB+BQ^{\mathsf T})$. Higher-mode terms are outside the
compressed range; they will be retained below.
The graph result (132) proves
\begin{equation}\label{eq:local-compression-actual}
 \mathcal D_{g,f}:=
 G_g^{-1}M_g^*\Pi_gC_{g,f}\Pi_gM_g
 \longrightarrow\mathcal D_{0,f}
\end{equation}
in the finite-dimensional operator norm. The original compressed
operator is self-adjoint for the Gram inner product $G_g$:
$G_g\mathcal D_{g,f}=\mathcal D_{g,f}^*G_g$.

\subsection{Compact physical supports and all leading spatial factors}

For real $h\in C_c(\mathbb R^3)$ use the exact local weights
\[
 f_{h,j}(e)=h(a_jm(e)),\qquad
 f_{h,j}(p)=\frac14\sum_{e\in\partial p}h(a_jm(e)),
 \quad L_j=j^2,\quad a_j=\frac1{100j}.
\]
These define a finite sum of the original local electric and magnetic
operators. A face adjacent to the sampled support is retained whenever
its four-edge average is nonzero. Put
\[
 N_j=2L_j+1,\quad \ell_j=a_jN_j,\quad
 H_h=\int_{\mathbb R^3}h(x)\,dx,\quad
 H_{h,ik}=\int_{\mathbb R^3}x_ix_kh(x)\,dx.
\]
The vertex endpoints remain at $\pm L_ja_j$; $\ell_j$ is the
specified difference-basis length $a_j(2L_j+1)$, not a replacement
for the endpoint distance.

\paragraph{Theorem 15.1.}
For the full comparison matrices of these local weights,
\begin{align}
 \frac{\ell_j^4}{a_j}\mathsf R_{h,j}^{\min}
 &\longrightarrow4\sqrt2\pi H_h I_3,
 &
 \frac{\ell_j^4}{a_j}\mathsf Q_{h,j}^{\min}
 &\longrightarrow4\sqrt2\pi H_h I_3, \label{eq:local-leading}\\
 \frac{\ell_j^6}{a_j}(\mathsf R_{h,j}^{\min})_{ik}
 &\longrightarrow2\sqrt2\pi^3 H_{h,ik}
 &&(i\ne k). \label{eq:local-quadrupole}
\end{align}
In particular the complete raw first-band mass obeys
\begin{equation}\label{eq:local-mass}
 \ell_j^8d_{L_j,a_j}(f_{h,j})\longrightarrow36\pi^2 H_h^2.
\end{equation}
Here $d(f)$ denotes $3\operatorname{tr}((\mathsf R_f^{\min})^2)/(8a^2)$
for arbitrary weights, extending the quadratic-weight notation of
(120).

\paragraph{Proof.}
Let $V_i$ be the original edge functions (117) and
$Y_i=d_1V_i/\sigma_j$. Put
\[
 A^h_{ik}=\sum_ef_{h,j}(e)V_i(e)V_k(e),\qquad
 F^h_{ik}=\sum_pf_{h,j}(p)Y_i(p)Y_k(p).
\]
Then, exactly,
$\mathsf R_h^{\min}=\sigma_j(F^h-A^h)$ and
$\mathsf Q_h^{\min}=\sigma_j(F^h+A^h)$.
Different $Y_i$ have disjoint face-plane supports, so $F^h_{ik}=0$
for $i\ne k$ with any face weight.
For the plane $(p,q)$ normal to $i$,
\[
 Y_i(n)^2=\frac4{N_j^3}
 \cos^2\left(\frac{\pi a_j(n_p+1/2)}{\ell_j}\right)
 \cos^2\left(\frac{\pi a_j(n_q+1/2)}{\ell_j}\right).
\]
The sign of the ordered curl disappears here by squaring; it was not
changed. On every fixed physical compact set the two cosine factors
tend uniformly to one. The quarter-face sample differs uniformly
from $h$ at the face center by at most its modulus of continuity
at distance $a_j/2$. The support increases by at most that distance.
The face-center Riemann sums with volume factor $a_j^3$ consequently
give
\[
 \ell_j^3 F^h_{ii}\longrightarrow4H_h.
\]
For clarity, uniform continuity bounds each sample error, the number
of nonzero samples times $a_j^3$ stays bounded in a fixed enlarged
compact cube, and the ordinary shifted cubic partition proves
convergence of the remaining Riemann sum. Thus the four edges and
the open endpoints contribute no unrecorded limit term.

For $V_i$ on direction $p$ its squared factor is
\[
 \frac2{N_j^3}
 \cos^2\left(\frac{\pi a_j(n_p+1/2)}{\ell_j}\right)
 \sin^2\left(\frac{\pi a_jn_q}{\ell_j}\right),
\]
and there is the other directional term with $p,q$ interchanged.
On a fixed support, $|\sin(\pi x/\ell_j)|\le\pi|x|/\ell_j$.
The same sample-count estimate gives $A^h_{ii}=O_h(\ell_j^{-5})$.
Cauchy--Schwarz with $|h|$ also bounds every
$A^h_{ik}=O_h(\ell_j^{-5})$.

More precisely, distinct normal indices $i,k$ share only the remaining
edge direction $d$. On that edge the exact product is
\[
 V_i(n,d)V_k(n,d)=
 -\frac2{N_j^3}
 \sin\left(\frac{\pi a_jn_i}{\ell_j}\right)
 \sin\left(\frac{\pi a_jn_k}{\ell_j}\right)
 \cos^2\left(\frac{\pi a_j(n_d+1/2)}{\ell_j}\right).
\]
Keeping the two original signs and using uniform convergence
$\ell_j\sin(\pi x/\ell_j)\to\pi x$ on the fixed support proves
\[
 \ell_j^5 A^h_{ik}\longrightarrow-2\pi^2H_{h,ik}.
\]
Finally $\ell_j\sigma_j/a_j=N_j\sqrt8\sin(\pi/(2N_j))
\to\sqrt2\pi$. These identities prove both matrix limits.
Insert the first into the exact raw mass formula to get
$3\cdot3(4\sqrt2\pi)^2H_h^2/8=36\pi^2H_h^2$.



This calculation retains spatial information beyond the leading
integral: (138) is an explicit map to three
mixed spatial moments. A vanishing integral does not imply that
the original weighted state or these next coefficients vanish.

\subsection{One positive-coupling diagonal and its raw common representation}

The finite-regulator convergence can be enforced for every bounded
edge weight at once at each fixed $j$. There are finitely many
original edges. Let $b^{(e)}$ be the weight one at $e$ and zero
elsewhere, with its prescribed face averages. Both the actual
centered operator and its comparison matrix are linear in this
weight. Choose
\[
 \epsilon_j=j^{-40}(1+\ell_j)^{-10}.
\]
At each fixed $j$, take the least positive integer $n_j$ such that
$g_j=2^{-n_j}<1/j$ satisfies the finite inequalities of
Theorem~14.2, $\|G_{g_j}-I\|<\epsilon_j$,
and
\[
 \|\mathcal D_{g_j,b^{(e)}}-\mathcal D_{0,b^{(e)}}\|
       <\epsilon_j/|\mathsf E_{L_j}|\quad(e\in\mathsf E_{L_j}).
\]
Also require the exact frame energy matrix to differ from
$0\oplus\Delta_jI$ by less than $\epsilon_j$, and the frame
coefficient of the full projected Fabel state to differ from
$(0,\mathsf R_{S_j}^{\min}/(4a_j))$ by less than $\epsilon_j$.
Projection, graph and eigenvalue convergence prove existence for
this finite list. This is a defined stricter coupling sequence;
it does not silently reuse a previously chosen $g_j$ for new
inequalities. All conclusions of Section~14 persist.

For every real edge weight $f$, finite linearity now gives the
uniform bound
\[
 \|\mathcal D_{g_j,f}-\mathcal D_{0,f}\|
 \le\epsilon_j\|f\|_\infty.
\]
The energy and Fabel coordinate requirements also follow from the
fixed-box frame and vector convergence already proved.
All actual states and operators retain their original coefficients
$\kappa_j=200jg_j^2$, $b_j=50j/g_j^2$, and full Wilson scalar.

\paragraph{Theorem 15.2.}
For every real $h\in C_c(\mathbb R^3)$, the actual Gram and
compressed operator coordinates satisfy
\[
 G_{g_j}\to I,\qquad
 \ell_j^4\mathcal D_{g_j,f_{h,j}}\to\mathcal T_h
\]
in operator norm on the fixed space $\mathcal K$, where
\begin{equation}\label{eq:local-T}
 \mathcal T_h(c,B)=
 \sqrt2\pi H_h\bigl(6\operatorname{tr}B,\ cI_3+4B\bigr).
\end{equation}
The actual excitation matrix $\mathcal A_j$ satisfies
\begin{equation}\label{eq:local-time}
 \mathcal A_j\to0,\qquad
 \ell_j\mathcal A_j\to
       \mathcal A_*:=0\oplus2\sqrt2\pi I.
\end{equation}
These limits include all finite products of the displayed compressed
operators and their compressed time evolution.

\paragraph{Proof.}
Apply (137) to each term of
(135), and use
$\ell_j^4\epsilon_j\|h\|_\infty\to0$.
The first component becomes $6\sqrt2\pi H_h\operatorname{tr}B$;
the other two become $\sqrt2\pi H_hcI$ and
$4\sqrt2\pi H_hB$. This proves (140), including its
adjoint relation in the original factor-$6$ metric.
The exact energy frame is block diagonal because the vacuum and
excited range are invariant and orthogonal. Its specified error
and $\ell_j\Delta_j\to2\sqrt2\pi$ give (141).
Operator-norm convergence in a fixed finite-dimensional space
preserves finite products. The exponential power series, uniformly
bounded on each fixed compact time interval, gives convergence of
the displayed time evolutions.



The map $h\mapsto\mathcal T_h$ is explicitly the integral functional
times the single displayed operator. On the traceless matrices its
eigenvalue is $4\sqrt2\pi H_h$; on the vacuum/trace coordinates its
matrix is $\sqrt2\pi H_h\left(\begin{smallmatrix}0&18\\1&4\end{smallmatrix}\right)$
in the basis $(1,0),(0,I)$, whose squared norms are $1,18$.
Its characteristic roots are
$\sqrt2\pi H_h(2+\sqrt{22})$ and
$\sqrt2\pi H_h(2-\sqrt{22})$. This is a finite-dimensional
self-adjoint compressed observable with all raw metric factors,
not an assertion that the full local-energy algebra has collapsed.

Let $v_{h,j}=P_{g_j}\Xi_{f_{h,j}}$. The exact vacuum component
of $\Xi$ vanishes. Equations (133)--(140)
therefore give its coordinate limit
\[
 \ell_j^4M_{g_j}^{-1}v_{h,j}
 \longrightarrow(0,\sqrt2\pi H_hI).
\]
This equation records the raw amplitude $\ell_j^{-4}$.
It defines no unit-vector rescaling. In particular, for real compact
$h,k$ and fixed physical Euclidean time $t\ge0$,
\begin{align}
 \ell_j^8\langle v_{h,j},e^{-t(H_j-\mathcal E_j)}v_{k,j}\rangle
 &\longrightarrow36\pi^2H_hH_k,\label{eq:local-fixed-time}\\
 \ell_j^8\langle v_{h,j},
       e^{-\ell_j\theta(H_j-\mathcal E_j)}v_{k,j}\rangle
 &\longrightarrow36\pi^2H_hH_k e^{-2\sqrt2\pi\theta}
 \quad(\theta\ge0).\label{eq:local-slow-time}
\end{align}
These follow by inserting the exact Gram $G_{g_j}$ and the exact
energy matrix in the inner product. The second equation records
the explicit time change $t=\ell_j\theta$; it is not used to replace
the original Hamiltonian or its physical energy quotient.
At fixed physical time this particular finite band becomes a
zero-energy sector. At times proportional to the box length its
nontrivial dynamics are the explicitly displayed finite matrix.
The same actual energy matrix gives the raw first-moment limit
\[
 \ell_j^9\langle v_{h,j},(H_j-\mathcal E_j)v_{h,j}\rangle
 \longrightarrow72\sqrt2\pi^3H_h^2.
\]
When $H_h\ne0$, the raw norm limit is strictly positive after its
displayed $\ell_j^8$ factor, so $v_{h,j}\ne0$ eventually and
\[
 \ell_j\frac{\langle v_{h,j},(H_j-\mathcal E_j)v_{h,j}\rangle}
                   {\|v_{h,j}\|^2}\longrightarrow2\sqrt2\pi.
\]
The vectors themselves retain their vanishing raw norms.

For the unchanged Fabel states $w_j$ of Section~14 let
$S_*=\eta\eta^{\mathsf T}$ as there, and define the full fixed matrix
\[
 (B_*)_{ii}=-\frac{25\sqrt2}{\pi}
                    (\operatorname{tr}S_*-S_{*,ii}),\qquad
 (B_*)_{ik}=\frac{400\sqrt2}{\pi^3}S_{*,ik}\quad(i\ne k).
\]
The exact lowest-mode coefficients give
$M_{g_j}^{-1}w_j=j^9(0,B_*+o(1))$ and
$6\operatorname{tr}(B_*^2)=\mathcal C_*$.
Consequently the original raw matrix elements obey
\begin{align}
 \frac{\ell_j^4}{j^{18}}
 \langle w_j,C_{g_j,f_{h,j}}w_j\rangle
 &\longrightarrow4\sqrt2\pi H_h\mathcal C_*,
 \label{eq:local-fabel-diag}\\
 \frac{\ell_j^4}{j^9}
 \langle\psi_j,C_{g_j,f_{h,j}}w_j\rangle
 &\longrightarrow-\frac{25275}{8}H_h.
 \label{eq:local-fabel-mixed}
\end{align}
The first follows from the band component $4\sqrt2\pi H_hB$ of
(140). For the sign and constant in the second,
$\operatorname{tr}S_*=337/64$ and
$\operatorname{tr}B_*=-8425\sqrt2/(32\pi)$; multiplication by
$6\sqrt2\pi H_h$ yields the result.
Thus the quotient of the first raw expectation by $\|w_j\|^2$
decays as $4\sqrt2\pi H_h/\ell_j^4$ when $H_h\ne0$,
even though the unscaled raw expectation itself grows.
Neither conclusion is inferred by dropping the Fabel amplitude.

\subsection{The map from a compact Navier--Stokes curvature profile}

Take any actual smooth compactly supported incompressible velocity
profile $u(s,\cdot)$ at a regular time $s$ in the Navier--Stokes
construction under study. Its viscosity, force, pressure and time
interval remain the ones of that construction. Set
\[
 \mathcal A_0=0,\quad\mathcal A_i=\lambda u_iT,\quad
 T=-i\sigma_3/2,\qquad
 h_s(x)=\sum_{i<k}-2\operatorname{tr}(F_{ik}(s,x)^2)
                =\lambda^2|\operatorname{curl}u(s,x)|^2.
\]
Here $F_{ik}=\lambda(\partial_iu_k-\partial_ku_i)T$ and
$-2\operatorname{tr}(T^2)=1$ prove the equality.
The weighted local-energy operator uses the full profile:
$f_{s,j}(e)=h_s(a_jm(e))$, with all prescribed face averages.
The composition
\[
 u(s,\cdot)\longmapsto h_s\longmapsto D_{f_{s,j}}
 \longmapsto\Xi_{f_{s,j}}\longmapsto P_{g_j}\Xi_{f_{s,j}}
\]
is now a specified map into the actual physical Hilbert space.
Its raw first-band masses and mixed-time spectral matrix elements
are (139), (142) and
(143), with
\[
 H_{h_s}=\lambda^2\|\operatorname{curl}u(s,\cdot)\|_2^2,
 \qquad
 H_{h_s,ik}=\lambda^2\int x_ix_k|\operatorname{curl}u(s,x)|^2\,dx.
\]
For $\lambda\ne0$ and a nonzero compact smooth incompressible profile,
$H_{h_s}>0$. Indeed integration by parts gives
$\int|\nabla u|^2=\int|\operatorname{curl}u|^2+
\int|\operatorname{div}u|^2$; if its curl vanishes, all derivatives
vanish and compact support forces $u=0$.
Thus this particular local-curvature state has nonzero first-band
mass at sufficiently late regulators, with all norm powers explicit.

This is a regular-time profile map. It does not assert that
$H_{h_s}$ diverges as a blowup time is approached: pointwise
concentration and its spatial integral remain distinct quantities
in the displayed exact map. The full sourced current
$j_i=(\lambda/g^2)(\Delta u_i-c^{-2}\partial_s^2u_i)T$
from Section~13 is unchanged. No equality between evolution in
the fluid parameter $s$ and the quantum time $t$ or $\theta$ is
used or claimed. The limits above apply to every fixed regular
profile; a simultaneous singular-fluid-time limit is not exchanged
with the regulator limits.
Only the nonzero compact incompressible profile, rather than its
singular endpoint, was used to prove positive band mass. The
gap-closing projected sequence therefore does not establish an
implication specific to Navier--Stokes breakdown.

\subsection{All terms leaving the seven-dimensional range}

Let $\widehat Q_g=I-\Pi_g$. On the original smooth spectral range,
the exact identity is
\begin{equation}\label{eq:local-leakage}
 \Pi_g C_{g,h}C_{g,k}\Pi_g
 =\Pi_gC_{g,h}\Pi_gC_{g,k}\Pi_g+
       (\widehat Q_gC_{g,h}\Pi_g)^*
                        (\widehat Q_gC_{g,k}\Pi_g).
\end{equation}
Each finite-range eigenvector is smooth, each weighted differential
operator preserves smoothness, and its formal adjoint is its
self-adjoint realization there. Inserting $I=\Pi_g+\widehat Q_g$
therefore proves the identity with no domain omission.
For $h=k$ the last term is positive. Products of the compressed
limits alone do not determine this term.

The fixed-box graph limit proves the exact comparison for every
leakage matrix element as well. Here are its full coordinates.
Extend a lowest-mode symmetric $B$ by zero to all modes and put
$\mathcal S_B= \sum_{\nu,\eta,\alpha}
B_{\nu\eta}a_{\nu,\alpha}^\dagger a_{\eta,\alpha}^\dagger$.
Let $\mathsf P_\perp$ remove the lowest $3$-by-$3$ block of a
symmetric mode matrix. From (134),
\begin{align}
 \widehat Q_0C_{0,f}\Phi_0
 &=\frac1{4a}\Phi(\mathsf P_\perp\mathsf R_f),
 \label{eq:local-leak-vac}\\
 \widehat Q_0C_{0,f}\Phi(B)
 &=\frac1{4a}\mathcal S_{\mathsf R_f}\mathcal S_B\Phi_0
 +\frac1{2a}
 \Phi\!\left(\mathsf P_\perp(\mathsf Q_fB+B\mathsf Q_f)\right).
 \label{eq:local-leak-pair}
\end{align}
In these equations $\Phi$ on a full mode matrix uses the same
colour-contracted formula over every spatial mode.
The two terms on the right in (148) have
respectively four and two creation operators, hence are orthogonal.
The annihilation term has landed in the vacuum and is included in
$\Pi_0$, so it is absent for that exact reason.
In particular
\[
 \|\widehat Q_0C_{0,f}\Phi_0\|^2
 =\frac3{8a^2}\left[
 \operatorname{tr}(\mathsf R_f^2)
           -\operatorname{tr}((\mathsf R_f^{\min})^2)\right].
\]
The four-creation contraction also has a completely explicit finite
sum. With combined indices $I=(\nu,\alpha)$ define
$T_{I_1I_2I_3I_4}
=(\mathsf R_f)_{\nu_1\nu_2}B_{\nu_3\nu_4}
\delta_{\alpha_1\alpha_2}\delta_{\alpha_3\alpha_4}$.
For another pair of matrices producing coefficients $U$,
\[
 \left\langle\mathcal S_{\mathsf R_f}\mathcal S_B\Phi_0,
              \mathcal S_{\mathsf R_k}\mathcal S_E\Phi_0\right\rangle
 =\sum_{I_1,\ldots,I_4}\sum_{J_1,\ldots,J_4}
 \overline{T_{I_1I_2I_3I_4}}U_{J_1J_2J_3J_4}
 \sum_{\pi\in S_4}\prod_{r=1}^4\delta_{I_rJ_{\pi(r)}}.
\]
Repeatedly commute the four annihilators through the four creators;
every surviving vacuum term is one of these $24$ pairings, proving
the formula with its multiplicities and colour factors.
The two-creation term has the already-proved factor-$6$ trace
inner product. These equations give every omitted-intermediate
covariance in (146) at fixed box and coupling limit.

We have constructed a common finite-dimensional representation of
the exact compressed local-energy maps, including their raw coupling
to Fabel and compact Navier--Stokes profiles. Its fixed physical-time
energy is zero and its leading spatial functional is the integral,
with explicitly retained next mixed moments. Section~16 evaluates
the full unprojected vacuum two-creation spectral measure in the
joint infinite-volume continuum limit, using the exact mode maps
above. Products on excited states still retain the four-creation
terms displayed here. Their continuum products and the survival of
nonabelian interactions require further calculation before this
sequence can identify a local interacting four-dimensional
Yang--Mills theory.

\section{The full local-energy spectral measure in the expanding-box limit}
\label{sec:full-local-spectrum}

We now retain every transverse mode of the original open box. The
observable remains the original weighted electric and magnetic energy,
with the original face average. Positive Euclidean time acts by the
original excitation semigroup. It is an explicitly recorded operation
on the vector; no vector is divided by its norm. The result below
constructs a common two-creation spectral representation, including its
continuous low-energy tail. The representation is identified explicitly
with three colour copies of a transverse free field. This identification
is part of the result and does not identify that representation with an
interacting four-dimensional Yang--Mills theory.

\subsection{All-mode coordinates and the observable to be taken to the limit}

Take real $h\in C_c^\infty(\mathbb R^3)$, and put
\[
 h_e=h(a m(e)),\qquad h_p=\frac14\sum_{e\in\partial p}h_e,\qquad
 \Xi_{h,g}=(D_h-\langle\psi_g,D_h\psi_g\rangle)\psi_g,
 \quad A_g=H_g-\mathcal E_g .
\]
All original edges and faces occur in $D_h$, even where the weight is
zero. In particular, faces adjacent to the support are retained. For
$t>0$ define the unchanged, time-filtered state
\begin{equation}\label{eq:full-damped-state}
 y_{h,g}(t)=e^{-tA_g/2}\Xi_{h,g},\qquad
 C_{g;h,k}(t)=\langle\Xi_{h,g},e^{-tA_g}\Xi_{k,g}\rangle.
\end{equation}
The inner product is conjugate-linear in the first argument. These
vectors are physical and orthogonal to the vacuum, since each operation
commutes with the original gauge action and with the vacuum projection.

Write $V_\nu=RO e_\nu$, as in Section~14, for the real transverse
edge eigenvectors with counting norm one and $d_1^{\mathsf T}d_1$
eigenvalue $\sigma_\nu^2$. Put $\omega_\nu=\sigma_\nu/a$. The full pair
matrix divided by the retained spacing is exactly
\begin{align}
 B_{h;\nu\eta}^{L,a}:=\frac{\mathsf R_{h;\nu\eta}}{a}
 &= \frac{\sum_p h_p(d_1V_\nu)_p(d_1V_\eta)_p}
                  {a^2\sqrt{\omega_\nu\omega_\eta}}
       -\sqrt{\omega_\nu\omega_\eta}\sum_e h_e V_\nu(e)V_\eta(e).
 \label{eq:full-pair-native}
\end{align}
Both terms have been carried from the full original operator; neither
is dropped in the limit. Section~14's graph theorem and its exact Fock
norm give, at fixed $L,a$, the finite-measure limit as $g\downarrow0$,
\begin{equation}\label{eq:full-finite-measure}
 d\nu_{h,0}^{L,a}(\omega)
 =\frac38\sum_{\nu,\eta}|B_{h;\nu\eta}^{L,a}|^2
                 \delta_{\omega_\nu+\omega_\eta}(d\omega),
 \qquad C_{0;h,h}^{L,a}(t)=\int e^{-t\omega}d\nu_{h,0}^{L,a}(\omega).
\end{equation}
The sum is ordered, so an off-diagonal pair occurs twice. Three colours
and the two bosonic contractions give $3\cdot2/4^2=3/8$.
Total mass converges at fixed box by strong convergence of the full
weighted vacuum vector. Finite spectral projections converge as proved
in Section~14. To obtain convergence for a bounded continuous spectral
function, first truncate to a finite union of spectral intervals and
then use the convergent vector norms to make the remaining mass small.
This proves in particular convergence of every expression
(149) at fixed $L,a$. No unbounded first moment is
inferred from weak convergence alone.

\subsection{Exact open-box modes, parity transform, and convergence of the full sum}

Here $N=2L+1$ and $\ell=aN$, retaining the distinction between this
basis length and the distance between the extreme vertices. For
$j=0,\ldots,N-1$ define on vertices
\[
 v_j(n)=\sqrt{\frac{2-\delta_{j0}}N}
       \cos\!\left(\frac{\pi j(n+L+1/2)}N\right),
 \quad s_j=2\sin\frac{\pi j}{2N}.
\]
For $j>0$, the edge function is
\[
 w_j(n)=\frac{v_j(n+1)-v_j(n)}{s_j}
       =-\sqrt{\frac2N}\sin\!\left(\frac{\pi j(n+L+1)}N\right).
\]
There is no $w_0$. Finite geometric sums give
$\sum v_jv_k=\delta_{jk}$ and $\sum w_jw_k=\delta_{jk}$ on their
respective original vertex and edge ranges. The displayed difference
identity gives $d_0v_j=s_jw_j$ and, by transposition,
$d_0^{\mathsf T}w_j=s_jv_j$ for $j>0$.

For a triple $\boldsymbol j$, let $I=\{i:j_i>0\}$ and
$s=(s_{j_1},s_{j_2},s_{j_3})$, with missing components zero. Each unit
vector $\epsilon\in\mathbb R^I$ satisfying $s\cdot\epsilon=0$
gives the exact transverse edge cochain
\begin{equation}\label{eq:full-mode-cochain}
 V_{\boldsymbol j,\epsilon}(n,i)
   =\epsilon_i w_{j_i}(n_i)\prod_{d\ne i}v_{j_d}(n_d),
 \qquad \sigma=|s|.
\end{equation}
Its curl on the $(i,k)$ face is
$(s_{j_i}\epsilon_k-s_{j_k}\epsilon_i)
w_{j_i}(n_i)w_{j_k}(n_k)\prod_{d\ne i,k}v_{j_d}(n_d)$.
The preceding orthogonality and difference identities prove its norm,
transversality, and curl eigenvalue. For each triple there are $|I|-1$
such polarizations if $|I|\ge2$, and none otherwise. Decomposing the
tensor product complex into these finite frequency blocks proves
completeness: each block of one-forms is $\mathbb R^I$, its gradient
line is $\mathbb Rs$, and its transverse complement is precisely the
one just constructed. This reproduces all $r$ columns of $RO$ up to
orthogonal changes inside eigenspaces, which leave
(151) unchanged.

\paragraph{Theorem 16.1.}
For any $a\downarrow0$, $\ell=a(2L+1)\longrightarrow\infty$, any
real $h\in C_c^\infty(\mathbb R^3)$, and every $t>0$, the complete
sum in (151) converges to
\begin{equation}\label{eq:full-momentum-integral}
 C_h(t)=\frac{3}{4(2\pi)^6}
 \int_{\mathbb R^3}\int_{\mathbb R^3}
 |p||q|(1+\widehat p\cdot\widehat q)^2
 |\widehat h(p+q)|^2 e^{-t(|p|+|q|)}\,dp\,dq,
 \quad \widehat h(k)=\int e^{-ik\cdot x}h(x)\,dx.
\end{equation}
The values assigned to unit directions at $p=0$ or $q=0$ do not
affect the integral. All positive-time spectral moments converge as
well. Real mixed covariances follow by polarization.

\paragraph{Proof.}
We give bounds for the complete sum before passing to any bounded
frequency region. Each component of $V$ has modulus at most
$2\sqrt2N^{-3/2}$, and each component of $d_1V$ has modulus at most
$4N^{-3/2}\sigma$. Only $O_h(a^{-3})$ original edges or faces have
nonzero weight once $a\le1$ and the box contains a fixed unit
neighbourhood of the support. Consequently
\begin{equation}\label{eq:full-uniform-pair-bound}
 |B_{h;\nu\eta}^{L,a}|
       \le K_h\ell^{-3}\sqrt{\omega_\nu\omega_\eta},
\end{equation}
where $K_h$ is finite and independent of $a,L$. For clarity, one may
take a cube $[-R_h,R_h]^3$ containing that neighbourhood, bound its
edge and face counts by $6(2R_h/a+3)^3$, use
$|h_p|\le\|h\|_\infty$, and insert the component bounds in
(150); this gives such a $K_h$ directly.
Because $2x/\pi\le\sin x\le x$ on $[0,\pi/2]$,
\[
 \frac{2|\boldsymbol j|}{\ell}\le\omega_{\boldsymbol j}
       \le\frac{\pi|\boldsymbol j|}{\ell}.
\]
Grouping integer triples in shells $m\le|\boldsymbol j|<m+1$,
whose counts are at most $K(m+1)^2$, proves for $\ell\ge1$
and fixed $u>0$ that
$\ell^{-3}\sum_\nu\omega_\nu e^{-u\omega_\nu}\le K_u$.
Indeed the bound is a constant times
$\ell^{-4}\sum_{m\ge1}(m+1)^3e^{-2um/\ell}$; comparison with the
integral on each interval $[m,m+1]$ bounds it uniformly for
$\ell\ge1$. It follows from (154) that
the part with $\omega_\nu+\omega_\eta>R$ is bounded by
$K_{h,t}e^{-tR/2}$, uniformly in the regulators. Multiplying by any
fixed power of the pair frequency preserves a tail tending to zero,
by applying this estimate with a smaller positive time.

The part with $\omega_\nu<\varepsilon$ tends to zero uniformly in
the limiting sense needed here: the same shell count gives
\[
 \limsup_{\ell\to\infty}\ell^{-3}
       \sum_{\omega_\nu<\varepsilon}\omega_\nu
       \le K\varepsilon^4.
\]
The analogous estimate holds in the other index. In a bounded
frequency region, triples having any $j_i=0$ contribute at most
$K_R\ell^2$ modes, so their contribution to the paired sum is
$O_{h,R}(\ell^{-1})$. Strips of width $\delta$ adjacent to coordinate
planes have limiting contribution $O_{h,R}(\delta)$ by counting the
allowed integers in that coordinate. These bounds permit us to work
first where every momentum coordinate is positive and bounded away
from zero, and both total frequencies are bounded above and below.

Set $k_i=\pi j_i/\ell$ and $\beta_i=j_i\bmod2$. At the physical
location of its component, (152) is exactly a
product of
$\cos(k_ix_i+\pi j_i/2)$ and $-\sin(k_ix_i+\pi j_i/2)$,
times $\sqrt8N^{-3/2}\epsilon_i$. For fixed bounded frequencies,
the lattice frequency vector $s/a$ converges uniformly to $k$.
Choose the two polarizations by Gram--Schmidt applied to a fixed
coordinate vector and the direction, on a region where these are not
parallel. Finitely many such regions cover the retained compact
frequency set; alternatively the polarization sum uses directly the
matrix $I-kk^{\mathsf T}/|k|^2$. Thus all formulas below are
independent of the choice across overlaps.

The elementary exponential transform of the product is
\begin{equation}\label{eq:full-parity-unitary}
 \ell^{3/2}a^{-3/2}V_{\boldsymbol j,\epsilon}(x)
   =\chi_{\boldsymbol j}\sum_{\rho\in\{-1,1\}^3}
       U_{\beta\rho}\,\epsilon^\rho e^{i(\rho k)\cdot x},
 \quad U_{\beta\rho}=\frac{i}{\sqrt8}e^{i\pi\rho\cdot\beta/2},
 \quad \epsilon^\rho_i=\rho_i\epsilon_i,
 \quad \chi_{\boldsymbol j}=(-1)^{\sum_i\lfloor j_i/2\rfloor}.
\end{equation}
The left side means the vector of component expressions evaluated
at a common physical $x$, whose restrictions are the original edge
values; no relocation of the original samples is made in the sums.
For each pair of signs, summing
$e^{i\pi(\rho_i-\rho_i')\beta_i/2}$ over $\beta_i=0,1$ gives two
when the signs agree and zero otherwise. Hence $U^*U=UU^*=I_8$.
The sign $\chi$ remains in every individual mode and drops out of
the squared modulus of its pair coefficient.

Ordinary Riemann sums in physical space now apply uniformly on the
retained bounded frequency set. Derivatives of each trigonometric
factor and of $h$ are bounded there, so replacing the integral on
each spacing-$a$ cell by its original sample has total error at most
$K_{h,R}a$. For the magnetic term, the quarter-edge face weight
differs from the value at the face centre by at most
$a\|\nabla h\|_\infty$, and the exact discrete curl coefficients
converge to $ik\times\epsilon$. This proves convergence of both
terms in (150), with uniform error
$o(1)\ell^{-3}$ per mode pair. There are $O_R(\ell^6)$ pairs,
so the squared-sum error tends to zero by the uniform coefficient
bound. No derivative or magnetic coefficient is discarded.

For each parity, the momentum grid has step $2\pi/\ell$ in every
coordinate. The eight grids can be aligned with one such grid:
their coordinate differences are at most $\pi/\ell$, and the
preceding compact-set uniform continuity makes the resulting error
tend to zero. Thereafter the sum over the two parities of a mode
pair is the Hilbert--Schmidt norm under $U\otimes U$ in
(155). Its unitarity replaces the two
standing-wave parity sums by the two sums over momentum signs.
The factor $\ell^{-6}$ in the squared pair coefficient and the
two grid densities $(\ell/(2\pi))^3$ leave exactly $(2\pi)^{-6}$.
The sign sums cover the whole momentum space, except null coordinate
planes already controlled above.

For real unit transverse polarizations $e_r(p),e_s(q)$, the remaining
plane-wave pair coefficient, without its Fourier measure factor, is
\begin{equation}\label{eq:full-plane-kernel}
 -\widehat h(-p-q)\left[
 \frac{ (p\times e_r(p))\cdot(q\times e_s(q))}{\sqrt{|p||q|}}
          +\sqrt{|p||q|}\,e_r(p)\cdot e_s(q)\right].
\end{equation}
The two factors of $i$ in the two curls produce the first minus
sign; the electric term supplies the second term with the same
minus sign. In particular the bracket vanishes for $q=-p$ with
matched polarizations, as required by the exact identity $D_1=H$.
For the polarization sum rotate $p$ to the third axis and $q$ to
the first--third plane, and put $c=\widehat p\cdot\widehat q$.
Take $e_1(p)=(1,0,0)$, $e_2(p)=(0,1,0)$,
$e_1(q)=(c,0,-\sqrt{1-c^2})$, $e_2(q)=(0,1,0)$.
The bracket divided by $\sqrt{|p||q|}$ is the diagonal matrix
$(1+c)I_2$. Thus its squared polarization sum is
$2|p||q|(1+c)^2$. Rotational invariance extends this calculation
to every direction; its endpoints follow by continuity.
Multiplication by $3/8$ proves (153).
The previously proved tails justify removal of all cutoffs and
also prove every positive-time moment assertion.



\subsection{Evaluation of the measure and a common physical state map}

\paragraph{Theorem 16.2.}
The locally finite raw measure associated to
(153) is absolutely continuous, with
\begin{align}
 d\nu_h(\omega)&=\rho_h(\omega)\,d\omega,&
 \rho_h(\omega)&=\frac{1}{320\pi^5}
       \int_{|k|\le\omega}|k|^4|\widehat h(k)|^2\,dk,
       \qquad \omega\ge0,\label{eq:full-density}\\
 C_h(t)&=\frac{1}{320\pi^5t}
       \int_{\mathbb R^3}|k|^4 e^{-t|k|}|\widehat h(k)|^2\,dk .
       \label{eq:full-laplace}
\end{align}
For nonzero compact smooth $h$, its support meets every interval
$(0,\varepsilon)$ and it has no atom at zero. Its total mass is
infinite, but $e^{-t\omega}d\nu_h$ has finite mass and all moments
for every $t>0$.

\paragraph{Proof.}
In (153) set $k=p+q$, $r=|p|$,
$s=|q|$, $K=|k|$ and $\omega=r+s$. For $K>0$, align $k$ with
the third axis. The change of the polar angle of $p$ into $s$
has absolute Jacobian $s/(Kr)$. Integration over the azimuth and
over $\delta(\omega-r-s)$ therefore gives the measure
$(2\pi/K)rs\,dr$. Set $u=r-s=2r-\omega$. The admissible range
is $-K\le u\le K$ precisely when $\omega\ge K$; its boundary
has zero measure. The cosine of the angle between $p$ and $q$
satisfies $2rs(1+c)=K^2-u^2$. The complete inner integral is
\[
 \frac{2\pi}{K}\int \frac{ (K^2-u^2)^2}{4}\,dr
 = \frac{\pi}{4K}\int_{-K}^K(K^2-u^2)^2\,du
 = \frac{4\pi K^4}{15}.
\]
The set $k=0$ is null in the outer integral, and the expression
extends there by zero. Tonelli's theorem applies to the
nonnegative integrand. Since
$[3/(4(2\pi)^6)](4\pi/15)=1/(320\pi^5)$, this proves
(157). Integrating $e^{-t\omega}$ from $K$
to infinity proves (158).

For completeness, compact support gives an everywhere convergent
Taylor series for $\widehat h$, by integrating the exponential
series on the fixed support. If all Taylor coefficients vanished,
this series would give $\widehat h=0$ everywhere and hence $h=0$.
One direct proof of the last implication is to convolve $h$ with
a Gaussian: the elementary Gaussian Fourier integral and Fubini
give a zero convolution for every positive Gaussian width; these
Gaussians are an approximate identity, so the convolution tends
to $h$. For $h\ne0$ let $m$ be the smallest degree with nonzero
homogeneous Taylor polynomial $P_m$. Uniformly on the unit sphere,
$\widehat h(r\theta)=r^mP_m(\theta)+O(r^{m+1})$.
Consequently
\begin{equation}\label{eq:full-small-energy}
 \rho_h(\omega)=
 \frac{\int_{S^2}|P_m(\theta)|^2\,d\theta}
       {320\pi^5(7+2m)}\,\omega^{7+2m}
       +O(\omega^{8+2m}).
\end{equation}
Its leading coefficient is positive because a nonzero homogeneous
polynomial cannot vanish on the whole sphere. This proves the
support assertion. Repeated integration by parts in $h$ makes
$\widehat h$ decay faster than every inverse power. Thus
$\rho_h(\omega)$ tends as $\omega\to\infty$ to the finite
positive constant
$(320\pi^5)^{-1}\int |k|^4|\widehat h(k)|^2\,dk$.
The raw total mass is therefore infinite. Boundedness of that
density proves all the positive-time integrability claims.



We specify the common state space and map, so that convergence of
numbers is not substituted for a spectral representation. Let
$\mathfrak h$ be the complex Hilbert space of square-integrable
transverse vector functions on $\mathbb R^3$, with three colour
coordinates and Lebesgue momentum measure. Let $\mathcal F$ be
its symmetric Fock space, with vacuum $\Omega$, and let
$A_{\rm fr}=d\Gamma(|p|)$. This definition means that on the
$n$-particle space the energy multiplies by
$|p_1|+\cdots+|p_n|$; the domain consists of vectors for which
the corresponding squared weighted norm is summable over $n$.
The physical subspace used here consists of the vectors fixed by
every simultaneous adjoint $\SU(2)$ rotation of the three colours.
Orthogonality
of those rotations follows from the original trace metric
$-2\tr(T_\alpha T_\beta)=\delta_{\alpha\beta}$.

Choose any measurable real transverse polarization basis. Define
\begin{align}
 R_h^{rs}(p,q)&=-\frac{\widehat h(-p-q)}{(2\pi)^3}
 \left[\frac{(p\times e_r(p))\cdot(q\times e_s(q))}{\sqrt{|p||q|}}
       +\sqrt{|p||q|}\,e_r(p)\cdot e_s(q)\right],\nonumber\\
 J_t h&=\frac14\sum_{\alpha=1}^3\sum_{r,s=1}^2
 \int e^{-t(|p|+|q|)/2}R_h^{rs}(p,q)
    a^\dagger_{r\alpha}(p)a^\dagger_{s\alpha}(q)\Omega\,dp\,dq.
 \label{eq:full-state-map}
\end{align}
The integral is the symmetric two-particle Hilbert-space integral:
the convention with creation symbols has norm twice the squared
norm of its symmetric kernel. Equivalently its two-particle wave
function is $\sqrt2/4$ times that kernel with equal colours.
This gives an explicit definition without distributional creation
operators. The preceding integral proves that it is well-defined,
that it is a colour singlet, and that
\[
 \|J_t h\|^2=C_h(t),\quad
 \langle J_t h,A_{\rm fr}^nJ_t h\rangle
       =\int\omega^n e^{-t\omega}\,d\nu_h(\omega),\quad
 e^{-sA_{\rm fr}/2}J_t h=J_{t+s}h.
\]
The positive density proves injectivity on the real test functions;
the same kernel defines a complex-linear injective map on complex
test functions, with the Hermitian Fourier modulus in its norm.
It intertwines spatial translations: replacing $h(x)$ by $h(x-b)$
multiplies its creation kernel by $e^{i(p+q)\cdot b}$, precisely
the specified second-quantized translation on the two momenta.
Polarization changes are pointwise orthogonal changes of the
one-particle coordinates and induce unitary maps on their symmetric
tensor squares, preserving the state and its raw norm.

\subsection{One actual positive-coupling sequence for every local profile}

Return to the original Hamiltonian, taking
$L_j=j^2$, $a_j=1/(100j)$ and $\ell_j=a_j(2L_j+1)$, for $j\ge2$.
The following construction strengthens, rather than silently reuses,
the coupling selection of Section~15. Enumerate the positive rational
times as $t_1,t_2,\ldots$. Keep all the finite inequalities already
used there. For each of the finitely many original single-edge
weights $\mathbf1_e$ and each $n\le j$, require additionally
\begin{equation}\label{eq:full-diagonal-selection}
 \left|C_{g;\mathbf1_e,\mathbf1_{e'}}(t_n)
       -C_{0;\mathbf1_e,\mathbf1_{e'}}^{L_j,a_j}(t_n)\right|
       <\frac{\varepsilon_j}{|\mathsf E_{L_j}|^2}
 \quad(e,e'\in\mathsf E_{L_j}),\qquad
 \varepsilon_j=j^{-40}(1+\ell_j)^{-10}.
\end{equation}
Each centered single-edge vector includes its adjacent quarter-weight
faces. At fixed $j$ its convergence was established above; mixed
terms follow by the real polarization identity. Hence the finite
list holds for every sufficiently small positive $g$. Define $n_j$
as the smallest positive integer for which $g_j=2^{-n_j}<1/j$
satisfies this list and all the earlier finite lists. The well-ordering
of the positive integers and fixed-box convergence prove existence.
This specifies actual positive couplings in the original model;
it does not replace $H_{g_j}$ by its quadratic part.

Linearity in the weights bounds the covariance error at each listed
time by $\varepsilon_j\|h\|_\infty\|k\|_\infty$. Theorem
16.1 therefore gives convergence for
every real compact smooth profile and every positive rational time
on this one sequence. For an arbitrary positive real time, squeeze
the diagonal covariance between two rational times: its spectral
measure is nonnegative, so it decreases with time. The limiting
function (158) is continuous on $(0,\infty)$
by dominated convergence. Taking rational bounds to the specified
time proves
\begin{equation}\label{eq:full-actual-convergence}
 C_{g_j;h,h}(t)\longrightarrow C_h(t),\qquad t>0.
\end{equation}
Real polarization proves mixed convergence for all times.

For precision this also proves weak convergence of the actual
damped raw spectral measures, not just of a few moments. Fix $t>0$
and let $\mu_j(d\omega)=e^{-t\omega}d\nu_{h,g_j}(\omega)$.
Their masses converge to $C_h(t)$, and their Laplace transforms
at $s>0$ converge to $C_h(t+s)$. The inequality
\[
 (1-e^{-sR})\mu_j([R,\infty))
       \le C_{g_j;h,h}(t)-C_{g_j;h,h}(t+s)
\]
first with small $s$ and then large $R$ proves tightness, by
continuity of $C_h$; finitely many initial measures can be covered
separately. On a compact interval, the algebra generated by
$e^{-s\omega}$ separates points and contains constants, so its
polynomials uniformly approximate continuous functions there.
An explicit version follows by setting $z=e^{-\omega}$ and using
Bernstein polynomials for the resulting continuous function on
$[e^{-R},1]$. Thus the transforms and tightness identify the weak
limit uniquely as $e^{-t\omega}d\nu_h$. Applying this assertion
at time $t/2$ to the bounded continuous function
$\omega^n e^{-t\omega/2}$ proves convergence of every damped
moment. In particular the original vectors satisfy
\begin{equation}\label{eq:full-actual-norm-energy}
 \|y_{h,g_j}(t)\|^2\longrightarrow C_h(t),\qquad
 \langle y_{h,g_j}(t),A_{g_j}y_{h,g_j}(t)\rangle
       \longrightarrow -C_h'(t).
\end{equation}
They lie in every excitation-power domain since the semigroup
multiplier damps all powers. For any finite list of profiles and
times, their Gram matrices and their semigroup matrix elements
converge to those of the vectors $J_t h$ defined in
(160). Indeed the mixed time argument is
$(t+s)/2$ plus the intervening positive evolution time, and the
mixed covariance follows by polarization. This constructs the
common spectral correspondence at the level actually proved.
No norm-preserving identification of the entire changing Hilbert
spaces, or convergence of products of several local operators,
has been asserted.

\subsection{Low-energy tail, original raw amplitudes, and curvature concentration}

Put $H_h=\int h$. For $H_h\ne0$, equation
(159) and the substitution $v=t\omega$
give
\begin{equation}\label{eq:full-large-time}
 \rho_h(\omega)\sim\frac{H_h^2}{560\pi^4}\omega^7,
 \quad C_h(t)\sim\frac{9H_h^2}{\pi^4t^8},\quad
 -C_h'(t)\sim\frac{72H_h^2}{\pi^4t^9},\quad
 \frac{\langle J_th,A_{\rm fr}J_th\rangle}{\|J_th\|^2}
       \sim\frac8t.
\end{equation}
To justify the substitution over the entire integration range,
split at a fixed small frequency. On its complement the positive
time exponential is exponentially small, while on the first piece
the uniform Taylor remainder above is bounded by the next power;
the integrals are the elementary gamma integrals
$\int_0^\infty v^n e^{-v}dv=n!$, obtained by repeated integration
by parts. This also proves that for general nonzero $h$ with
first degree $m$, the quotient is $(8+2m)/t+o(t^{-1})$.
These are raw amplitudes: the norm tends to zero as time grows.
They are nevertheless nonzero at every finite time, with a ratio
tending to zero. The positive-coupling convergence in
(163) transfers both raw quantities at
each fixed time, and a further explicitly selected subsequence
transfers this time sequence: take $t=n$ and the least increasing
$j_n$ for which both errors are less than $1/n$ times their
positive limiting quantities. Existence follows from that equation.
Thus actual vectors $y_{h,g_{j_n}}(n)$ also have the quotient
$8/n+o(n^{-1})$ when $H_h\ne0$, with their unaltered norms and
energies asymptotic to the displayed limits.

The map has a useful exact continuity estimate for concentration.
For $h\in C_c^\infty$, $|\widehat h|\le\|h\|_1$ in
(158), so
\begin{equation}\label{eq:full-L1-bound}
 \|J_t h\|^2\le \frac{9}{\pi^4t^8}\|h\|_1^2.
\end{equation}
The constant is the integral
$4\pi\int_0^\infty K^6e^{-tK}dK/(320\pi^5t)$, with every
Fourier factor retained. The same estimate for differences extends
$J_t$ uniquely to $L^1(\mathbb R^3)$, by approximation with
compact smooth functions. More generally, a finite signed measure
defines the kernel (160) by its bounded Fourier
transform; the identical estimate with total variation proves
existence. Point measures have the exact raw norm given by equality
in (165).

For an explicit concentration family let $\rho$ be compact smooth
and real, $x_*$ fixed, and
$h_\epsilon(x)=A\epsilon^{-\beta}\rho((x-x_*)/\epsilon)$,
where $A,\beta\in\mathbb R$, $\epsilon>0$, and all these parameters
are retained. Direct change of variables gives
\begin{align}
 \widehat h_\epsilon(k)
   &=A\epsilon^{3-\beta}e^{-ik\cdot x_*}\widehat\rho(\epsilon k),
       \nonumber\\
 \|J_t h_\epsilon\|^2
   &= \frac{A^2\epsilon^{6-2\beta}}{320\pi^5t}
       \int |k|^4e^{-t|k|}|\widehat\rho(\epsilon k)|^2\,dk,
       \label{eq:full-concentration}\\
 \epsilon^{2\beta-6}\|J_t h_\epsilon\|^2
   &\longrightarrow \frac{9A^2(\int\rho)^2}{\pi^4t^8}.
       \nonumber
\end{align}
Dominated convergence uses $|\widehat\rho|\le\|\rho\|_1$.
It proves convergence of the correspondingly scaled kernels in
the same Fock space, including their energy forms, to the point
measure kernel. This is a displayed coordinate and amplitude map,
not a deletion of $A$ or a division by the vector norm. For
$\int\rho\ne0$, the limiting energy quotient at fixed $t$ is
$8/t$ as $\epsilon\downarrow0$, while the raw norm vanishes,
stays finite, or diverges according to the exponent $6-2\beta$.

Finally insert an actual regular smooth compact incompressible
Navier--Stokes profile from the preceding sections:
\begin{equation}\label{eq:full-NS-chain}
 u(s)\longmapsto h_s=\lambda^2|\nabla\times u(s)|^2
 \longmapsto D_{h_s}\longmapsto\Xi_{h_s,g_j}
 \longmapsto e^{-tA_{g_j}/2}\Xi_{h_s,g_j}
 \longleftrightarrow J_t h_s .
\end{equation}
Every map on the left is an exact coordinate operation on the
original objects; the last correspondence is the proved common
Gram and spectral limit. For $\lambda\ne0$ and a nonzero profile,
$H_{h_s}=\lambda^2\|\nabla\times u(s)\|_2^2>0$, using the
compact-support integration identity proved in Section~15. Hence
(164) applies with its exact squared
coefficient $\lambda^4\|\nabla\times u(s)\|_2^4$, and
(165) bounds its raw norm with the same
coefficient. To record the spatial scaling explicitly, for any
fixed nonzero compact smooth divergence-free $U$, put
$u_\epsilon(x)=\epsilon^{-\alpha}U((x-x_*)/\epsilon)$.
Its curvature weight has $A=\lambda^2$,
$\beta=2\alpha+2$, $\rho=|\nabla\times U|^2$, so
\[
 \|u_\epsilon\|_2^2=\epsilon^{3-2\alpha}\|U\|_2^2,
 \quad \|\nabla\times u_\epsilon\|_2^2
          =\epsilon^{1-2\alpha}\|\nabla\times U\|_2^2,
 \quad \|J_t h_\epsilon\|^2
       \sim \frac{9\lambda^4\|\nabla\times U\|_2^4}{\pi^4t^8}
                    \epsilon^{2-4\alpha}.
\]
This is an explicit family of regular profiles; it is not asserted
to be the singular profile of the supplied flow. The original
spacetime connection and current remain
$\mathcal A_i=\lambda u_iT$, $\mathcal A_0=0$ and
$j_i=(\lambda/g^2)(\Delta u_i-c^{-2}\partial_s^2u_i)T$.
No identification of the fluid time $s$ with the quantum time $t$
has been made. A singular endpoint can be passed through this
map only by evaluating its actual Fourier-weighted curvature
measures; a pointwise infinity does not specify those measures.

The full spectral correspondence proved here is stronger than the
lowest-band correspondence: its limiting measure has nonzero
weight at every positive energy scale and a continuous tail at
zero. Its Hamiltonian is exactly the explicitly defined
$d\Gamma(|p|)$, and its states are the equal-colour transverse
two-creation kernels above. Convergence of higher local products
and survival of a nonabelian interaction on a continuum scaling
trajectory have not been established by this calculation. Those
are further calculations in the original counterexample programme;
the free spectral correspondence itself is not a disproof of the
interacting Yang--Mills mass-gap assertion.

\section{A nonzero nonabelian vertex and the actual vacuum response}
\label{sec:nonabelian-vertex}

The free spectral correspondence in Section~16 does not remove the
nonabelian terms of the original operator. We calculate those terms
in the same exact tree and Haar coordinates, including the kinetic
contribution, and exhibit a nonzero physical three-creation channel.
Every statement in this section fixes the original finite box and
$a>0$ before the small-positive-coupling limit. Bounds depending on
that box are not asserted to be uniform in the continuum limit.

\subsection{Exact logarithmic fields and their retained density terms}

Use the tree and chords of Section~14. To avoid confusing a coefficient
matrix with the chart $\mathcal B_g$, write $t_{ce}$ for the original
additive tree matrix and $b_{ce}$ for the following second integer
matrix. If $e$ is a chord, set
$l_{ce}=\delta_{ce}$ and $r_{ce}=0$. If $e$ is a tree edge, set
$l_{ce}=\epsilon(s(c),e)$ and $r_{ce}=-\epsilon(t(c),e)$, where
$\epsilon(v,e)$ records whether the root-to-$v$ tree path uses $e$.
Then
\begin{equation}\label{eq:vertex-integer-matrices}
 t_{ce}=l_{ce}+r_{ce},\qquad b_{ce}=r_{ce}-l_{ce},\qquad
 \mathcal X_{e,\alpha}=
       \sum_c(l_{ce}L_{c,\alpha}+r_{ce}R_{c,\alpha}).
\end{equation}
This is exactly the original field obtained by varying that edge,
recomputing the rooted chord products and differentiating. In particular
each downstream tree holonomy has derivative $T_\alpha$, and the
inverse target holonomy contributes the displayed minus sign. Product
Haar measure makes each left and right field divergence-free.

Identify $y^\alpha T_\alpha$ with the real colour vector $y$.
The Pauli multiplication law is
\begin{equation}\label{eq:vertex-pauli}
 (u\cdot T)(v\cdot T)
       =-\tfrac14(u\cdot v)I+\tfrac12(u\times v)\cdot T,
 \quad [u\cdot T,v\cdot T]=(u\times v)\cdot T.
\end{equation}
Write $\operatorname{ad}_y v=y\times v$. Differentiation of the
exponential power series, grouping the positions of the differentiated
factor, gives
\[
 (d\exp)_Y(V)e^{-Y}=\int_0^1 e^{s\operatorname{ad}_Y}V\,ds.
\]
Thus the exact logarithmic coordinate derivatives of left and right
multiplication are respectively
\[
 F_-(\operatorname{ad}_y)T_\alpha,
 \quad F_+(\operatorname{ad}_y)T_\alpha,
 \qquad F_-(z)=\frac{z}{e^z-1},\quad F_+(z)=F_-(z)+z.
\]
The values at zero are their removable continuations. On a sufficiently
small original chart these are analytic matrix functions. Multiplying
their series by the exponential series proves
$F_-(z)=1-z/2+z^2/12+O(z^4)$ and
$F_+(z)=1+z/2+z^2/12+O(z^4)$, including the zero cubic coefficient.
The exact functions, rather than a truncated replacement, define all
remainders below.

Put $d=3r$ and keep the exact chart (112), with
$Z_c=\exp(gx_c\cdot T)$. On compactly supported smooth functions of
$x$, whose support lies inside that chart for small $g$, define
\[
 Y_{e,\alpha}(g)=g\mathcal B_g\mathcal X_{e,\alpha}\mathcal B_g^*.
\]
Let $D_{e,\alpha}(g)$ be the vector field whose $c$-component is
$l_{ce}F_-(g\operatorname{ad}_{x_c})e_\alpha+
r_{ce}F_+(g\operatorname{ad}_{x_c})e_\alpha$.
Density conjugation gives the exact equality
\begin{equation}\label{eq:vertex-exact-field}
 Y_{e,\alpha}(g)=D_{e,\alpha}(g)
                 -\tfrac12D_{e,\alpha}(g)\log\mathcal J(gx).
\end{equation}
Indeed apply the derivative to $\mathcal J(gx)^{-1/2}f(x)$
before multiplying by the chart density. Direct expansion of the
retained sine density gives
\[
 \log\mathcal J(gx)=\log(16\pi^2)^{-r}
       -\frac{g^2}{12}\sum_c|x_c|^2
       -\frac{g^4}{1440}\sum_c|x_c|^4+O(g^6|x|^6).
\]
For each $c$, all positive powers of $\operatorname{ad}_{x_c}$
are orthogonal to $x_c$. Consequently the radial contraction in
(170) retains exactly the $t_{ce}$ term.
The coefficients through order two are therefore
\begin{align}
 D^0_{e,\alpha}&=\sum_c t_{ce}\partial_{c,\alpha},\nonumber\\
 D^1_{e,\alpha}&=\tfrac12\sum_c b_{ce}
                   (x_c\times e_\alpha)\cdot\nabla_c,\nonumber\\
 D^2_{e,\alpha}&=\tfrac1{12}\sum_c t_{ce}
          [x_c\times(x_c\times e_\alpha)]\cdot\nabla_c,
 \qquad M^2_{e,\alpha}=\tfrac1{12}\sum_c t_{ce}x_c^\alpha,
       \label{eq:vertex-fields}\\
 Y_{e,\alpha}(g)&=D^0_{e,\alpha}+gD^1_{e,\alpha}
                   +g^2(D^2_{e,\alpha}+M^2_{e,\alpha})+O(g^4).
       \nonumber
\end{align}
Here $M^2$ is multiplication. Its sign can also be verified without
the density series: the divergences of $D^0,D^1$ are zero and
$\operatorname{div}D^2=\sum_ct_{ce}x_c^\alpha/6$.
Thus $D^2+M^2$ is skew-adjoint on flat compact test functions,
as required by the exact Haar isometry. This check retains the
term that would be missed by treating the chart density as constant.

\subsection{Every face coefficient and the full first two vertices}

In the original ordered four-link word of a face, let
$\zeta_1,\ldots,\zeta_4$ be the signed chord vectors, with zero
for each tree link. Reversed traversal supplies the minus sign in
that factor, without reversing the order of the four factors.
For $k\ge1$ the exact Taylor coefficient of the full potential is
\begin{equation}\label{eq:vertex-Wk}
 W_{p,k}(x)=-\sum_{n_1+\cdots+n_4=k}
       \tr\!\left(\prod_{m=1}^4
                   \frac{(\zeta_m\cdot T)^{n_m}}{n_m!}\right),
 \qquad 2-\tr\prod_{m=1}^4 e^{g\zeta_m\cdot T}
           =\sum_{k\ge1}g^kW_{p,k}.
\end{equation}
The constant term is zero because the original scalar $2$ cancels
the trace of the identity at this configuration, not because the
scalar Wilson term was removed from $H$.

The first coefficients can be written as the following complete
recursion. Start with $A_0=B_0=C_0=0$ and set
\begin{align}
 A_m&=A_{m-1}+\zeta_m,\nonumber\\
 B_m&=B_{m-1}+\tfrac12 A_{m-1}\times\zeta_m,\nonumber\\
 C_m&=C_{m-1}+\tfrac12 B_{m-1}\times\zeta_m
   +\tfrac1{12}\{A_{m-1}\times(A_{m-1}\times\zeta_m)
          +\zeta_m\times(\zeta_m\times A_{m-1})\}.
       \label{eq:vertex-BCH}
\end{align}
To verify the recursion, multiply the two exponential series through
degree three using (169), and then use
$\log(I+Q)=Q-Q^2/2+Q^3/3+O(Q^4)$. The linear coefficient is the
sum, the quadratic cross coefficient is half the commutator, and
the cubic coefficients are respectively one half and one twelfth
as displayed. This finite multiplication proves that the logarithm
of the face word is $(gA_4+g^2B_4+g^3C_4)\cdot T+O(g^4)$.
Since $2-2\cos(|z|/2)=|z|^2/4-|z|^4/192+O(|z|^6)$, substitution
proves, with $A=A_4,B=B_4,C=C_4$,
\begin{equation}\label{eq:vertex-W234}
 W_{p,1}=0,\quad W_{p,2}=\tfrac14|A|^2,\quad
 W_{p,3}=\tfrac12 A\cdot B,\quad
 W_{p,4}=\tfrac14|B|^2+\tfrac12 A\cdot C-\tfrac1{192}|A|^4.
\end{equation}
In particular the nested-commutator and negative fourth-order
trace contributions remain alongside the positive square.

Write $\{S,T\}=ST+TS$. The original operator in the exact chart is
\begin{align}
 \widetilde H_g&=-\frac2a\sum_{e,\alpha}Y_{e,\alpha}(g)^2
       +\frac1{2g^2a}\sum_p
                      \left(2-\tr\prod_{m=1}^4e^{g\zeta_m\cdot T}\right),
       \label{eq:vertex-full-H}\\
 H^{(1)}&=-\frac2a\sum_{e,\alpha}\{D^0_{e,\alpha},D^1_{e,\alpha}\}
                           +\frac1{2a}\sum_pW_{p,3},\nonumber\\
 H^{(2)}&=-\frac2a\sum_{e,\alpha}
       \left((D^1_{e,\alpha})^2+
             \{D^0_{e,\alpha},D^2_{e,\alpha}+M^2_{e,\alpha}\}\right)
                           +\frac1{2a}\sum_pW_{p,4}.
       \label{eq:vertex-H12}
\end{align}
Its zeroth coefficient is exactly $H_{\rm osc}$ of Section~14.
Define the remainder to be the difference of
(175) and $H_{\rm osc}+gH^{(1)}+g^2H^{(2)}$;
this definition retains every higher term.

We record the remainder control used below. Fix a small closed original
ball $|y|\le\rho$ inside the logarithmic chart, with a slightly larger
ball still inside it. Analyticity bounds all derivatives through any
fixed finite order on that larger ball. Taylor's integral remainder
for each coefficient of $Y$, and its first $x$ derivative, is bounded
by a constant times $g^3(1+|x|)^m$ on $|gx|\le\rho$, for some
fixed integer $m$; the zero cubic coefficient of $Y$ gives a stronger
bound where needed. For the potential, fifth-order Taylor remainder
before division by $g^2$ is bounded by $C g^5|x|^5$, so its operator
remainder is bounded by $C g^3|x|^5/a$. These constants are finite
sums of suprema of derivatives of the explicit functions above and
depend on the fixed box, $a,\rho$. They are not uniform-limit claims.

Let $\chi$ be a smooth colour-invariant cutoff equal to one on
$|y|\le\rho/2$ and zero before $|y|=\rho$. For any fixed polynomial
times the comparison Gaussian $F$, the exact chart vector
$\chi(gx)F(x)$, extended by zero, is smooth on the original compact
manifold after applying $\mathcal B_g^*$. The preceding coefficient
bounds and all Gaussian polynomial moments give
\begin{equation}\label{eq:vertex-remainder-bound}
 \|[\widetilde H_g-H_{\rm osc}-gH^{(1)}-g^2H^{(2)}]
                       \chi(gx)F\|\le C_F g^3.
\end{equation}
Derivatives of the cutoff are supported at $|x|\ge\rho/(2g)$.
Their polynomial factors are bounded by powers of $g^{-1}$ and
their Gaussian integrals by $C_Ng^N$ for every fixed $N$: bound
the smallest Gaussian quadratic eigenvalue below and integrate
$r^m e^{-cr^2}$ beyond that radius. This proves the same estimate
when the cutoff is commuted past any displayed coefficient operator.

\subsection{A physical three-creation image of the complete cubic term}

Let $V$ denote the matrix of all original transverse edge eigenvectors
and $\Sigma=\operatorname{diag}\sigma_\nu$. Let $\iota$ insert chord
coordinates with zero tree entries. The comparison Gaussian in $x$
has the matrix
\begin{equation}\label{eq:vertex-K}
 K=\iota^{\mathsf T}V\Sigma V^{\mathsf T}\iota
   =G^{-1/2}O\Sigma O^{\mathsf T}G^{-1/2},\qquad
 \Phi_0(x)=N_0\exp\left(-\tfrac18\sum_\alpha(x^\alpha)^{\mathsf T}Kx^\alpha\right).
\end{equation}
Here $N_0$ is the already specified unit-vacuum constant in this
coordinate measure. No trial state is changed by it. To prove the
matrix identity, use $T\iota=I$ and
$V=T^{\mathsf T}G^{-1/2}O$. Thus
$V^{\mathsf T}\iota=O^{\mathsf T}G^{-1/2}$, which proves
(178), positivity of $K$, and
$T^{\mathsf T}K=V\Sigma V^{\mathsf T}\iota$.

For each colour put $p_c=(Kx)_c$, as a three-component vector,
and set
\[
 v_e=\sum_c t_{ce}p_c,\qquad
 \theta_e=\sum_c b_{ce}x_c\times p_c.
\]
Direct differentiation of the full Gaussian gives
$D^0_{e,\alpha}\Phi_0=-v_e^\alpha\Phi_0/4$ and
$D^1_{e,\alpha}\Phi_0=\theta_e^\alpha\Phi_0/8$.
In $\sum_\alpha D^0_{e,\alpha}\theta_e^\alpha$ every derivative
contracts two equal indices of the alternating colour tensor, so
the sum is zero. In $\sum_\alpha D^1_{e,\alpha}v_e^\alpha$ the
derivative forces the component $(x_c\times e_\alpha)_\alpha$,
which is zero. The full first vertex therefore satisfies
\begin{equation}\label{eq:vertex-cubic-polynomial}
 H^{(1)}\Phi_0=P_3(x)\Phi_0,\qquad
 P_3(x)=\frac1{8a}\sum_e v_e\cdot\theta_e
                      +\frac1{4a}\sum_p A_p\cdot B_p.
\end{equation}
This calculation includes the kinetic part of the original nonlinear
operator, rather than identifying the Wilson cubic term alone with
the quantum vertex. Every term is invariant under simultaneous
adjoint colour rotations and is odd under $x\mapsto-x$.

Expand this exact alternating polynomial as
\[
 P_3(x)=\sum_{i<j<k}c_{ijk}\,
                 x_i\cdot(x_j\times x_k).
\]
Repeated chord indices give zero, and collecting distinct indices with
their permutation signs determines each coefficient uniquely. Put
\[
 D=\sqrt2G^{1/2}O\Sigma^{-1/2},\qquad
 d_I=\sum_J c_J\det D_{J,I},\quad |I|=|J|=3,
 \quad E_I=\frac{\sigma_{i_1}+\sigma_{i_2}+\sigma_{i_3}}a.
\]
The coordinate identity $x_c^\alpha=\sum_\nu D_{c\nu}
(a_{\nu\alpha}+a^\dagger_{\nu\alpha})$ then gives exactly
\begin{equation}\label{eq:vertex-three-state}
 P_3\Phi_0=\sum_{i<j<k}d_{ijk}\sum_{\alpha\beta\gamma}
   \epsilon_{\alpha\beta\gamma}
   a^\dagger_{i\alpha}a^\dagger_{j\beta}a^\dagger_{k\gamma}\Phi_0.
\end{equation}
Each attempted one-creation contraction identifies two colours in
$\epsilon_{\alpha\beta\gamma}$ and vanishes. The remaining
coefficient is the alternating determinant of the three coordinate
columns, proving the formula. Distinct ordered triples have disjoint
mode occupations; the six nonzero colour assignments are orthogonal
and each has norm one. Consequently
\begin{align}
 \|P_3\Phi_0\|^2&=6\sum_I|d_I|^2
       =6c^*\bigwedge\nolimits^3(2K^{-1})c,\nonumber\\
 \langle P_3\Phi_0,A_0P_3\Phi_0\rangle&=6\sum_I E_I|d_I|^2,
 \qquad A_0=H_{\rm osc}-\mu_0 .\label{eq:vertex-raw-moments}
\end{align}
The second equality in the first line follows also by direct Gaussian
contraction: two ordered colour determinants have inner product six
times the determinant of their cross-covariance submatrix. Their
covariance is exactly $2K^{-1}$ by (178).
Thus the entire map preserves the original raw Gram factors and
the physical three-colour singlet constraint.

\subsection{A finite rational certificate that the full vertex is nonzero}

\paragraph{Theorem 17.1.}
On the original $L=2$ box, $H^{(1)}\Phi_0$ is a nonzero physical
three-creation vector. For every original $a>0$,
\begin{equation}\label{eq:vertex-lower-bounds}
 \|H^{(1)}\Phi_0\|^2>
              \frac{14641}{125000\sqrt3\,a^2},\qquad
 \|A_0^{-1}H^{(1)}\Phi_0\|^2>
              \frac{14641}{13500000\sqrt3}.
\end{equation}

\paragraph{Proof.}
Use the following three positively oriented original chords, in the
displayed order:
\[
 c_1=((-1,1,-2),2),\qquad c_2=((-1,2,-2),3),\qquad
 c_3=((-1,1,-1),2).
\]
Directions here are $1,2,3$ as in the original manuscript. Assign
$x_{c_1}=e_1,x_{c_2}=e_2,x_{c_3}=e_3$ and all other chords zero.
The face $((-1,1,-2);2,3)$ has first two factors with those positive
vectors, third factor with $-e_3$, and fourth vector zero. All other
faces involve fewer than these three nonzero colour vectors.
Equation~(174) gives the magnetic coefficient
$-1/(8a)$.

Let $K_I$ be the three-by-three submatrix of $K$ in this ordered
list and put $Q_I=(bV\Sigma V^{\mathsf T}\iota)_{I,I}$.
The complete kinetic coefficient multiplied by $a$ is
\begin{equation}\label{eq:vertex-certificate-contraction}
 \frac18\sum_{j=1}^3(Q_I)_{j,\cdot}
                   \cdot\bigl(e_j\times(K_I)_{j,\cdot}\bigr).
\end{equation}
This follows from $T^{\mathsf T}K=V\Sigma V^{\mathsf T}\iota$
in (179); all original tree paths in
$b$ are retained.

Here are short rational enclosures for every entry. Each entry lies
strictly between the corresponding integer below divided by $10^6$
and that integer plus one divided by $10^6$:
\begin{equation}\label{eq:vertex-certificate-matrices}
 10^6K_I:\quad
 \begin{pmatrix}
 1273974&432279&-317748\\
 432279&1273974&-386480\\
 -317748&-386480&1563504
 \end{pmatrix},\quad
 10^6Q_I:\quad
 \begin{pmatrix}
 -845265&-977768&417911\\
 473487&-1342407&1208367\\
 421454&840903&-1200597
 \end{pmatrix}.
\end{equation}
We specify the exact finite calculation producing these bounds, so
they are rational certificates rather than rounded numerical evidence.
For $N=5$, use every frequency triple in
$\{0,1,2,3,4\}^3$ with at least two positive entries. There are
$125-1-12=112$ such blocks. Set
$s_j=2\sin(\pi j/10)$ and $\sigma=\sqrt{\sum_i s_{j_i}^2}$.
For an original edge $e=(n,i)$ put
\[
 \phi_e(\boldsymbol j)=w_{j_i}(n_i)\prod_{d\ne i}v_{j_d}(n_d),
\]
with zero when $j_i=0$, and with the exact $v,w$ of Section~16.
The contribution of that block to the full edge spectral matrix is
\begin{equation}\label{eq:vertex-projector-sum}
 S_{e,e'}(\boldsymbol j)=\sigma\phi_e(\boldsymbol j)
 \phi_{e'}(\boldsymbol j)
       \left(\delta_{i i'}-\frac{s_{j_i}s_{j_{i'}}}{\sigma^2}\right).
\end{equation}
The two matrices in (184) are
respectively the indicated entries of $\sum_{\boldsymbol j}S$
and $b\sum_{\boldsymbol j}S$. This is the complete transverse
projection in each exact frequency block, so no choice or rounding
of polarization vectors occurs.

Every trigonometric value here is an integer multiple of $\pi/10$.
Start with
$\sin(\pi/10)=(\sqrt5-1)/4$ and
$\cos(\pi/10)=\sqrt{(5+\sqrt5)/8}$ and use the angle-addition
recurrence up to multiples $0,\ldots,19$. These radical identities
follow, for example, by $\cos5\theta=0$ and the larger positive
root $\cos^2\theta=(5+\sqrt5)/8$; the bound $0<\theta<\pi/6$
selects $\theta=\pi/10$. Also retain the factors $\sqrt{1/5}$
and $\sqrt{2/5}$ in the original vertex and edge functions.

The included program \texttt{certify\_nonabelian\_vertex.py}
evaluates precisely this finite sum using endpoints that are integers
divided by $10^{40}$. Addition and negation are exact. For multiplication
take the minimum and maximum of the four endpoint products and round
outward by integer floor and ceiling division. For a nonzero-sign
divisor apply the same rule to the four rational quotients. To enclose
a positive square root with denominator $M=10^{40}$, take the integer
floor of $\sqrt{lM}$ for the lower numerator and the integer ceiling
of $\sqrt{uM}$ for the upper numerator, for input numerators $l,u$.
The integer-square inequalities are checked directly. Monotonicity
of each operation proves inductively that every resulting interval
contains the exact radical value. Integer tree paths and the complete
112-block loops are specified in that source. The full interval
endpoints, short bounds above, and final inequality are recorded in
\texttt{NONABELIAN\_VERTEX\_CERTIFICATE.json}.

Inserting even the wider short intervals
(184) into
(183), and adding the exact
$-1/8$ magnetic term, proves by rational arithmetic that the
full coefficient $c$ obeys
\begin{equation}\label{eq:vertex-strict-coefficient}
 -\frac{243}{1000}<ac<-\frac{242}{1000}.
\end{equation}
Sorting the three chord indices changes only its sign and not its
absolute value in the alternating coefficient vector of
(181).

Finally all original frequencies are at most $\sqrt{12}$, since
each $s_j\le2$. Hence $0<K\le\sqrt{12}I$ by
(178), and $2K^{-1}\ge I/\sqrt3$.
Diagonalizing this positive matrix proves
$\bigwedge^3(2K^{-1})\ge I/(3\sqrt3)$: its eigenvalues are
the products of three of the original eigenvalues. Thus
(181) is at least
$2\sum_I|c_I|^2/\sqrt3$, which with
(186) gives the first bound in
(182). Each $E_I\le6\sqrt3/a$, so
the inverse excitation multiplies the squared coefficient by at
least $a^2/108$. This proves the second bound, with every physical
spacing factor retained.



\subsection{The actual vacuum, raw quasimode Gram factor, and an observable response}

The finite-box comparison spectrum has an isolated simple physical
ground state and the actual low eigenvalues and projections converge,
as proved in Section~14. We now use those facts with the controlled
remainder, rather than assuming a formal perturbation series for the
actual vacuum. Define
\begin{align}
 \phi_1&=-A_0^{-1}P_3\Phi_0,\nonumber\\
 E^{(2)}&=\langle\Phi_0,H^{(2)}\Phi_0\rangle
                  -6\sum_I\frac{|d_I|^2}{E_I},\nonumber\\
 \phi_2&=-A_0^{-1}(I-|\Phi_0\rangle\langle\Phi_0|)
                       (H^{(1)}\phi_1+H^{(2)}\Phi_0).
       \label{eq:vertex-quasimode-coefficients}
\end{align}
All inverses in this formula act on the physical vacuum complement
of the fixed oscillator. These are finite explicit Fock calculations:
the inputs are polynomial Gaussians, the coefficient operators are
polynomial differential operators, and dividing each nonvacuum Fock
coefficient by its sum of positive frequencies defines the inverse.
They retain all intermediate physical modes produced by those
polynomials. Oddness of $P_3$ gives the exact first energy coefficient
$\langle\Phi_0,H^{(1)}\Phi_0\rangle=0$.

Set $F_g=\chi(gx)(\Phi_0+g\phi_1+g^2\phi_2)$ without dividing
by its norm. Substitution of (187)
into the exact operator, using (177), gives
\begin{equation}\label{eq:vertex-quasimode-residual}
 \|[\widetilde H_g-\mu_0-g^2E^{(2)}]F_g\|=O(g^3),
 \qquad \|F_g\|^2=1+g^2\|\phi_1\|^2+O(g^3).
\end{equation}
For the residual, orders zero, one and two cancel by the displayed
equations; their remaining polynomial-Gaussian terms have finite norm,
and cutoff commutators have the Gaussian tail bound already proved.
For the Gram factor, both $\phi_1,\phi_2$ are orthogonal to the
vacuum by definition. Expanding the unaltered square and using the
same tail bound proves the second equation.

The spectral theorem applied to this nonzero quasimode places an
actual eigenvalue within $O(g^3)$ of $\mu_0+g^2E^{(2)}$: otherwise
the residual norm would exceed that spectral distance times
$\|F_g\|$. The fixed-box spectral separation identifies it as the
unique actual vacuum eigenvalue. The orthogonal remainder of $F_g$
from its actual vacuum line has norm $O(g^3)$, by the same spectral
theorem and the fixed positive separation from that line. Positivity
and the already proved vacuum convergence fix the positive overlap.
Its squared magnitude equals $\|F_g\|^2$ minus that orthogonal
remainder's squared norm. Therefore the overlap is $1+O(g^2)$,
and we obtain the actual statements
\begin{equation}\label{eq:vertex-actual-vacuum}
 \mathcal E_g=\mu_0+g^2E^{(2)}+O(g^3),\qquad
 \mathcal B_g\psi_g=\Phi_0+g\phi_1+O(g^2).
\end{equation}
The only unit vector here is the originally specified actual vacuum.
The raw trial vector and its nontrivial Gram factor in
(188) have been left unchanged.
In particular the second lower bound of
(182) proves a nonzero first derivative
of the actual vacuum in this exact coordinate frame.

There is also an explicit bounded physical observable that detects it.
Choose even radial cutoffs $\chi_R(x)$, equal to one on $|x|\le R$
and zero for $|x|\ge2R$, with $0\le\chi_R\le1$, and set
$O_R(x)=\chi_R(x)P_3(x)$. Its multiplication norm is finite at
each fixed $R$. Lift it by the original logarithmic coordinate
$x=y/g$ and set it zero outside the chart; for $2Rg<\rho$ this
defines a globally smooth gauge-invariant multiplier $O_{R,g}$.
Its prescribed dependence on $g$ is part of this exact observable
map. Equation~(189) gives
\begin{equation}\label{eq:vertex-observable-response}
 \langle\psi_g,O_{R,g}\psi_g\rangle
   =-2g\langle\chi_RP_3\Phi_0,A_0^{-1}P_3\Phi_0\rangle+O_R(g^2).
\end{equation}
The constant term is zero by parity. Multiplication is bounded, so
the norm error in (189) controls every
remaining term of this expectation. Gaussian domination gives
$\chi_RP_3\Phi_0\to P_3\Phi_0$ in norm as $R\to\infty$.
The limiting inner product is
$6\sum_I|d_I|^2/E_I>0$ by the proved nonzero vertex. Fix the least
integer $R$ for which the inner product exceeds half that positive
quantity. Its existence follows from that norm convergence. This
produces a fixed bounded observable prescription with a strictly
negative nonzero first response for the actual positive-coupling
vacuum. Neither a Gaussian replacement of that vacuum nor an
assumption of interaction survival in the continuum is used.

\subsection{Exact relation to the full nonlinear coupling derivative}

The retained source \emph{Local-energy spacetime kernel and the exact
nonlinear coupling response} proves differentiation on the original
fixed Hilbert space. With $K_{\rm el}=\sum_eE_e$ and
$W=\sum_p(2-W_p)$ its derivative is
\[
 V_g=\partial_gH_g=\frac{4g}{a}K_{\rm el}-\frac1{g^3a}W,
 \quad \psi'_g=-R_g(V_g-\mathcal E'_g)\psi_g,
 \quad R_g=(A_g|_{\psi_g^\perp})^{-1}P_{\psi_g^\perp}.
\]
At each fixed $g>0$, smoothness of the simple vacuum follows directly
from the resolvent identity on the common elliptic $H^2$ domain and
a contour around that eigenvalue. Differentiating the eigen-equation
and its original unit-norm constraint gives the displayed formula;
the reduced inverse is bounded at that fixed regulator. The complete
source, including all unbounded-operator domains and Duhamel terms,
is retained with this manuscript.

There is an exact morphism from this derivative to the vertices
calculated above. Differentiate the full density-bearing chart, at
fixed original $F$, on any compact subset inside its logarithmic
domain. The density terms in the two derivatives agree, giving
\begin{equation}\label{eq:vertex-chart-connection}
 (\partial_g\mathcal B_g)F=\frac1g\mathcal D\mathcal B_gF,
 \qquad \mathcal D=x\cdot\nabla+\frac d2,
 \qquad \mathcal D^*=-\mathcal D
\end{equation}
on compact test functions. Explicitly, differentiating $g^{d/2}$
gives $d/(2g)$, differentiating $\mathcal J(gx)^{1/2}$ gives
$x\cdot\nabla_y\log\mathcal J(gx)/2$, and differentiating
$F(gx)$ gives $x\cdot\nabla_yF(gx)$. Expanding
$g^{-1}x\cdot\nabla_x(\mathcal B_gF)$ reproduces the latter two
terms exactly. Integration by parts gives the stated skew-adjoint
identity. It follows that
\begin{align}
 \partial_g\widetilde H_g
   &=\mathcal B_g V_g\mathcal B_g^*
                         +\frac1g[\mathcal D,\widetilde H_g],
       \label{eq:vertex-response-connection}\\
 \partial_g(\mathcal B_g\psi_g)
   &=\frac1g\mathcal D(\mathcal B_g\psi_g)
                 -\mathcal B_gR_g(V_g-\mathcal E'_g)\psi_g.
       \nonumber
\end{align}
These equalities hold locally on the stated smooth coordinate
domains; no derivative of a discontinuous boundary extension is
taken. Products with the fixed interior cutoffs used above give
the corresponding global test-vector identities.
Thus $H^{(1)}$ is the regular first coefficient of the full
coordinate derivative in (192).
Omitting its commutator would confuse a changing coordinate frame
with the fixed-Hilbert-space derivative and leave false singular
terms. This explicit relation retains both derivatives and shows
how they enter the same original operator.

\subsection{The first nonlinear coefficient of the original local covariance}

The odd vacuum response above does not by itself give a nonzero
first-order response for an even energy covariance. We calculate that
covariance, retaining its full quadratic coefficient. For the fixed
real original edge weights $f$, let $D_f^{(n)}$ be the coefficients
of $\mathcal B_gD_{f,g}\mathcal B_g^*$. Their exact formulas are
(175)--(176), with each electric
edge term multiplied by $f_e$ and each face term by its original
quarter-boundary average $f_p$. Write $D_f^{(0)}=D_0$,
$D_f^{(1)}=D_1$, $D_f^{(2)}=D_2$ only in this subsection; these
symbols denote the weighted operators, not the individual fields
$D^n_{e,\alpha}$ above.

We first establish the order of the expansion on which this calculation
depends. The same exact coefficient functions give
\[
 Y^{(3)}_{e,\alpha}=0,\qquad
 Y^{(4)}_{e,\alpha}=
 -\frac1{720}\sum_ct_{ce}(\operatorname{ad}_{x_c})^4e_\alpha
                       \cdot\nabla_c
 +\frac1{720}\sum_ct_{ce}|x_c|^2x_c^\alpha .
\]
The fourth density coefficient follows by differentiating the retained
$-g^4|x_c|^4/1440$ term; its radial contraction with every positive
adjoint power is zero. Set $Y^{(2)}=D^2+M^2$. Then
\begin{align*}
 H^{(3)}&=-\frac2a\sum_{e,\alpha}\{Y^{(1)}_{e,\alpha},Y^{(2)}_{e,\alpha}\}
                             +\frac1{2a}\sum_pW_{p,5},\\
 H^{(4)}&=-\frac2a\sum_{e,\alpha}
       \{(Y^{(2)}_{e,\alpha})^2+
                   \{Y^{(0)}_{e,\alpha},Y^{(4)}_{e,\alpha}\}\}
                             +\frac1{2a}\sum_pW_{p,6}.
\end{align*}
The coefficients $W_{p,5},W_{p,6}$ are the full finite sums
(172); no fifth or sixth power in an oriented factor
is omitted. These formulas specify all coefficients needed for the
following fourth-order quasimode.

Put $\phi_0=\Phi_0$. For $1\le n\le4$, recursively set
\begin{align}
 E^{(n)}&=\left\langle\Phi_0,
                   \sum_{k=1}^nH^{(k)}\phi_{n-k}\right\rangle,\nonumber\\
 \phi_n&=-A_0^{-1}P_{\Phi_0^\perp}
             \sum_{k=1}^n(H^{(k)}-E^{(k)})\phi_{n-k}.
       \label{eq:vertex-fourth-quasimode}
\end{align}
These agree with the preceding coefficients at $n=1,2$.
Induction proves that all inputs are finite physical polynomial
Gaussians. It also proves their parity: conjugation by
$\mathcal P F(x)=F(-x)$ multiplies $H^{(k)}$ by $(-1)^k$,
as follows directly from the homogeneous coefficient degrees in
the fields and potential. Hence $\phi_n$ has parity $(-1)^n$
and $E^{(1)}=E^{(3)}=0$. Taylor's integral remainder through the
next order, with the same Gaussian cutoff estimates as
(177), proves that
$F_g^{[4]}=\chi(gx)\sum_{n=0}^4g^n\phi_n$ has residual $O(g^5)$
for its displayed eigenvalue polynomial. Its component off the actual
vacuum line has both $L^2$ norm and $\widetilde H_g$ graph norm
$O(g^5)$ by the fixed positive spectral separation and its residual.

This stronger graph estimate justifies applying the unbounded local
operator to the actual vacuum expansion. Indeed on the original
common smooth domain the commuting edge Casimirs give
\begin{equation}\label{eq:vertex-weighted-graph-bound}
 \|D_{f,g}u\|\le\|f\|_\infty
           \left(\|H_gu\|+\frac{4|\mathsf P_L|}{g^2a}\|u\|\right).
\end{equation}
To verify it, the joint edge-Casimir expansion gives
$\|\sum f_eE_eu\|\le\|f\|_\infty\|\sum E_eu\|$.
Use $\kappa\sum E_e=H_g-bW$,
$\|W\|\le4|\mathsf P_L|$, and
$\|W_f\|\le\|f\|_\infty4|\mathsf P_L|$ to obtain
the displayed sum, including both potential contributions.
Consequently the $O(g^5)$ vacuum-complement remainder becomes only
$O_f(g^3)$ after applying $D_{f,g}$. All polynomial terms themselves
have the stronger local coefficient estimates already proved.

The raw overlap squared of $F_g^{[4]}$ with the actual vacuum is
$1+g^2\|\phi_1\|^2+O(g^3)$. Comparing coefficients with the
original unit-vacuum constraint therefore fixes
\[
 \chi_2=\phi_2-\tfrac12\|\phi_1\|^2\Phi_0,
 \qquad \mathcal B_g\psi_g
       =\Phi_0+g\phi_1+g^2\chi_2+O(g^3).
\]
The expansion also holds after applying the weighted operator, with
the corresponding coefficients, by
(194). The term along $\Phi_0$
is thus explicitly retained; the raw quasimode was not normalized.

Define the following finite Fock vectors and scalars:
\begin{align}
 m_0&=\langle\Phi_0,D_0\Phi_0\rangle,\nonumber\\
 m_2&=\langle\Phi_0,D_2\Phi_0\rangle
       +2\operatorname{Re}\langle\phi_1,D_1\Phi_0\rangle
       +2\operatorname{Re}\langle\chi_2,D_0\Phi_0\rangle
       +\langle\phi_1,D_0\phi_1\rangle,\nonumber\\
 \xi_0&=(D_0-m_0)\Phi_0,\qquad
 \xi_1=D_1\Phi_0+(D_0-m_0)\phi_1,\nonumber\\
 \xi_2&=(D_2-m_2)\Phi_0+D_1\phi_1+(D_0-m_0)\chi_2.
       \label{eq:vertex-local-state-coefficients}
\end{align}
Parity gives $m_1=0$, $\xi_0,\xi_2$ even and $\xi_1$ odd.
The preceding graph argument proves the actual expansions
$m_{f,g}=m_0+g^2m_2+O_f(g^3)$ and
$\mathcal B_g\Xi_{f,g}=\xi_0+g\xi_1+g^2\xi_2+O_f(g^3)$
in the required norms. This includes the motion of the vacuum and
its scalar subtraction.

Put $S_0(t)=e^{-tA_0}$ and $A_2=H^{(2)}-E^{(2)}$. Define
\begin{align}
 S_1(t)&=-\int_0^t S_0(t-s)H^{(1)}S_0(s)\,ds,\nonumber\\
 S_2(t)&=-\int_0^t S_0(t-s)A_2S_0(s)\,ds\nonumber\\
 &\quad+\int_{0\le r\le s\le t}
 S_0(t-s)H^{(1)}S_0(s-r)H^{(1)}S_0(r)\,dr\,ds .
       \label{eq:vertex-semigroup-coefficients}
\end{align}
These are finite-dimensional integrals on each polynomial input and
the finite Fock spaces its indicated insertions generate. To prove
the expansion for the actual semigroup, let $U_0=S_0(t)F$,
$U_1=S_1(t)F$, $U_2=S_2(t)F$ for a fixed physical polynomial
Gaussian $F$. Differentiation of these finite integrals gives
$(\partial_t+A_0)U_0=0$,
$(\partial_t+A_0)U_1=-H^{(1)}U_0$ and
$(\partial_t+A_0)U_2=-H^{(1)}U_1-A_2U_0$.
Their coefficients stay bounded on each fixed compact time interval.
The cutoff sum $\chi(gx)(U_0+gU_1+g^2U_2)$ therefore has residual
$O(g^3)$ for the actual equation
$(\partial_t+\widetilde H_g-\mathcal E_g)U=0$, uniformly on
that interval. Variation of constants and the actual excitation
semigroup's contraction bound make its difference from the exact
evolution at most the time integral of this residual. The cutoff
tail is smaller than every power of $g$. This proves the actual
semigroup expansion on these vectors without a bounded-operator
perturbation assumption.

It follows that, at every fixed $L,a,f,t\ge0$, the original nonlinear
covariance has the expansion
\begin{align}
 C_{g;f}(t)&=\langle\Xi_{f,g},e^{-tA_g}\Xi_{f,g}\rangle
       =C_{0;f}(t)+g^2 C_{2;f}(t)+O_{L,a,f,t}(g^3),
       \label{eq:vertex-C2}\\
 C_{2;f}(t)&=2\operatorname{Re}\langle\xi_2,S_0(t)\xi_0\rangle
       +\langle\xi_1,S_0(t)\xi_1\rangle\nonumber\\
 &\quad+2\operatorname{Re}\langle\xi_1,S_1(t)\xi_0\rangle
       +\langle\xi_0,S_2(t)\xi_0\rangle .\nonumber
\end{align}
There is no linear term: its two external-vector terms pair an odd
vector with an even one, and its semigroup insertion is odd between
two even vectors. The displayed second coefficient retains the
quartic operator, both cubic insertions, every intermediate mode,
the actual second vacuum coefficient and the centering correction.

For an entirely algebraic evaluation, take the original occupation
basis with energies $\lambda_n=\sum_{\nu,\alpha}n_{\nu\alpha}
\sigma_\nu/a$. Its creation and annihilation matrix entries are
$\sqrt{n_{\nu\alpha}+1}$ and $\sqrt{n_{\nu\alpha}}$ with the
corresponding occupation changed by one. The polynomial formulas
above determine every matrix entry from those rules. The time
integrals are exactly
\begin{align}
 I_t(\alpha,\beta)&=
 \begin{cases}(e^{-t\alpha}-e^{-t\beta})/(\beta-\alpha),&\alpha\ne\beta,\\
 t e^{-t\alpha},&\alpha=\beta,
 \end{cases}\nonumber\\
 (S_1)_{nm}&=-H^{(1)}_{nm}I_t(\lambda_n,\lambda_m),\nonumber\\
 (S_2)_{nm}&=-(A_2)_{nm}I_t(\lambda_n,\lambda_m)
       +\sum_kH^{(1)}_{nk}H^{(1)}_{km}
                      J_t(\lambda_n,\lambda_k,\lambda_m),
       \label{eq:vertex-time-matrices}
\end{align}
where for three distinct arguments
$J_t(\alpha,\beta,\gamma)=
\sum_{z\in\{\alpha,\beta,\gamma\}}
e^{-tz}/\prod_{w\ne z}(w-z)$.
For exactly two equal arguments, symmetry gives
\[
 J_t(\alpha,\alpha,\gamma)=
 \frac{[t(\gamma-\alpha)-1]e^{-t\alpha}+e^{-t\gamma}}
                  {(\gamma-\alpha)^2},\qquad
 J_t(\alpha,\alpha,\alpha)=\tfrac12t^2e^{-t\alpha}.
\]
These identities follow by integrating the two exponentials on the
simplex in (196); the repeated
cases follow by direct integration of the resulting linear factor.
Only finitely many occupations enter (197), because
all of its input vectors and insertions are polynomials of bounded
degree. The formula does not discard other physical intermediate
states that are reached by those insertions.

As an exact check on every subtraction, take $f_e=1$ on the full
finite box. Then $D_n=H^{(n)}$, $D_0=H_{\rm osc}$,
$m_0=\mu_0$ and $m_2=E^{(2)}$. The equations defining
$\phi_1,\phi_2$ give $\xi_0=\xi_1=\xi_2=0$ individually.
Thus both sides of (197) vanish, as the exact
identity $\Xi_{1,g}=0$ requires. A formula omitting the vacuum
response or its raw Gram coefficient would not pass these
coefficient identities.

We have proved that the original nonlinear Hamiltonian has a nonzero
colour-singlet cubic vertex, with an actual vacuum and observable
response, even though Section~16's selected covariance limit admits
the displayed free representation. The exact finite-box response
and the full intermediate-state terms are now available for the
joint regulator calculation. Their survival at fixed physical
scales, higher local products, and an interacting continuum
mass-gap violation are not established by the fixed-box bounds
of this section. The original counterexample programme remains
unfinished.

\section{The interacting first band in the original tensor coordinates}
\label{sec:interacting-tensor-band}

We now relate the complete first interaction matrix to the original
Fabel tensor and to each weighted curvature state. The result is an
exact onto map at every sufficiently small positive coupling in each
fixed original box, together with the first interaction coefficient of
its raw norms, energies and spectral evolution. All estimates in this
section fix $L\ge2$ and $a>0$. Their volume dependence is retained.
The source \emph{The cubic color singlet and the first vacuum-energy
correction} is included in full as
\texttt{references/CUBIC\_SINGLET\_INPUT.md}; its complete inspected
revision and coefficients are used below.

\subsection{The original weighted graph estimate and the exact input dictionary}

To avoid confusing the Gaussian matrix $K$ of (178)
with an operator, put $K_{\mathrm{el}}=\sum_eE_e$ and
$W=\sum_p(2-\operatorname{tr}U_p)$. Each of the four distinct edge
variables of an original face is in the fundamental representation.
Directly, $\sum_\alpha T_\alpha^2=-3I/4$; differentiating its
trace twice in each edge therefore proves
\[
 K_{\mathrm{el}}\operatorname{tr}U_p=3\operatorname{tr}U_p,
 \qquad K_{\mathrm{el}}W=3W-6|\mathsf P_L|.
\]
The constant and the number of original faces have not been removed.
For complex smooth $u$, integration by parts in Haar probability
measure gives
\[
 \operatorname{Re}\langle K_{\mathrm{el}}u,Wu\rangle
 =\sum_{e,\alpha}\int W|X_{e,\alpha}u|^2
       +\tfrac12\int(K_{\mathrm{el}}W)|u|^2.
\]
For the mixed derivative use
$\operatorname{Re}(\overline u Xu)=X|u|^2/2$ and integrate its
$XW$ factor once more. Since the original $\kappa b=a^{-2}$,
we obtain the exact identity
\begin{align}
 \|H_gu\|^2={}&\|\kappa K_{\mathrm{el}}u\|^2+\|bWu\|^2
   +\frac2{a^2}\sum_{e,\alpha}\int W|X_{e,\alpha}u|^2\nonumber\\
 &+\frac1{a^2}\int(3W-6|\mathsf P_L|)|u|^2.
 \label{eq:band-exact-graph}
\end{align}
Each $2-\operatorname{tr}U_p\ge0$. Joint Peter--Weyl decomposition
of the commuting nonnegative edge Casimirs gives
$\|\sum f_eE_eu\|\le\|f\|_\infty\|K_{\mathrm{el}}u\|$.
The retained quarter-boundary average has $|f_p|\le\|f\|_\infty$,
so $|W_f|\le\|f\|_\infty W$ pointwise. The triangle inequality,
$(x+y)^2\le2(x^2+y^2)$ and (199) prove
\begin{equation}\label{eq:band-weighted-graph}
 \|D_{f,g}u\|\le\sqrt2\,\|f\|_\infty
 \left(\|H_gu\|^2+\frac{6|\mathsf P_L|}{a^2}\|u\|^2\right)^{1/2}.
\end{equation}
Smooth functions are a graph core for the original compact elliptic
operator, so approximation extends this to its full operator domain.
This proves a bound independent of $g$ at fixed $L,a$. In particular
the lower-order raw quasimode already controls each weighted vacuum
vector through order $g^2$ with an $O(g^3)$ error. The fourth-order
fields, face coefficients and quasimode in Section~17 remain valid
and are retained. Its earlier $g^{-2}$ estimate is a weaker valid
bound; (200) proves why that loss is unnecessary
for this operator. The factor $|\mathsf P_L|/a^2$ remains explicit.

Here is the coordinate map to the retained input. Its matrices,
distinguished by a subscript $\mathrm{in}$, satisfy
\begin{equation}\label{eq:band-input-dictionary}
 B_{\mathrm{in}}=K/4,\quad M_{\mathrm{in}}=2K^{-1},\quad
 \mathsf A_{\mathrm{in}}=G^{1/2}O,\quad
 t^{\mathrm{in}}_{ec}=t_{ce},\quad s^{\mathrm{in}}_{ec}=-b_{ce}.
\end{equation}
Substituting into the input Gaussian gives exactly
$\Phi_0=N_0\exp(-\sum_\alpha x^{\alpha\mathsf T}Kx^\alpha/8)$.
Its variable $p_{\mathrm{in}}$ is our $Kx/4$.
The kinetic determinant $-2\sum s^{\mathrm{in}}t^{\mathrm{in}}
\det(x,p_{\mathrm{in}},p_{\mathrm{in}})/a$ therefore equals
the kinetic term $\sum v_e\cdot\theta_e/(8a)$ in
(179), by the cyclic determinant identity.
The ordered face expansion is the same Pauli product, so its magnetic
term agrees too. If $\Theta_{ijk}$ is the input coefficient of
$\det(z_i,z_j,z_k)\Phi_0$ for $i<j<k$, its precise map to
our creation coefficient is
\begin{equation}\label{eq:band-theta-map}
 d_{ijk}=\sqrt{\frac8{\sigma_i\sigma_j\sigma_k}}\,\Theta_{ijk}.
\end{equation}
Indeed $z_i^\alpha=\sqrt{2/\sigma_i}(a_{i\alpha}+a^\dagger_{i\alpha})$,
and these three distinct modes have no vacuum contractions. Thus
the two raw norms are exactly
$6|d_{ijk}|^2=48|\Theta_{ijk}|^2/(\sigma_i\sigma_j\sigma_k)$;
division by the identical excitation
$(\sigma_i+\sigma_j+\sigma_k)/a$ proves agreement of the
vacuum-energy and inverse-vector formulas, with no hidden scale.

\subsection{The complete first interaction matrix in a raw frame}

Let $\sigma=\sigma_*$ be the lowest spatial frequency, let the three
lowest modes be numbered $1,2,3$ as in Section~14, and put
$\Delta=2\sigma/a$. In the fixed ordered six coordinates
$\mathcal I=(11,22,33,12,13,23)$ use the following comparison vectors:
\begin{align}
 S_{ii}&=\frac1{\sqrt6}\sum_\alpha
                  (a^\dagger_{i\alpha})^2\Phi_0
       =\frac{\sigma}{2\sqrt6}(|z_i|^2-6/\sigma)\Phi_0,\nonumber\\
 S_{ij}&=\frac1{\sqrt3}\sum_\alpha
           a^\dagger_{i\alpha}a^\dagger_{j\alpha}\Phi_0
       =\frac{\sigma}{2\sqrt3}(z_i\cdot z_j)\Phi_0\quad(i<j).
 \label{eq:band-six-vectors}
\end{align}
The six and three independent colour contractions prove that these
vectors are orthonormal. They are exactly the input's six vectors,
not a change of any original trial state. Section~14 proves their
span is the full first physical eigenspace. Denote its projection
by $P_*$ and write $A_0=H_{\mathrm{osc}}-\mu_0$.

For complete evaluation let $h_n\Phi_0$ be the product Hermite
occupation basis on all $3r$ coordinates, with
$\omega_n=\sum_{i,\alpha}n_{i\alpha}\sigma_i/a$.
The probabilists' polynomial is
\[
 \operatorname{He}_k(u)=k!\sum_{j=0}^{\lfloor k/2\rfloor}
 \frac{(-1)^j u^{k-2j}}{2^j j!(k-2j)!},\qquad
 h_n(z)=\prod_{i,\alpha}
 \frac{\operatorname{He}_{n_{i\alpha}}(\sqrt{\sigma_i/2}z_i^\alpha)}
      {\sqrt{n_{i\alpha}!}}.
\]
Define the explicitly evaluable coefficients
\begin{align}
 c_{I,n}&=\langle h_n\Phi_0,H^{(1)}S_I\rangle,\quad
 V_{IJ}=\langle S_I,H^{(2)}S_J\rangle,\nonumber\\
 \mathsf K_{IJ}&=V_{IJ}
   -\sum_{|n|=3\ \mathrm{or}\ 5}
                  \frac{\overline{c_{I,n}}c_{J,n}}{\omega_n-\Delta}
   -E^{(2)}\delta_{IJ}.
 \label{eq:band-full-K}
\end{align}
These symbols specify finite coefficients, not unevaluated arbitrary
operator data: substitute the full polynomial fields and face words
of (176), then the displayed Hermite polynomial and
$x=G^{1/2}Oz$. Each monomial expectation is zero in odd degree;
in degree $2k$ it is the sum over all $(2k-1)!!$ pairings, with
factor $2\delta_{ij}\delta_{\alpha\beta}/\sigma_i$ for each pair.
Differentiating the Gaussian exponential proves this prescription.
The retained input gives the expanded face and kinetic contractions
and the identical matrix in its equations (24)--(36).
Degree and parity allow only $1,3,5$ for $H^{(1)}S_I$;
the degree-one projection vanishes because it is a colour vector,
whereas the original vector and its projection are fixed by all
colour rotations. Hence all denominators actually used in
(204) are at least $\sigma/a>0$. Every reached
spatial mode is included. In particular outside-mode contributions
are not removed from this formula. The input's occupation-number
proof cancels the terms acting exclusively on outside vacuum modes
between the two absolute energies; both sides of that identity are
retained in its equations (42)--(44).

For use without altering the raw state norms, construct an explicit
frame as follows. Set
\[
 \eta_I=-\sum_{|n|=3\ \mathrm{or}\ 5}
              \frac{c_{I,n}}{\omega_n-\Delta}h_n\Phi_0,
 \qquad \mathsf K^{\mathrm{abs}}=\mathsf K+E^{(2)}I.
\]
Let $\zeta_I\perp\operatorname{Ran}P_*$ solve
\begin{equation}\label{eq:band-zeta}
 (A_0-\Delta)\zeta_I
 =-H^{(1)}\eta_I-H^{(2)}S_I+\sum_JS_J\mathsf K^{\mathrm{abs}}_{JI}.
\end{equation}
Its right side is perpendicular to that cluster by
(204); expand its finite Hermite polynomial and
divide every other coefficient by $\omega_n-\Delta$.
The physical eigenspace at $\Delta$ is exactly the displayed six,
so the division is defined on this invariant input. This also
specifies every coefficient of $\zeta_I$. No negative vacuum
denominator is changed or suppressed.

Let $P_g=1_{(\Delta/2,11\Delta/10)}(A_g)$ be the actual physical
projection of Section~14. Define the six columns
\begin{equation}\label{eq:band-raw-frame}
 \mathcal R_g e_I=P_g\mathcal B_g^*\chi(gx)
                         (S_I+g\eta_I+g^2\zeta_I),\qquad
 \mathcal G_g=\mathcal R_g^*\mathcal R_g.
\end{equation}
The cutoff coefficient equations give residual $O(g^3)$ for
$(\mu_0+\Delta)I+g^2\mathsf K^{\mathrm{abs}}$.
The same spectral-separation argument as Section~17, applied to
six columns, bounds their off-cluster part and its original graph
norm by $O(g^3)$. One can first orthogonally diagonalize the real
symmetric finite matrix in the residual; the distance from each
approximate eigenvalue to the other actual clusters is bounded
below at this fixed box. The off-cluster resolvent then gives
the asserted bound for every column and hence their operator norm.

Parity makes $\eta_I$ odd, whereas $S_I,\zeta_I$ are even.
The perpendicular choice in (205) therefore gives
\begin{align}
 \mathcal G_g&=I+g^2\mathcal G_2+O(g^3),\quad
             (\mathcal G_2)_{IJ}=\langle\eta_I,\eta_J\rangle,\nonumber\\
 \mathcal R_g^*A_g\mathcal R_g
       &=\Delta\mathcal G_g+g^2\mathsf K+O(g^3),\nonumber\\
 \mathcal A_g^{\mathrm{coord}}
       :=\mathcal G_g^{-1}\mathcal R_g^*A_g\mathcal R_g
       &=\Delta I+g^2\mathsf K+O(g^3).
 \label{eq:band-raw-Gram-energy}
\end{align}
For the energy identity multiply the residual by the bounded frame
adjoint and subtract the actual vacuum
$\mathcal E_g=\mu_0+g^2E^{(2)}+O(g^3)$.
Invertibility of $\mathcal G_g$ follows from its convergence to $I$.
The frame is onto the actual six-dimensional cluster. These formulas
use its inverse Gram matrix, not a unit-norm alteration of its
columns. The coordinate operator is self-adjoint for the retained
inner product $q^*\mathcal G_g r$, since its energy form is Hermitian.

\subsection{An exact onto tensor map at positive coupling}

Keep the full original weight $f_S(e)=m(e)^{\mathsf T}S m(e)$,
including all off-diagonal coefficients of the real symmetric tensor.
Complexification below is linear only; real tensors give real states.
Define
\begin{equation}\label{eq:band-tensor-map}
 \mathcal T_g:\operatorname{Sym}_3(\mathbb C)\longrightarrow
       \operatorname{Ran}P_g,\qquad
 \mathcal T_gS=P_g\Xi_{f_S,g}.
\end{equation}
The original vacuum subtraction is linear in the weight, so this
is an exact linear map at each $g$. Let $\mathcal T_0$ denote its
comparison limit. In the six-vector coordinates its full coefficient
$q(S)$, with $R=\mathsf R_S^{\min}$ from (119), is
\begin{equation}\label{eq:band-q}
 q_{ii}(S)=\frac{\sqrt6}{4a}R_{ii},\qquad
 q_{ij}(S)=\frac{\sqrt{12}}{4a}R_{ij}\quad(i<j),\qquad
 \|q(S)\|^2=\frac3{8a^2}\operatorname{tr}(R^*R).
\end{equation}
The diagonal contraction has norm squared six. An off-diagonal
matrix entry occurs twice in $\Phi(R)$ and its single colour
contraction has norm squared three. These facts prove every factor
in (209) and recover the unchanged mass
$d_{L,a}(S)$ of (120).

In the coordinates $(S_{11},S_{22},S_{33},S_{12},S_{13},S_{23})$,
write $F$ for the matrix of $q$. Its diagonal block is
$c_d(\boldsymbol 1\boldsymbol 1^{\mathsf T}-I_3)$ and its other
block is $c_oI_3$, where
\[
 c_d=\frac{\sqrt6\sigma\delta_L}{4a},\qquad
 c_o=\frac{\sqrt{12}\sigma\mathfrak m_L^2}{4a},\qquad
 \det F=2c_d^3c_o^3\ne0.
\]
The determinant follows from eigenvalues $2,-1,-1$ of the first
displayed integer matrix. More explicitly, given $q$, set
$R_{ii}=4a q_{ii}/\sqrt6$ and $R_{ij}=4a q_{ij}/\sqrt{12}$;
then its unique preimage is
\begin{equation}\label{eq:band-F-inverse}
 S_{ii}=\frac{\operatorname{tr}R/2-R_{ii}}{\sigma\delta_L},
 \qquad S_{ij}=\frac{R_{ij}}{\sigma\mathfrak m_L^2}\quad(i<j).
\end{equation}
Here $\delta_L<0$ and $\mathfrak m_L^2>0$ in the actual box;
neither sign has been changed by the coordinate map.

Put
$Q_g=\mathcal G_g^{-1}\mathcal R_g^*\mathcal T_g$.
Section~17's weighted vector expansion, now controlled also by
(200), and the column expansion above give
$Q_g=F+O(g^2)$. There is no linear term: both candidate linear
inner products pair one odd vector and one even vector. Thus
$\det Q_g\to\det F\ne0$. For each fixed $L,a$ it follows that
there is $g_0(L,a)>0$ such that $\mathcal T_g$ is a bijection
onto $\operatorname{Ran}P_g$ for $0<g<g_0(L,a)$.
Its exact inverse, independent of a basis choice in that actual
Hilbert space, is
\begin{equation}\label{eq:band-exact-inverse}
 \mathcal T_g^{-1}u=(\mathcal T_g^*\mathcal T_g)^{-1}
                              \mathcal T_g^*u,\qquad u\in\operatorname{Ran}P_g.
\end{equation}
For the adjoint use the displayed six tensor-entry coordinates with
their ordinary Euclidean inner product. Multiplying (211)
by $\mathcal T_g$ gives the identity on its range, and multiplying
it by $\mathcal T_g$ on the right gives the identity on all six
tensor coordinates. This proves both directions of the exact morphism.

In particular let $h$ be any real original edge profile and retain its
original averaged face weights. Define
\begin{equation}\label{eq:band-h-tensor}
 S_g(h)=(\mathcal T_g^*\mathcal T_g)^{-1}\mathcal T_g^*
                                   P_g\Xi_{h,g}.
\end{equation}
Then $\mathcal T_gS_g(h)=P_g\Xi_{h,g}$ exactly. The kernel of this
map consists exactly of profiles whose actual projected state is zero;
it includes every constant profile because $\Xi_{1,g}=0$.
It does not assert injectivity on the whole space of profiles.
At comparison coupling, (210) gives the same map
with $R=\mathsf R_h^{\min}$ computed using all original weighted
edges and faces. Also $S_g(h)=S_0(h)+O(g^2)$ for each fixed profile.
For a regular Navier--Stokes field at fluid time $s$, insert the
actual sampled profile $h(e)=\lambda^2|\nabla\times u(s,a m(e))|^2$.
This supplies an exact coordinate map from that finite-box curvature
state to its tensor representative. The curvature map's units,
sampling, physical spacing and original fluid time remain in $h$;
the quantum projection does not identify that time with a semigroup time.

\subsection{Raw interaction energies and the resolved spectral splitting}

For fixed $S$ write $u_g=\mathcal T_gS$ and let its coordinates in
the unchanged frame be
$q_g=\mathcal G_g^{-1}\mathcal R_g^*u_g=q+g^2b_S+O(g^3)$,
where $q=FS$. In terms of Section~17's complete weighted-state
coefficients for $f=f_S$, its second coefficient is explicitly
\[
 (b_S)_I=\langle S_I,\xi_2\rangle+\langle\eta_I,\xi_1\rangle
       +\langle\zeta_I,\xi_0\rangle
       -\sum_J(\mathcal G_2)_{IJ}q_J.
\]
Indeed expand $\mathcal R_g^*\Xi_{f_S,g}$, then multiply by
$\mathcal G_g^{-1}=I-g^2\mathcal G_2+O(g^3)$.
The projector in the first factor makes this the same as
$\mathcal R_g^*P_g\Xi_{f_S,g}$ exactly. Its off-cluster column
error contributes $O(g^3)$ against the bounded fixed-box weighted
state norm. This retains every intermediate state and vacuum correction.
The exact norm is $q_g^*\mathcal G_gq_g$, so (207)
proves
\begin{align}
 \|u_g\|^2&=n_0+g^2n_2+O(g^3),\quad n_0=q^*q,\nonumber\\
 n_2&=2\operatorname{Re}(q^*b_S)+q^*\mathcal G_2q,\nonumber\\
 \langle u_g,A_gu_g\rangle
       &=\Delta(n_0+g^2n_2)+g^2q^*\mathsf Kq+O(g^3).
 \label{eq:band-raw-expansions}
\end{align}
For $S\ne0$ the proved inverse gives $n_0>0$. Dividing these two
unaltered scalar quantities yields
\begin{equation}\label{eq:band-tensor-Rayleigh}
 \frac{\langle\mathcal T_gS,A_g\mathcal T_gS\rangle}
      {\|\mathcal T_gS\|^2}
 =\Delta+g^2\frac{S^*F^*\mathsf KFS}{S^*F^*FS}+O(g^3).
\end{equation}
The $n_2$ terms cancel in this ratio only after appearing explicitly
in both quantities in (213); they were not
dropped from the state. The positive definite tensor Gram matrix is
$F^*F$. Therefore the generalized eigenvalues of
$(F^*\mathsf KF,F^*F)$ equal every eigenvalue of $\mathsf K$:
multiply the equation by $(F^*)^{-1}$ and put $q=FS$.
In fact the exact inverse (211) represents
every actual first-band eigenvector by a unique tensor. Minimizing
the exact quotient over $S\ne0$ thus gives the actual first physical
gap, with the complete coefficient (204).

This also resolves the splitting in a correlation limit. Fix
$\tau\ge0$ and retain the original quantum time $t=\tau/g^2$.
The exact frame intertwining is
$A_g\mathcal R_g=\mathcal R_g\mathcal A_g^{\mathrm{coord}}$;
it follows because the range is invariant and its inverse coordinate
map is $\mathcal G_g^{-1}\mathcal R_g^*$. Consequently for any
fixed tensors $S,T$,
\begin{equation}\label{eq:band-long-time}
 e^{\tau\Delta/g^2}
 \langle\mathcal T_gS,e^{-\tau A_g/g^2}\mathcal T_gT\rangle
 \longrightarrow (FS)^*e^{-\tau\mathsf K}(FT).
\end{equation}
To prove the limit, the left side equals
$q_g(S)^*\mathcal G_g
 \exp[-\tau(\mathcal A_g^{\mathrm{coord}}-\Delta I)/g^2]q_g(T)$.
The matrix in its exponential tends in norm to $\mathsf K$ by
(207). The matrix exponential converges
uniformly on each bounded $\tau$ interval: integrate the exact
finite-matrix Duhamel identity between the two exponentials and
bound it by $\tau e^{\tau(\|\mathsf K\|+1)}O(g)$.
Together with $\mathcal G_g\to I$ and $q_g\to F$, this proves
(215). Both states are exactly in the actual
band throughout this identity; no full-space approximation is
multiplied by the large prefactor. The large prefactor is displayed and
does not redefine the original energy origin or state norm. No
positivity of $\mathsf K$ is assumed in this matrix limit.

Finally the complete Fabel tensor of Section~13 satisfies
$z^6K_\tau\to S_*=\eta\eta^{\mathsf T}$,
$\eta=(1/4,3/8,-9/4)$, as $z=\sqrt{1-8\tau}\downarrow0$.
Here this $\tau$ is the original material parameter, distinct from
the semigroup parameter in (215). At fixed
$L,a,g$, linearity of (208) proves
\begin{align}
 z^6\mathcal T_gK_\tau&\longrightarrow\mathcal T_gS_*,\nonumber\\
 z^{12}\|\mathcal T_gK_\tau\|^2&\longrightarrow\|\mathcal T_gS_*\|^2>0,
 \quad
 z^{12}\langle\mathcal T_gK_\tau,A_g\mathcal T_gK_\tau\rangle
       \longrightarrow\langle\mathcal T_gS_*,A_g\mathcal T_gS_*\rangle.
 \label{eq:band-cusp-raw}
\end{align}
The last passage is valid since the actual band is finite-dimensional
and $A_g$ is bounded on it. Its limiting ratio has correction
$(FS_*)^*\mathsf K(FS_*)/\|FS_*\|^2$ by
(214). This is the exact way the given
material direction samples the full interaction matrix. Its sign
cannot be inferred from a scalar coordinate sign in the source map.
No simultaneous change of $L,a,g,z$ is substituted for these ordered
limits: controlling their combined remainders and the nonlinear
continuum states remains necessary for the requested mass-gap violation.

\subsection{All three interaction channels in the source tensor}

The retained input \emph{Exact cubic-box symmetry of the first
physical gap correction}, included in
\texttt{references/CUBIC\_BOX\_SYMMETRY\_INPUT.md}, now determines
the entire structure of (204). We give the exact
basis dictionary before applying it. The input's interval coordinate
is $l=n+L$. Substitution in its cosine and sine functions gives
precisely our $v_1(n),w_1(n)$ of Section~14, without a scale or sign
change. Its spatial modes are ordered $(V_{12},V_{13},V_{23})$.
Our normal-index modes satisfy
\[
 (V_1,V_2,V_3)=(V_{23},-V_{13},V_{12}),\qquad
 z_{\mathrm{in}}=Jz,\quad
 J=\begin{pmatrix}0&0&1\\0&-1&0\\1&0&0\end{pmatrix}.
\]
Both statements follow by comparing the two signed components of
(117). In the six coordinates, this sends the
input's diagonal states to $(S_{33},S_{22},S_{11})$ and its
off-diagonal states to $(-S_{23},S_{13},-S_{12})$.
It is an orthogonal signed permutation and retains every raw norm.

For completeness the source symmetry acts on the original links,
not only on these three modes. Extend $U_{-e}=U_e^{-1}$ and for a
signed coordinate permutation $R$ set $(\alpha_RU)_{Re}=U_e$.
This is a Haar-preserving permutation and inversion of the original
edge factors, with inverse $\alpha_{R^{-1}}$. Inversion preserves
the bi-invariant Casimir. It maps each face word to a cyclically
reordered face word or its inverse, whose SU(2) trace is identical.
The pullback therefore commutes with the complete $H_g$ and with
the full vertex gauge action, carried by $h_{Rn}=h_n$. It fixes the
actual positive vacuum and the first-band projector.

The chosen root can move under this operation. Its exact induced
chord map is obtained by applying $\alpha_R$ to the tree-identity
representative and recomputing every root-to-vertex product and
rooted chord; call it $\beta_R$. This finite product map and its
inverse act on the physical functions. Its derivative at the
identity is the original oriented cochain permutation modulo
gradients. Conjugation by $\mathcal B_g$ retains the Haar and
dilation density. On a compact $x$ set,
$g^{-1}\log\beta_R(\exp(gx))=D\beta_R(I)x+O(g|x|^2)$;
bounded analytic derivatives on an interior chart and Gaussian
tail estimates give convergence on each polynomial Gaussian.
The detailed inverse-chart bounds and all root changes are proved
in Sections~1--3 of the retained source.

On the three-mode basis the limiting action is $\bigwedge^2R$:
a coordinate permutation takes $V_{ij}$ to $V_{\pi(i)\pi(j)}$,
and a reflection with coordinate signs $\varepsilon_i$ multiplies
it by $\varepsilon_i\varepsilon_j$. The latter follows directly
from $v_1(-n)=-v_1(n)$ and the even edge function together with
its reflected orientation. Its action on the six states is the
symmetric square. In the raw frame (206) the
exact coordinate representation of the symmetry commutes with
$\mathcal A_g^{\mathrm{coord}}$ and tends to that symmetric square.
Subtract $\Delta I$ in (207), divide the
commutator by $g^2$, and pass to the finite-dimensional limit.
Thus $\mathsf K$ commutes with every one of these actions.
Reflections fix all three diagonal states but give the three
off-diagonal states distinct nontrivial sign characters. This
forces all cross-block entries and all distinct off-diagonal-state
entries to be zero. Permutations equate the three diagonal entries,
their three mutual entries, and the three off-diagonal entries.
The displayed signed change $J$ preserves the resulting matrix:
\begin{equation}\label{eq:band-three-channel-K}
 \mathsf K=\begin{pmatrix}
 d&o&o&0&0&0\\o&d&o&0&0&0\\o&o&d&0&0&0\\
 0&0&0&c&0&0\\0&0&0&0&c&0\\0&0&0&0&0&c
 \end{pmatrix}.
\end{equation}
Here $d,o,c$ are exactly its three full matrix entries in
(204), each with original units $a^{-1}$.
The coefficient $c$ in this display is not the conversion constant
in the fluid-to-spacetime coordinate map.

Define $k_s=d+2o$, $k_t=d-o$ and $k_o=c$. The three orthogonal
projections in six coordinates are respectively onto the constant
diagonal vector, the diagonal sum-zero plane, and the off-diagonal
three-space. Direct multiplication proves that their eigenvalues
are $k_s,k_t,k_o$. For any original tensor $S$, the unchanged squared
norms of its three components are
\begin{align}
 N_s(S)&=\frac{\sigma^2\delta_L^2}{2a^2}|\operatorname{tr}S|^2,\nonumber\\
 N_t(S)&=\frac{3\sigma^2\delta_L^2}{8a^2}
                 \sum_i|S_{ii}-\operatorname{tr}S/3|^2,\nonumber\\
 N_o(S)&=\frac{3\sigma^2\mathfrak m_L^4}{4a^2}
                 \sum_{i<j}|S_{ij}|^2.
 \label{eq:band-channel-masses}
\end{align}
Indeed the scalar projection of $q$ has mass
$|\sum_iq_{ii}|^2/3$, with $\sum_iq_{ii}=2c_d\operatorname{tr}S$;
the sum-zero projection has entries
$-c_d(S_{ii}-\operatorname{tr}S/3)$; the remaining entries
are $c_oS_{ij}$. This proves (218) and
$N_s+N_t+N_o=d_{L,a}(S)$, retaining the three positive raw masses.
Consequently the full coefficient in (214) is
\begin{equation}\label{eq:band-channel-correction}
 \frac{k_sN_s(S)+k_tN_t(S)+k_oN_o(S)}{N_s(S)+N_t(S)+N_o(S)}.
\end{equation}
Similarly the limit in (215) for equal tensors
is exactly $e^{-\tau k_s}N_s+e^{-\tau k_t}N_t+e^{-\tau k_o}N_o$.
The tensors $I_3$, $\operatorname{diag}(1,-1,0)$ and
$e_1e_2^{\mathsf T}+e_2e_1^{\mathsf T}$ isolate these three
coefficients respectively, with their displayed raw masses.

For the full specified Fabel endpoint $S_*=\eta\eta^{\mathsf T}$,
the diagonal entries are $(4,9,324)/64$ and the off-diagonal
entries are $(6,-36,-54)/64$. Substitution gives
\begin{equation}\label{eq:band-Fabel-channel-masses}
 N_s(S_*)=\frac{\sigma^2}{a^2}\frac{113569}{8192}\delta_L^2,
 \quad N_t(S_*)=\frac{\sigma^2}{a^2}\frac{100825}{16384}\delta_L^2,
 \quad N_o(S_*)=\frac{\sigma^2}{a^2}\frac{1593}{2048}\mathfrak m_L^4.
\end{equation}
For the second numerator the three traceless entries are
$(-325,-310,635)/192$; for the third, the squared off-diagonal
sum is $531/512$. These direct rational calculations prove every
coefficient in the display. The cusp direction therefore reaches
all three interacting channels with nonzero raw mass. Its full
interaction correction is their weighted value in
(219), not a sign assigned from one
source coordinate. Evaluating $d,o,c$ and their simultaneous-regulator
behavior remains part of the original spectral calculation.

\subsection{An explicit smooth fluid representative of every real band state}

The onto tensor map admits a further exact right inverse through
regular Navier--Stokes curvature profiles, with the forcing constructed
rather than suppressed. Fix the original positive $\nu,a,g$, with
$0<g<g_0(L,a)$, and fix real $\lambda\ne0$. Let $w$ be a real
vector in the actual first physical band; reality refers to complex
conjugation in the original group-coordinate Hilbert space. Since
the Hamiltonian, its positive vacuum and its spectral projector are
real, the inverse (211) gives a real tensor
$S=\mathcal T_g^{-1}w$. Put
\begin{equation}\label{eq:band-positive-edge-profile}
 M_S=\max_{e\in\mathsf E_L}|f_S(e)|,\qquad h_e=f_S(e)+M_S\ge0.
\end{equation}
This addition retains its original units, since $M_S$ has the same
units as the edge weight. It changes neither centered state nor
projection: $D_{h,g}=D_{f_S,g}+M_SH_g$, so
$\Xi_{h,g}=\Xi_{f_S,g}$ exactly on the original actual vacuum.
The corresponding face weights still use the original quarter-edge
average, which adds the same scalar $M_S$ to every face.

Let $p_e=a m(e)$ be the distinct physical edge midpoints and set
$\rho=a/8$. The distance between two distinct such points is at
least $a/\sqrt2$: their coordinates are integers or half-integers
times $a$, with exactly one half-integer coordinate. If the
half-integer coordinates have the same position, a nonzero
difference has magnitude at least $a$; if their positions differ,
two components have magnitude at least $a/2$. This proves that
the closed radius-$\rho$ balls are pairwise disjoint.

For an explicit cutoff define $b(t)=e^{-1/t}$ for $t>0$ and
$b(t)=0$ for $t\le0$, and
\[
 \chi_e(x)=\frac{b(1-|x-p_e|^2/\rho^2)}
 {b(1-|x-p_e|^2/\rho^2)+b(|x-p_e|^2/\rho^2-1/4)}.
\]
The denominator is positive everywhere. The vanishing of all
one-sided derivatives of $b$ at zero proves smoothness;
$\chi_e=1$ for $|x-p_e|\le\rho/2$ and zero for
$|x-p_e|\ge\rho$. Choose a fixed unit spatial vector $e_3$ and put
\begin{equation}\label{eq:band-fluid-right-inverse}
 \omega_e=\frac{\sqrt{h_e}}{|\lambda|}e_3,\qquad
 \mathcal V_e(x)=-\tfrac14\chi_e(x)|x-p_e|^2\omega_e,
 \qquad U_S(x)=\nabla\times\sum_e\mathcal V_e(x).
\end{equation}
All quantities in this finite sum are explicit original-coordinate
functions. Curl of a smooth compactly supported vector field is
divergence-free, so $U_S\in C_c^\infty(\mathbb R^3;\mathbb R^3)$
and $\nabla\cdot U_S=0$. Near $p_e$, disjointness removes all other
summands and the cutoff equals one. There
$U_S(x)=\omega_e\times(x-p_e)/2$, by direct differentiation of
the quadratic vector potential. Differentiating its three components
once more gives $\nabla\times U_S(p_e)=\omega_e$. Consequently
\begin{equation}\label{eq:band-fluid-sampling}
 \lambda^2|\nabla\times U_S(p_e)|^2=h_e\quad\hbox{for every original edge }e,
 \qquad P_g\Xi_{\lambda^2|\nabla\times U_S|^2,g}=w.
\end{equation}
The state notation on the right uses exactly these samples, not a
different rule for the face weights. This proves the full right-inverse
identity, including its nonlinear square root and added constant.
For $w=0$ the formula gives $S=0$, $M_S=0$ and $U_S=0$.

These profiles occur in explicit smooth forced Navier--Stokes flows.
Fix $s_0>0$ and choose a real smooth function $\theta$ with compact
support in $(0,\infty)$, equal to one near $s_0$. Set
\begin{align}
 u(s,x)&=\theta(s)U_S(x),\qquad p(s,x)=0,\nonumber\\
 f(s,x)&=\theta'(s)U_S(x)
    +\theta(s)^2(U_S\cdot\nabla)U_S(x)-\nu\theta(s)\Delta U_S(x).
 \label{eq:band-forced-fluid}
\end{align}
Direct substitution proves, with the original positive viscosity,
$\partial_su+(u\cdot\nabla)u=-\nabla p+\nu\Delta u+f$ and
$\nabla\cdot u=0$ on $\mathbb R^3$ for all $s\ge0$.
Both velocity and force are smooth, spatially compactly supported
and time-compact, with zero initial data for this chosen $\theta$;
in particular their energy is finite at every time. The pressure,
viscous and transport terms are all displayed. The fixed-colour
connection of the earlier bridge still has its full source
$j_i=(\lambda/g^2)(\Delta u_i-c^{-2}\partial_s^2u_i)T$;
the construction has not imposed a sourceless Yang--Mills equation.
At $s=s_0$, (223) gives the prescribed
actual quantum band state, with precisely its original norm and
energy. No state has been divided by its norm.

This proves surjectivity for the specified map from smooth forced
fluid profiles to the real first physical band. The displayed force
depends on $w,L,a,g$ through $S$ and its local vector potentials.
It is not the fixed force or singular profile stipulated in the
counterexample programme. To transfer that particular Navier--Stokes
breakdown one must calculate its own samples and retained limits
using (212); the right inverse proves that the
curvature representation itself does not exclude the band states
whose full interacting energies are now given in these coordinates.

\section{Exact non-Abelian map of the released viscous source}
\label{sec:nonabelian-fluid-profile}

This section composes the actual released positive-viscosity source
profile with a three-colour connection. It is a coordinate-level
calculation on the source field; it is not a replacement of that field by
an Euler solution, and it is not a claim that the resulting source-free
Yang--Mills theory has been constructed. The source files retained for
provenance are
\begin{sloppypar}
\texttt{released\_ns\allowbreak\_reader\allowbreak\_corrected\allowbreak\_20260909.tex},
\texttt{endpoint\_derivation.tex},
and
\texttt{axis\_continuation\_body.tex}.
The first has 854321 bytes, the second 25620 bytes, and the third 41080
bytes in this workspace. The latter two are exact source continuations;
their construction theorems are attributed to those sources. The
identities below are independently checked by
\texttt{check\_nonabelian\_flow.py}, whose receipt records 41 exact
symbolic tests.
\end{sloppypar}

\subsection{Spaces, units, and the exact map}

Let \(I\) be the source's preterminal time interval and let
\[
 u\in C^\infty(I\times\mathbb R^3;\mathbb R^3),\qquad
 \partial_i u_i=0.
\tag{225}
\]
be the actual localized field in the retained source. Pointwise formulas
need no support hypothesis; in each spatial integral below the source's
compact support is retained. Keep \(c>0\), \(g>0\), and
\(\lambda\in\mathbb R\setminus\{0\}\), with \(X^0=cs\), \(X^i=x^i\) and
\(\eta_{\mu\nu}=\operatorname{diag}(-1,1,1,1)\). Let
\(T_a=-i\sigma_a/2\), so
\[
 [T_a,T_b]=\varepsilon_{abd}T_d,\qquad
 -2\operatorname{tr}(T_aT_b)=\delta_{ab}.
\tag{226}
\]
Identify colour coefficient vectors with the oriented spatial basis, but
do not identify the two fields. Define
\[
 {\bf A}_0=0,\qquad {\bf A}_i=\lambda\,e_i\times u,\qquad
 A_\mu={\bf A}_\mu\cdot T.
\tag{227}
\]
The exact inverse on this displayed image is
\[
 u=-\frac1{2\lambda}\sum_{i=1}^3e_i\times{\bf A}_i.
\tag{228}
\]
Indeed \(\sum_i e_i\times(e_i\times u)=-2u\). The earlier Cartan map
\(C_i=\lambda u_iT_3\) and (227) are related by two explicit compositions:
\[
 C_i=\lambda u_iT_3\longmapsto {\bf A}_i=\lambda(e_i\times u),
 {\bf A}\longmapsto
 C_i=\lambda\left[-\frac1{2\lambda}\sum_j e_j\times{\bf A}_j\right]_iT_3.
\tag{229}
\]
Each composition is the identity on its stated image. This proves the
type morphism between the two encodings while making no assertion that
either representative map is injective on all gauge orbits.

\subsection{Curvature, its invariant norm, and the current}

For \(F_{\mu\nu}=\partial_\mu A_\nu-\partial_\nu A_\mu+[A_\mu,A_\nu]\),
write \({\bf E}_i={\bf F}_{0i}\) and
\({\bf B}_i=\frac12\varepsilon_{ijk}{\bf F}_{jk}\). Direct expansion of
(227), retaining the commutator, gives
\[
 {\bf E}_i=\frac{\lambda}{c}e_i\times u_s,\qquad
 B_i^a=\lambda(\partial_a u_i-\delta_{ia}\partial_k u_k)
               +\lambda^2u_i u_a.
\tag{230}
\]
The epsilon contractions used here are
\[
 \varepsilon_{ijk}\varepsilon_{akb}
 =\delta_{ib}\delta_{ja}-\delta_{ia}\delta_{jb},\qquad
 ((e_j\times u)\times(e_k\times u))_a
 =\varepsilon_{jkb}u_bu_a.
\tag{231}
\]
Thus incompressibility gives
\[
 B_i^a=\lambda\partial_a u_i+\lambda^2u_i u_a,\quad
 |{\bf B}|^2=\lambda^2|\nabla u|^2+
 2\lambda^3u_i u_a\partial_a u_i+\lambda^4|u|^4
 =\lambda^2|\nabla u|^2+\lambda^3\partial_a(|u|^2u_a)
 +\lambda^4|u|^4.
\tag{232}
\]
For compact support,
\[
 \int|{\bf B}|^2\,dx
 =\lambda^2\|\nabla u\|_2^2+\lambda^4\|u\|_4^4
 =\lambda^2\|\operatorname{curl}u\|_2^2+\lambda^4\|u\|_4^4,\qquad
 \int|{\bf E}|^2\,dx=\frac{2\lambda^2}{c^2}\|u_s\|_2^2.
\tag{233}
\]
The second equality uses integration by parts and the cubic divergence in
(232) has not been discarded pointwise. Moreover
\(\operatorname{tr}B=\lambda^2|u|^2\), hence
\[
 |{\bf B}|^2\ge\frac{\lambda^4}{3}|u|^4.
\tag{234}
\]
The trace is frame-dependent; the norm and the inequality are invariant.

Define the induced current by the full sourced equation
\[
 D^\mu F_{\mu\nu}=g^2j_\nu,\qquad D_\mu V=\partial_\mu V+[A_\mu,V].
\tag{235}
\]
The exact lower-index Gauss component and spatial component are
\[
 g^2{\bf j}_0=-\frac{\lambda}{c}\operatorname{curl}u_s
              -\frac{\lambda^2}{c}u\times u_s,
\tag{236}
\]
\[
 g^2{\bf j}_i=-\frac{\lambda}{c^2}e_i\times u_{ss}
 -\lambda\nabla\omega_i
 -\lambda^2\left(2\omega_i u+
 \varepsilon_{ijk}u_k\partial_j u+e_i\times((u\cdot\nabla)u)\right)
 -\lambda^3|u|^2(e_i\times u),\quad \omega=\operatorname{curl}u.
\tag{237}
\]
To prove (237), first use the unreduced exact form
\[
 g^2{\bf j}_i=-\frac{\lambda}{c^2}e_i\times u_{ss}
 -\varepsilon_{ijk}\{\partial_j{\bf B}_k+
 \lambda(e_j\times u)\times{\bf B}_k\},
\tag{238}
\]
then insert (230) and contract all three epsilon indices. The coefficient
\(2\omega_i u\) consists of two terms: differentiating the quadratic
part of (230), and commuting its linear part. Antisymmetry and Jacobi give
\[
 g^2D^\nu j_\nu
 =\tfrac12[D^\nu,D^\mu]F_{\mu\nu}
 =\tfrac12[F^{\nu\mu},F_{\mu\nu}]=0.
\tag{239}
\]

If the same actual field satisfies the retained forced viscous equation,
\[
 u_s+(u\cdot\nabla)u-\nu\Delta u+\nabla p=f,\qquad \nu>0,
\tag{240}
\]
set
\[
 N=\nu\Delta u-(u\cdot\nabla)u-\nabla p+f=u_s,\qquad
 u_{ss}=\nu\Delta N-(N\cdot\nabla)u-(u\cdot\nabla)N-\nabla p_s+f_s.
\tag{241}
\]
Equations (236)--(238) with (241) are the exact fluid-to-current
morphism. Viscosity, pressure, forcing and both differentiated transport
terms remain present.

\subsection{The released axis jet and an actual curvature consequence}

The released source continuation fixes \(0<h<1/100\), \(C>0\), and
\(\tau=1-s>0\), and proves, on a common preterminal axis neighborhood,
\[
 u_1=a(\sigma,z,s)x_1-b(\sigma,z,s)x_2,\quad
 u_2=a(\sigma,z,s)x_2+b(\sigma,z,s)x_1,\quad
 u_3=w(\sigma,z,s),\quad \sigma=(x_1^2+x_2^2)/2,
\tag{242}
\]
\[
 b(0,0,s)=C^{-1}\tau^{-1-h},\qquad
 \partial_3w(0,0,s)=-2a(0,0,s).
\tag{243}
\]
The first equality is the source's actual axis quotient; the second is
incompressibility, not an imposed symmetry reduction. At the axis the
full spatial gradient and velocity are therefore
\[
 \nabla u=\begin{pmatrix}a&-b&0\\ b&a&0\\0&0&-2a\end{pmatrix},
 \qquad u=(0,0,w).
\tag{244}
\]
Substitution into (230), with no coefficient suppressed, gives
\[
 |{\bf B}|^2_{\rm axis}
 =2\lambda^2(a^2+b^2)+(-2\lambda a+\lambda^2w^2)^2
 \ge 2\lambda^2C^{-2}\tau^{-2-2h}.
\tag{245}
\]
Thus the source's axis vorticity growth is invariant curvature growth for
this exact non-Abelian representative, including \(a\) and \(w\). Since
\(\partial_s b=(1+h)C^{-1}\tau^{-2-h}\),
\[
 \operatorname{curl}u_s(0,0,s)=
 2(1+h)C^{-1}\tau^{-2-h}e_3,\qquad
 u(0,0,s)\times u_s(0,0,s)=0,
\]
and hence
\[
 g^2{\bf j}_0(0,0,s)
 =-\frac{2\lambda(1+h)}{cC}\tau^{-2-h}e_3.
\tag{246}
\]
This is an exact nonzero sourced Yang--Mills component. It prevents the
axis calculation from being misreported as a sourceless solution.

The source's viscosity map is retained with no parameter replacement:
\[
 x=\sqrt{\nu}\,y,\quad
 u_\nu(x,s)=\sqrt{\nu}\,u_*(y,s),\quad
 p_\nu(x,s)=\nu p_*(y,s),\quad f_\nu(x,s)=\sqrt{\nu}f_*(y,s).
\tag{247}
\]
Its inverse divides by the displayed factors after \(y=x/\sqrt{\nu}\).
The chain rule gives
\(\mathcal R_\nu(u_\nu,p_\nu,f_\nu)=
\sqrt{\nu}\mathcal R_1(u_*,p_*,f_*)\) and preserves zero initial data,
support and the full forcing residual. This map is between the source's
viscosity-one presentation and every stated \(\nu>0\); it does not set
physical viscosity to one.

\subsection{Composition with the actual quantum state map}

At every regular source time, retain the positive density and the original
quarter-face weighted operator from Sections~16--18. The exact chain is
\[
 u\longmapsto h=|{\bf B}|^2
 \longmapsto D_{h,g}
 \longmapsto \Xi_{h,g}=(D_{h,g}-\langle\psi_g,D_{h,g}\psi_g\rangle)\psi_g
 \longmapsto e^{-tA_g/2}\Xi_{h,g}.
\tag{248}
\]
Here \(h\) is nonnegative, but its cubic and quartic terms in (232) remain
in its Fourier transform:
\[
 \widehat h(k)=\lambda^2\widehat{|\nabla u|^2}(k)
 +i\lambda^3 k\cdot\widehat{|u|^2u}(k)
 +\lambda^4\widehat{|u|^4}(k).
\tag{249}
\]
The nine ordered products in \(|\widehat h|^2\) remain in the full
two-creation spectral measure. The exact Section16 comparison map
\(J_t h\) consequently has
\[
 \|J_t h\|^2=\frac1{320\pi^5t}\int_{\mathbb R^3}|k|^4e^{-t|k|}
 |\widehat h(k)|^2\,dk,\qquad t>0.
\tag{250}
\]
If \(H=\int h\ne0\), its point-concentration comparison has raw norm
\(9H^2/(\pi^4t^8)\) and raw excitation \(72H^2/(\pi^4t^9)\), with quotient
\(8/t\). These are unchanged raw states; no state is divided by its norm.
Equations (248)--(250) are a free comparison representation with the
original interacting finite-box state map preceding it. They do not prove
that the interacting continuum Hamiltonian has this spectrum.

\subsection{Scope of the propagated result}

The exact map (227)--(249) proves a non-Abelian, gauge-invariant curvature
and current consequence for the released positive-viscosity profile,
including all original source terms. It does not turn a sourced classical
connection into a source-free quantum Yang--Mills solution, and it does
not by itself establish a zero mass gap. The remaining interacting
question is whether the actual finite-box states in (248), with their
full \(D_{h,g}\), Gram corrections, and regulator-dependent interaction
matrix, admit a jointly controlled continuum limit. A negative or
divergent classical term is not substituted for that spectral proof.

\section{Released positive-viscosity profile: gauge-invariant curvature measures and the quantum-state map}
\label{sec:released-profile-measures}

The preceding non-Abelian embedding gives pointwise axis information.  The released positive-viscosity source supplies the stronger spatial estimates needed to pass from an axis value to an actual gauge-invariant spatial integral.  We use its original coordinates, viscosity and terminal variable
\[
  \tau=1-s,\qquad \nu>0,\qquad c>0,\qquad g>0,
\]
and retain the fixed nonzero bridge coefficient $\lambda$.  The source reader and its endpoint continuation prove that, for positive constants $C_\Omega$ and $C_{\dot u}$ fixed by the source profile,
\[
 \|\operatorname{curl}u_\nu(1-\tau)\|_2^2
 \ge \nu^{3/2}C_\Omega\tau^{-1/2-3h},\qquad
 \|\partial_su_\nu(1-\tau)\|_2^2
 \ge \nu^{5/2}C_{\dot u}\tau^{-3/2-3h},
 \tag{250}
\]
for the source range $0<h<1/100$ and sufficiently small $\tau>0$.  These are lower bounds on the full spatial integrals, not evaluations on the measure-zero axis.

Use the exact reducible connection in the original Cartesian chart,
\[
 A_0=0,\qquad A_i=\lambda u_{\nu,i}T,
 \qquad T=-i\sigma_3/2,
 \qquad -2\operatorname{tr}(T^2)=1.
 \tag{251}
\]
Since every component has the same Lie generator, $[A_\mu,A_\nu]=0$ and
\[
 F_{0i}=\lambda c^{-1}\partial_su_{\nu,i}T,
 \qquad
 F_{ij}=\lambda(\partial_i u_{\nu,j}-\partial_j u_{\nu,i})T.
 \tag{252}
\]
The two positive component norms are therefore exactly
\[
\begin{aligned}
 \mathcal K_B(s)&:=-2\int\sum_{i<j}\operatorname{tr}(F_{ij}^2)\,dx
      =\lambda^2\|\operatorname{curl}u_\nu(s)\|_2^2,\\
 \mathcal K_E(s)&:=-2\int\sum_i\operatorname{tr}(F_{0i}^2)\,dx
      =\lambda^2c^{-2}\|\partial_su_\nu(s)\|_2^2.
\end{aligned}
\tag{253}
\]
Combining (250) and (253) gives the exact terminal lower bounds
\[
 \mathcal K_B(1-\tau)\ge\lambda^2\nu^{3/2}C_\Omega\tau^{-1/2-3h},\qquad
 \mathcal K_E(1-\tau)\ge\lambda^2c^{-2}\nu^{5/2}C_{\dot u}\tau^{-3/2-3h}.
 \tag{254}
\]
Hence both component norms diverge as $\tau\downarrow0$, while the source energy identity gives
\[
 \int_0^1\mathcal K_B(s)\,ds<\infty,
 \qquad
 \int_0^1\mathcal K_E(s)\,ds=\infty.
 \tag{255}
\]
The first statement uses the original positive-viscosity dissipation bound; the second follows by integrating the nonintegrable exponent in (254).  These are Euclidean component norms.  They do not assert a sign for the Lorentzian action, whose metric has signature $(-,+,+,+)$.

The full sourced Yang--Mills equation remains coordinate-level precise.  With $D^\mu F_{\mu\nu}=g^2j_\nu$,
\[
 j_0=0,
 \qquad
 j_i=\frac{\lambda}{g^2}
 \left(\Delta u_{\nu,i}-c^{-2}\partial_s^2u_{\nu,i}\right)T,
 \tag{256}
\]
and substitution of the complete positive-viscosity equation gives
\[
 j_i=\frac{\lambda}{g^2}\left[\Delta u_{\nu,i}
 -c^{-2}\left(\nu\Delta(\partial_su_{\nu,i})
 -(\partial_su_\nu\!\cdot\!\nabla)u_{\nu,i}
 -(u_\nu\!\cdot\!\nabla)\partial_su_{\nu,i}
 -\partial_i\partial_sp_\nu+\partial_sf_{\nu,i}\right)\right]T.
 \tag{257}
\]
Every viscosity, pressure, force and transport term is displayed.  Taking the spatial divergence of (256) gives $D^\nu j_\nu=0$ from incompressibility, so the sourced map is internally consistent but is not a source-free Yang--Mills solution.

At each $s<1$, the source profile is smooth and compactly supported.  The input to the spatial observable map is the pointwise gauge-invariant density, not the already integrated scalar.  Define
\[
\begin{aligned}
 h_s^{B}(x)&:=-2\sum_{i<j}\operatorname{tr}(F_{ij}(s,x)^2)\\
            &=\lambda^2|\operatorname{curl}u_\nu(s,x)|^2,
\end{aligned}
\]
\[
\begin{aligned}
 h_s^{E}(x)&:=-2\sum_i\operatorname{tr}(F_{0i}(s,x)^2)\\
            &=\lambda^2c^{-2}|\partial_su_\nu(s,x)|^2.
\end{aligned}
\]
\[
 h_s:=h_s^{B}\quad\text{or}\quad h_s:=h_s^{B}+h_s^{E},\qquad
 H_{h_s}:=\int_{\mathbb R^3}h_s(x)\,dx.
\]
Thus $H_{h_s}=\mathcal K_B(s)$ for the magnetic choice and $H_{h_s}=\mathcal K_B(s)+\mathcal K_E(s)$ for the total choice.  For each lattice spacing $a>0$ and box size $L$, write an oriented edge as $e=(n,i)$ and define its physical midpoint $x_e=a(n+\tfrac12e_i)$.  The exact sampling morphism and face extension are
\[
 I_{L,a}:C_c^\infty(\mathbb R^3)\longrightarrow\mathbb R^{\mathsf E_L},\qquad (I_{L,a}h)_e=h(x_e),\qquad
 (I^P_{L,a}h)_p=\frac14\sum_{e\in\partial p}(I_{L,a}h)_e.
\]
It sends the pointwise continuum density to the original edge and face coefficients, with no averaging or normalization of $h$.  The same map has an explicit physical-volume comparison.  If $h\in C_c^1(\mathbb R^3)$ is supported in a compact set $K$ and the box contains all edge midpoints meeting $K_{a\sqrt3/2}$, then for each direction $i$ the midpoint cubes form a partition and
\[
 \left|a^3\sum_{e=(n,i)}h(x_e)-\int_{\mathbb R^3}h(x)\,dx\right|
 \le \frac{\sqrt3}{2}a\,\operatorname{vol}(K_{a\sqrt3/2})\,\|\nabla h\|_\infty.
\]
Indeed, integrate $h(x_e)-h(x)$ over the cube centered at $x_e$ and use the mean-value bound $|h(x_e)-h(x)|\le(\sqrt3a/2)\|\nabla h\|_\infty$; cubes outside $K_{a\sqrt3/2}$ contribute zero.  Thus the exact midpoint morphism preserves the unnormalized curvature mass in the joint lattice limit, with a displayed $O(a)$ error for each fixed profile.  In the following display, $D_{h,g}$ abbreviates this $D_{L,a,g}[h]$ at the displayed fixed regulator.  The resulting operator is exactly
\[
 D_{L,a,g}[h]=\kappa\sum_{e\in\mathsf E_L}(I_{L,a}h)_eE_e
 +b\sum_{p\in\mathsf P_L}(I^P_{L,a}h)_p(2-W_p),
\]
so the full edge/face sets and the physical scales remain present.  This is the domain-correct composition $h\mapsto I_{L,a}h\mapsto D_{L,a,g}[h]\mapsto\Xi_{L,a,g}[h]$ used below.  Apply the already proved exact observable map
\[
 h_s\longmapsto D_{h_s,g}
 \longmapsto \Xi_{h_s,g}=
 (D_{h_s,g}-\langle\psi_g,D_{h_s,g}\psi_g\rangle)\psi_g
 \longmapsto e^{-tA_g/2}\Xi_{h_s,g}.
 \tag{258}
\]
The centering is exact and no division by the raw norm occurs.  Gauge invariance is a projection, not a normalization: if $\mathcal G_L$ is the original vertex gauge group with product Haar probability, then
\[
 (\Pi_G f)(U)=\int_{\mathcal G_L}f(g\mathbin{\cdot}U)\,dg,
 \qquad \Pi_G^*=\Pi_G,\quad \Pi_G^2=\Pi_G,\quad \|\Pi_Gf\|\le\|f\|.
\]
Haar invariance proves the two identities and the contraction inequality.  The operators $D_{L,a,g}[h]$ and the vacuum commute with this action, so $\Pi_G\Xi_{L,a,g}[h]=\Xi_{L,a,g}[h]$ exactly; no scalar rescaling is introduced.  Vacuum centering removes only the component along $\psi_g$, whereas division by a state norm would define a separate normalized spectral probability and is not performed here.  The common free spectral map gives, for each fixed $s<1$ and $t>0$,
\[
 \|J_th_s\|^2\le \frac{9}{\pi^4t^8}\,\|h_s\|_1^2,
 \qquad
 H_{h_s}=\int h_s\,dx=\lambda^2\|\operatorname{curl}u_\nu(s)\|_2^2>0
\tag{259}
\]
whenever the profile is nonzero.  Thus (254) proves divergence of the input mass $H_{h_s}$ at the terminal endpoint, but (259) also shows why that divergence alone does not prove a zero interacting mass gap: the centered quantum vector, its actual interacting norm, and the joint regulator/volume limit still require estimates not supplied by the classical source.  The exact propagation established here is the map (251)--(258), with all source terms and units retained.

The inspected source provenance is
\texttt{ns\_inner\_profile\_curvature\_current\_bridge.md}.
The released reader and endpoint/axis continuation are recorded in the lane
provenance files.  No theorem about a source-free four-dimensional Yang--Mills
mass gap is inferred from the sourced profile.



\section{A support-uniform lower bound for the free curvature state}
\label{sec:support-uniform-free-state}

The compact support statement of the released source yields a quantitative consequence that uses the unnormalized state directly.  Let $K\subset\mathbb R^3$ be a fixed compact support for $h_s\ge0$, and choose $R>0$ with $K\subset\{x:|x|\le R\}$.  Write $H_s=\int h_s(x)\,dx$.  For every $k$ with $|k|\le(2R)^{-1}$,
\[
 |\widehat h_s(k)-H_s|
 =\left|\int_K(e^{-ik\cdot x}-1)h_s(x)\,dx\right|
 \le |k|R H_s\le \tfrac12H_s,
\]
so $|\widehat h_s(k)|\ge H_s/2$.  Inserting this exact bound into the already proved free norm formula gives
\[
 \|J_t h_s\|^2
 \ge \frac{H_s^2}{1280\pi^5t}
       4\pi\int_0^{(2R)^{-1}}r^6e^{-tr}\,dr
 =:C_{R,t}H_s^2>0.\tag{260}
\]
The corresponding free energy form obeys the previously established total-variation estimate
\[
 \mathcal E_t(J_t h_s)\le \frac{72}{\pi^4t^9}H_s^2.\tag{261}
\]
Consequently
\[
 \frac{\mathcal E_t(J_t h_s)}{\|J_t h_s\|^2}
 \le \frac{72}{\pi^4t^9C_{R,t}}.
 \tag{262}
\]
For fixed $R$, the exact integral in (260) satisfies $C_{R,t}\sim9/(4\pi^4t^8)$ as $t\to\infty$, hence the right side of (262) is asymptotic to $32/t$.  This proves a zero-quotient sequence for the free common-state representation for every nonzero nonnegative compactly supported curvature weight, including the released profile at each preterminal time.  Along the terminal sequence, (254) makes $H_s$ diverge while the same quotient bound remains independent of $H_s$; no norm division was used.

The conclusion still concerns the explicitly identified free comparison Hamiltonian.  Transferring (260)--(262) to the interacting $D_{h_s,g}$ requires the existing positive-coupling norm/energy convergence at each fixed $s,t$ and a jointly chosen $(s,t,g,L)$ trajectory.  The classical source proves neither that joint limit nor a source-free Yang--Mills solution, so this exact free-state zero quotient is a necessary bridge calculation rather than the final mass-gap contradiction.

\section{Diagonal transfer along the proved weak-coupling path}
\label{sec:diagonal-profile-transfer}

The fixed-profile convergence theorem can be combined with the released terminal sequence without normalizing the curvature.  Choose $s_n=1-1/n$ (discarding finitely many indices so that $s_n$ lies in the smooth source interval), set $h_n=\lambda^2|\operatorname{curl}u_\nu(s_n)|^2$, and let $t_n=n$.  Every $h_n$ is nonnegative, smooth and supported in the same compact set.  The support-uniform estimate (260)--(262) gives a positive free norm and a free quotient bounded by $32/n+o(n^{-1})$.

For each $n$, the established positive-coupling theorem supplies a tail of the original sequence
\[
 L_j=j^2,\qquad a_j=(100j)^{-1},\qquad g_j=2^{-m_j}>0,
 \tag{263}
\]
for which the actual centered vectors $y_{h_n,g_j}(t_n)$ have both their raw norm and their energy within any prescribed relative error of the corresponding free quantities.  The profile-dependent error is finite because $h_n$ is smooth and compactly supported at each fixed $n$.  Recursively choose $j_n>j_{n-1}$ in that tail with both relative errors below $1/n$; this is a genuine diagonal subsequence and does not alter $h_n$, $t_n$, $a_j$, or the source units.  Then
\[
 \|y_{h_n,g_{j_n}}(t_n)\|^2
   =\|J_{t_n}h_n\|^2(1+o(1)),\qquad
 \langle y_{h_n,g_{j_n}}(t_n),A_{g_{j_n}}y_{h_n,g_{j_n}}(t_n)\rangle
   =\mathcal E_{t_n}(J_{t_n}h_n)(1+o(1)),
 \tag{264}
\]
and therefore
\[
 \frac{\langle y_{h_n,g_{j_n}}(t_n),A_{g_{j_n}}y_{h_n,g_{j_n}}(t_n)\rangle}
  {\|y_{h_n,g_{j_n}}(t_n)\|^2}
 \le \frac{32}{n}+o(n^{-1})\longrightarrow0.
 \tag{265}
\]
The vectors are vacuum-orthogonal by the exact centering in (258), and no quotient or state norm was used in constructing them.  Equations (263)--(265) are consequently an actual zero-quotient sequence for the changing finite-regulator Hamiltonians on the proved weak-coupling continuum path, with the Navier--Stokes curvature mass retained and diverging along $s_n\uparrow1$.

This does not yet prove failure of the interacting four-dimensional Yang--Mills mass-gap assertion.  The selected couplings satisfy $g_{j_n}\to0$, and the existing convergence theorem controls the free comparison path; it does not prove survival of a fixed renormalized non-Abelian interaction or identify the resulting limit with the target interacting theory.  The diagonal construction isolates the exact remaining requirement: replace (263) by a physical-scale trajectory on which the full interacting Wilson/Haar terms remain present, then repeat the relative-error argument with those terms.  Until that trajectory is proved, (265) is a rigorous weak-coupling counterexample candidate, not a completed Millennium disproof.

\section{Fixed-coupling diagnostic for the competing magnetic channel}
\label{sec:fixed-coupling-diagnostic}

The root Yang--Mills lane supplies an exact fixed-coupling calculation that fixes the next obstruction.  On the unchanged geometric path $L_j=j^2$, $a_j=(100j)^{-1}$, retain $g_j=g_0>0$, so $\xi_0=(4g_0^4)^{-1}$ and $M_{L_j}=12j^4(2j^2+1)$.  For the explicitly defined cusp magnetic translation with $\theta_j=2\pi a_j^2/D_j$, the full interacting vacuum estimate is
\[
 \|\chi_j\|\le \frac{6144\pi^2\xi_0}{10^8}\,j^{-2},
 \qquad
 d_j=\|\chi_j\|^2\le
 \left(\frac{6144\pi^2\xi_0}{10^8}\right)^2j^{-4}.
 \tag{266}
\]
This estimate is obtained from the exact all-angle operator inequality and the full positive vacuum; no constant-vacuum replacement or face deletion is made.  It says that the unnormalized centered measures of this fixed-coupling channel converge to zero in total mass.  Division by $d_j$ leaves a normalized spectral probability, but (266) supplies no low-energy estimate for that probability.

The same dictionary gives $\xi_0M_{L_j}\to\infty$ whenever $g_0$ is fixed and $a_j\downarrow0$.  Therefore the small-interaction regime used by the weak-coupling diagonal is unavailable on this fixed-coupling path.  The exact fixed-coupling result neither proves nor rules out a different non-Abelian state with a vanishing physical quotient; it records why the present NS curvature-state sequence cannot simply be relabelled as such a state.  A fixed-renormalized-coupling counterexample still requires a new state construction whose norm and energy remain controlled when $\xi M_L$ diverges.





\clearpage
\section*{References}
\paragraph{[1]}
G.~Dusson, I.~M.~Sigal, and B.~Stamm,
\emph{The Feshbach--Schur map and perturbation theory},
in \emph{Partial Differential Equations, Spectral Theory, and Mathematical
Physics}, EMS Series of Congress Reports (2021), pp.~65--88.
\href{https://doi.org/10.4171/ECR/18-1/5}{doi:10.4171/ECR/18-1/5};
\href{https://arxiv.org/abs/2105.02058v1}{arXiv:2105.02058v1},
Theorem~1.2 and Eqs.~(1.10)--(1.16).
\paragraph{[2]}
D.~Bakry and M.~Ledoux,
\emph{A logarithmic Sobolev form of the Li--Yau parabolic inequality},
Revista Matem\'atica Iberoamericana \textbf{22} (2006), no.~2, 683--702.
\href{https://doi.org/10.4171/RMI/470}{doi:10.4171/RMI/470},
Eq.~(1.7) and its proof, pp.~686--687.
\paragraph{[3]}
B.~Simon, \emph{A Feynman--Kac Formula for Unbounded Semigroups},
\href{https://arxiv.org/abs/math-ph/9907022v1}{arXiv:math-ph/9907022v1}
(1999), Theorem~1.1 and Eqs.~(1.1)--(1.3).
\paragraph{[4]}
H.~Mori, \emph{Transport, Collective Motion, and Brownian Motion},
Progress of Theoretical Physics \textbf{33} (1965), 423--455.
\href{https://doi.org/10.1143/PTP.33.423}{doi:10.1143/PTP.33.423}.
\paragraph{[5]}
J.~Kogut and L.~Susskind,
\emph{Hamiltonian formulation of Wilson's lattice gauge theories},
Physical Review D \textbf{11} (1975), 395--408.
\href{https://doi.org/10.1103/PhysRevD.11.395}{doi:10.1103/PhysRevD.11.395}.
\paragraph{[6]}
F.~Peter and H.~Weyl,
\emph{Die Vollst\"andigkeit der primitiven Darstellungen einer
geschlossenen kontinuierlichen Gruppe},
Mathematische Annalen \textbf{97} (1927), 737--755.
\href{https://doi.org/10.1007/BF01447892}{doi:10.1007/BF01447892}.
\paragraph{[7]}
D.~A.~Yarotsky, \emph{Quasi-particles in weak perturbations of
non-interacting quantum lattice systems},
\href{https://arxiv.org/abs/math-ph/0411042v1}{arXiv:math-ph/0411042v1}
(2004), Section 2.
\paragraph{[8]}
J.~C.~Baez, \emph{Spin Network States in Gauge Theory},
Advances in Mathematics \textbf{117} (1996), 253--272.
\href{https://doi.org/10.1006/aima.1996.0012}{doi:10.1006/aima.1996.0012};
\href{https://arxiv.org/abs/gr-qc/9411007v1}{arXiv:gr-qc/9411007v1}.

\end{document}


